% =====================================================================
% hetlearn technical appendix
% Standalone article. Compile with pdflatex + bibtex (see README.txt).
% =====================================================================
\documentclass[11pt,a4paper]{article}
\PassOptionsToPackage{obeyspaces,spaces}{url}

\usepackage[T1]{fontenc}
\usepackage[utf8]{inputenc}
\usepackage{lmodern}
\usepackage{textcomp}
\usepackage[margin=2.4cm]{geometry}
\usepackage{amsmath}
\usepackage{amssymb}
\usepackage{booktabs}
\usepackage{array}
\usepackage{longtable}
\usepackage{enumitem}
\usepackage{xcolor}
\usepackage{listings}
\usepackage[authoryear,round]{natbib}
\usepackage{url}
\usepackage[colorlinks=true,linkcolor=blue!50!black,citecolor=blue!50!black,urlcolor=blue!50!black]{hyperref}

% ---- inline code: option names, file names, MATLAB identifiers --------------------------
% \code{...} prints its argument literally and can break at punctuation and spaces.
% \codem{...} prints the same but cannot break: use it in headings, inside \texorpdfstring
% and inside math. Do not put %, # or a backslash inside either command.
\DeclareUrlCommand{\code}{\urlstyle{tt}}
\makeatletter\g@addto@macro{\UrlBreaks}{\do\:}\makeatother   % break error ids at colons
\DeclareRobustCommand{\codem}[1]{\texttt{\detokenize{#1}}}
\setlength{\emergencystretch}{2em}

% ---- cfg examples ---------------------------------------------------------------------
\lstdefinestyle{cfg}{language=Matlab,
  basicstyle=\ttfamily\small,
  columns=fullflexible,
  keepspaces=true,
  upquote=true,
  showstringspaces=false,
  breaklines=true,
  frame=single,
  framerule=0.2pt,
  rulecolor=\color{black!35},
  backgroundcolor=\color{black!3},
  commentstyle=\color{black!60}\itshape,
  keywordstyle=\bfseries,
  aboveskip=4pt,
  belowskip=4pt}
\lstset{style=cfg}

% ---- option entries --------------------------------------------------------------------
\newenvironment{optspec}%
  {\begin{description}[style=sameline,leftmargin=2.7cm,labelwidth=2.5cm,font=\normalfont\itshape]}%
  {\end{description}}

% ---- notation ----------------------------------------------------------------------------
\newcommand{\RE}{\mathrm{RE}}
\newcommand{\BR}{\mathrm{BR}}
\newcommand{\fRE}{f^{\RE}}
\newcommand{\fBR}{f^{\BR}}
\newcommand{\Eagg}{\tilde{E}}
\newcommand{\T}{\mathbf{T}}
\newcommand{\K}{\mathbf{K}}
\newcommand{\Rs}{\mathbf{R}}
\newcommand{\Z}{\mathbf{Z}}
\newcommand{\Hs}{\mathbf{H}}
\newcommand{\D}{\mathbf{D}}
\newcommand{\FE}{\mathit{FE}}
\newcommand{\transp}{^{\prime}}

\setcounter{tocdepth}{2}
\setlength{\parskip}{0.35\baselineskip}
\setlength{\parindent}{0pt}

\title{\textbf{hetlearn: technical appendix}\\[4pt]
\large Heterogeneous expectations with adaptive learning and switching\\ in IRIS models}
\author{Tolga \"Ozden\thanks{Code and documentation:
  \url{https://github.com/tolgaozden/iris_toolbox_learning}, released under the MIT License.}}
\date{Version of 23 October 2026. Tested with IRIS 20230622 and MATLAB R2025a.}

\begin{document}
\maketitle

\begin{abstract}
\noindent
This appendix documents the mathematics, the assumptions and every user-facing option of the
\code{hetlearn} toolbox. The toolbox takes a linear IRIS model (nonlinear models have a code path but
no test) with forward-looking terms and replaces each one-step expectation by a population mix of rational
(RE) agents and adaptive learners whose shares switch on past forecast performance, as in
\citet{ozden2025}. It writes the extended model file, runs a history stage (learning on the RE
smoother, the filtered RE states, an inversion of the learning model, or a time-varying Kalman
filter), and then forecasts quarter by quarter with the beliefs and shares in force, under a
projection facility. On top of the forecast stage sit conditional forecasts, Ramsey policy,
density forecasts, shock decompositions and a recursive real-time forecast evaluation. Each
option is given with its equations, the assumptions behind it, its default, its allowed values,
its interactions with other options and a short configuration example. Each equation names the
file that implements it.
\end{abstract}

\tableofcontents
\clearpage

% =====================================================================
\section{Introduction}\label{sec:intro}
% =====================================================================

\subsection{What the toolbox does}

\code{hetlearn} is a MATLAB package (\code{toolbox/+hetlearn/}) built on the IRIS toolbox
\citep{iris}. A user supplies a model package: an IRIS model file, a parameter function, an
IRIS-readable database, the observables and a filter range. The toolbox

\begin{enumerate}[nosep]
  \item rewrites the model so that every lead $x_{t+k}$ of a transition variable becomes an
        explicit expectation variable, and gives every forward-looking variable an expectation
        block that mixes an RE forecast and a learners' forecast (Section~\ref{sec:generator});
  \item runs an RE Kalman filter and smoother over the data (Section~\ref{sec:history});
  \item lets the learners form beliefs over the history, by recursive least squares (RLS) on
        a perceived law of motion (PLM), the reduced-form forecasting rule they estimate, while the population shares follow a discrete-choice
        rule on past squared forecast errors (Sections~\ref{sec:plm}--\ref{sec:switching});
  \item forecasts one quarter at a time from a forecast origin: in each quarter the model is
        solved with the beliefs and shares in force and simulated for that quarter only, and the
        beliefs and shares are updated with the new data point (Section~\ref{sec:forecast}).
\end{enumerate}

The learning and switching block follows \citet{ozden2025}, equations (12)--(16) and
Appendix~D of that paper. The endogenous gain follows \citet*[eq.~9]{koz2024}, hereafter KOZ, who apply the
adaptive step-size algorithm of \citet{kostyshyna2012}; the gain rule \code{'koz'} is named after
them.

\subsection{Public functions}\label{sec:functions}

The toolbox is used through these functions; \code{help hetlearn.<name>} gives the full text of each.

\begin{longtable}{>{\raggedright\arraybackslash}p{8.4cm}>{\raggedright\arraybackslash}p{6.6cm}}
\toprule
Call & Purpose \\
\midrule
\endhead
\code{res = hetlearn.run(cfg, ...)} & one scenario; options \code{'save'}, \code{'results_dir'}, \code{'quiet'}, \code{'model'} (preloaded), \code{'filter'} (an RE filter output) \\
\code{T = hetlearn.run_files(names, ...)} & runs scenario functions in order; a failure is reported and the next one runs \\
\code{T = hetlearn.compare(ids, vars, range)} & levels \code{<var>__<id>} and differences \code{<var>__<id>_minus_<id1>} of saved runs \\
\code{D = hetlearn.decompose(res, vars, range, opts)} & Section~\ref{sec:decompose} \\
\code{E = hetlearn.evaluate(name, variants, origins, opts)} & Section~\ref{sec:evaluate} \\
\code{T = hetlearn.evaluate_tables(errors, opts)} & tables from saved errors \\
\code{o = hetlearn.defaults()} & the option tree \\
\code{opts = hetlearn.load_options(cfg, master)} & merge and validate \\
\code{spec = hetlearn.load_spec(name)} & package contract \\
\code{[d, file] = hetlearn.load_database(spec)} & the package database \\
\code{[m, info] = hetlearn.load_model(spec, opts)} & generate, load, assign, solve the extended model \\
\code{H = hetlearn.re_filter(spec, opts, ...)} & RE filter and smoother; \code{'force'}, \code{'end'}, \code{'model'} \\
\code{meta = hetlearn.extend_model(spec, plm, ...)} & the generator \\
\code{[A, info] = hetlearn.re_anchor_coefficients(spec, P, ...)} & anchored coefficients \\
\code{t = hetlearn.find_peak(db, name, from, to)} & date of a series' maximum in a window \\
\bottomrule
\end{longtable}

\subsection{Scope of this appendix}

This appendix covers the core toolbox: two agent types (RE agents and learners), adaptive
learning, switching and every stage built on them. It leaves out two parts of the repository:

\begin{itemize}[nosep]
  \item the bounded-rationality extensions of \code{docs/DESIGN_br_extensions.md}: the options
        \code{expectations.types}, \code{expectations.transforms}, \code{shares.init},
        \code{shares.fixed}, the type kinds \code{'rule'} and \code{'re_mean'}, and the N-type
        paths of the engine (\code{+engine/ntypes.m}). A later chapter will document the
        further expectation rules. Every one of these options is off by default, and a run that sets none of them
        takes the two-type code described here;
  \item the Bayesian estimation module (\code{estimation/}, package \code{+hlest}), which has
        its own copy of the toolbox and no scenario uses it.
\end{itemize}

The time-varying Kalman history (\code{history.mode = 'kalman'}) belongs to the
extensions but works with two types, so Section~\ref{sec:kalman} documents its two-type form.

\subsection{How an option is documented}

Each option has its own entry with the same fields:

\begin{optspec}
  \item[Default] the value in \code{hetlearn.defaults()} (\code{defaults.m}). The project
        baseline in \code{config/master_options.m} overrides some defaults for the SW07
        scenarios; those values are noted where they differ, but the entries document the
        toolbox defaults.
  \item[Values] the values \code{hetlearn.load_options} accepts (\code{load_options.m}).
  \item[Math] the equations the option selects or parameterizes.
  \item[Assumptions] what the option takes as given about agents, information or the model.
  \item[Interactions] other options that change its effect, combinations that are errors, and
        options it makes irrelevant.
  \item[Code] the files that implement it. Paths are relative to \code{toolbox/+hetlearn/};
        \code{+engine/forecast.m} is the function \code{hetlearn.engine.forecast}.
  \item[Example] one or two lines of a scenario's \code{cfg}.
\end{optspec}

Appendix~\ref{app:index} lists every option with its default and the section that documents it.

\subsection{Option precedence and merging}\label{sec:precedence}

\code{opts = hetlearn.load_options(cfg, master)} (\code{load_options.m}) builds the option tree in
three layers: the defaults of \code{hetlearn.defaults()}, then \code{config/master_options.m}, then
the scenario's \code{cfg}, each layer overriding the previous one field by field.
\code{master_options()} is found on the path or in \code{<root>/config/}; the optional second argument replaces it (the tests use their own
\code{tests/scenarios/master_options.m}). Rules:
\begin{itemize}[nosep]
  \item the merge is recursive; a field that does not exist in the defaults is
        \code{hetlearn:options:unknown}, naming the full path (typos are caught);
  \item \emph{replaced as a whole}: \code{conditions} and \code{policy.loss} (a loss struct
        without \code{dr} gives \code{dr} weight 0);
  \item \emph{free-form}, merged without name checks: \code{beliefs.init.values},
        \code{beliefs.plm.by_var}, \code{beliefs.plm.derived} (must hold strings),
        \code{gain.by_var}, \code{shares.fixed_n2} (as a struct), \code{shocks.add},
        \code{report};
  \item strings become char; a char where a cellstr is expected becomes a one-element cellstr
        (\code{expectations.variables} other than \code{'all'}, \code{shares.master},
        \code{compare_with}, \code{shocks.exclude}, \code{density.shocks},
        \code{history.inversion.shocks}, \code{gain.koz.master}, \code{report.variables});
  \item the removed options \code{projection.max_root}, \code{.method}, \code{.scope},
        \code{.on_failure}, \code{switching.error_scale}, \code{.error_scale_ref} and the removed
        bounded-rationality options raise \code{hetlearn:options:removed}, which names the
        replacement. A saved \code{res.cfg} that still holds a removed field at
        its old default value loads without error.
\end{itemize}

Validation errors: \code{hetlearn:options:type} (wrong type), \code{:value} (outside the
enumeration or range, or a forbidden combination), \code{:missing} (a required companion option),
\code{:unsupported} (combinations not derived or built), \code{:notImplemented}
(\code{history.mode = 'hetero_filter'}). Checks that need the model (names of variables, shocks and
targets) are made later by the generator, the PLM builder and the engine.

\subsection{Conventions}

\begin{itemize}[nosep]
  \item Quarters are IRIS dates. The forecast origin $T_0$ is the last quarter of data used;
        the horizon is $T_0+1,\dots,T_0+H$.
  \item Model units are those of the package. In the \citet{smets2007} model (SW07), the main
        package of the repository, and in the toy model of the tests (\code{tests/toy_nk/}),
        inflation and the policy rate are in percent per quarter; reports multiply by
        \code{spec.report.scale} (4 annualizes). Learning series are always in model units.
  \item The SW07 data, seven US observables over the filter range 1965Q1--2026Q2 with every
        observable missing in 2020Q1--2020Q3, are described in \code{models/sw07/README.md}.
  \item This appendix calls the adaptive agents learners; the generated names use \code{BR}
        (bounded rationality) for them. Share names: $n_1$ is the RE share, $n_2 = 1-n_1$ the
        learners' share. Options give the learners' share (\code{shares.init_n2},
        \code{shares.fixed_n2}); Section~\ref{sec:sharenames} lists every share name.
  \item Tolerances quoted in this appendix are the ones the code checks.
\end{itemize}

% =====================================================================
\section{Notation and timing}\label{sec:notation}
% =====================================================================

\subsection{Symbols}

\begin{longtable}{>{\raggedright\arraybackslash}p{3.2cm}>{\raggedright\arraybackslash}p{7.2cm}>{\raggedright\arraybackslash}p{4.4cm}}
\toprule
Symbol & Meaning & Name in the code \\
\midrule
\endhead
$t$, $T_0$, $H$ & quarter, forecast origin, horizon & \code{forecast.origin}, \code{forecast.horizon} \\
$\mathcal{F}$ & forward-looking transition variables of the model & \code{meta.forward_vars} \\
$\mathcal{L}\subseteq\mathcal{F}$ & learning set: variables whose expectation is the RE/learner mix & \code{expectations.variables}, \code{info.learnVars} \\
$x\in\mathcal{L}$ & a learning variable & \code{x} \\
$\Eagg_t x_{t+1}$ & aggregate one-step expectation of $x$ in the model & \code{x_exp1} \\
$\fRE_t(x)$ & RE agents' forecast of $x_{t+1}$ made in $t$ & \code{x_exp1_RE} \\
$\fBR_t(x)$ & learners' forecast of $x_{t+1}$ made in $t$ & \code{x_exp1_BR} \\
$n_{1,t}$, $n_{2,t}$ & discrete-choice RE and learner shares computed at the end of $t$ & \code{res.learning.<x>.n1}, \code{n2} \\
$n^{m}_{1,t}$ & RE share in the model in quarter $t$ & parameter \code{n1_x}; \code{n1_model} \\
$\phi_t$ & PLM coefficients at the end of $t$ (in force in $t+1$) & \code{S.phi.<target>} \\
$R_t$ & RLS second-moment matrix & \code{S.R.<target>} \\
$z_{t-1}$ & regressor vector of a PLM equation, entries dated $t-1$ or earlier & \code{Z} in \code{update_learning.m} \\
$g_t$ & gain of the RLS update at the end of $t$ & \code{res.learning.<target>.gain} \\
$\gamma_0$ & constant gain & \code{gain.value} \\
$\zeta^{\RE}_t$, $\zeta^{\BR}_t$ & fitness of the two predictors & \code{zetaRE}, \code{zetaBR} \\
$\chi$, $\omega$ & intensity of choice, memory & \code{switching.intensity}, \code{switching.memory} \\
$\FE^{i}_t$ & squared one-step forecast error of predictor $i$ & \code{FE_RE}, \code{FE_BR} \\
$s_t$ & predetermined state vector of the original model & \code{meta.re_states} (IRIS \code{xbVector}) \\
$k_x$, $c_x$ & coefficients of the anchored RE forecast & parameters \code{are_x__*} \\
$\T,\K,\Rs,\Z,\Hs,\D$ & state-space matrices of a solved model & \code{sspace(m, 'triangular', false)} \\
$\varepsilon_t$ & transition and measurement shocks & IRIS \code{eList} \\
\bottomrule
\end{longtable}

The solved first-order form of any IRIS model used below is
\begin{equation}\label{eq:sspace}
  \begin{bmatrix} x^f_t \\ x^b_t \end{bmatrix} = \T\, x^b_{t-1} + \K + \Rs\, \varepsilon_t,
  \qquad
  y_t = \Z\, x^b_t + \Hs\, \varepsilon_t + \D,
\end{equation}
with $x^f$ the non-predetermined and $x^b$ the predetermined transition variables (including
IRIS auxiliary lags such as \code{y{-1}}) and $y$ the measurement variables. This is the form
\code{sspace(m, 'triangular', false)} returns; the fast one-period step
(Section~\ref{sec:enginestep}), the anchored coefficients (Section~\ref{sec:anchored}), the
time-varying decomposition (Section~\ref{sec:decompose}) and the Kalman history
(Section~\ref{sec:kalman}) all read it.

\subsection{The sequence of one quarter}\label{sec:timing}

Every stage that simulates the model (forecast, inversion history, density draws, evaluation
runs, the Ramsey loss continuation) runs the same sequence in quarter $t$
(\code{+engine/forecast.m}, \code{+engine/history.m}):

\begin{enumerate}[nosep]
  \item \emph{Parameters.} The belief parameters are set from $\phi_{t-1}$
        (\code{+plm/param_values.m}); the share parameters from
        \code{+engine/shares_for_mode.m}: $n^m_{1,t}$ and $n^m_{2,t} = 1-n^m_{1,t}$.
        Under switching $n^m_{1,t} = n_{1,t-1}$, the share computed at the end of the previous
        quarter \citep[eq.~16]{ozden2025}.
  \item \emph{Solve.} If the parameters differ from those of the last solve, the model is
        assigned and solved through the projection facility (\code{+engine/facility.m},
        call \code{'solve'}). Normally the check of the actual law of motion (ALM, the model
        with the beliefs and shares in force; Section~\ref{sec:facility}) at the end of $t-1$
        has already solved this model, and its solution is reused.
  \item \emph{Simulate.} A one-period simulation of quarter $t$ only, from the databank holding everything up to
        $t-1$, with the shocks of quarter $t$ and the conditions active in $t$
        (\code{+engine/step.m}).
  \item \emph{Read the forecasts.} $\fRE_t(x)$ and $\fBR_t(x)$ are read from the simulated
        values of \code{x_exp1_RE} and \code{x_exp1_BR}. The learners' value is checked
        against the PLM formula with $\phi_{t-1}$ (tolerance $10^{-9}$, error
        \code{hetlearn:engine:timing}).
  \item \emph{Update at the end of $t$} (\code{+engine/update_learning.m}): forecast errors
        $\FE^i_t$ of the forecasts made in $t-1$; fitness and shares $n_{1,t}$; RLS
        candidates for $\phi_t$; the PLM part of the projection facility; the ALM check of the
        model in force in $t+1$; the update of the gain state.
\end{enumerate}

Two assumptions hold throughout:
\begin{itemize}[nosep]
  \item \emph{Anticipated utility} \citep{kreps1998}. When the model is solved in quarter
        $t$, the beliefs $\phi_{t-1}$ and shares $n^m_t$ are treated as permanent. Agents do
        not anticipate that beliefs or shares will change. The model is re-solved when they do.
  \item \emph{Unanticipated shocks and conditions.} Each quarter is a new one-period
        simulation, so every shock and every condition is a surprise in the quarter it hits.
\end{itemize}

\subsection{Information sets}\label{sec:infosets}

Two options date the information behind the two forecasts made in $t$:

\begin{center}
\begin{tabular}{>{\raggedright\arraybackslash}p{5.0cm}>{\raggedright\arraybackslash}p{4.0cm}>{\raggedright\arraybackslash}p{5.2cm}}
\toprule
 & \code{'lagged'} & \code{'current'} \\
\midrule
\code{beliefs.info} (learners) & data up to $t-1$ (default) & data up to $t$ \citep[eq.~12]{ozden2025} \\
\code{expectations.re_info} (RE agents) & $E_{t-1}x_{t+1}$ & $E_t x_{t+1}$ (default) \\
\bottomrule
\end{tabular}
\end{center}

Under both settings of \code{beliefs.info} the beliefs used in $t$ are $\phi_{t-1}$ and the RLS
update at the end of $t$ uses regressors dated $t-1$ or earlier. Only the data that enter the
forecast change. Sections~\ref{sec:beliefsinfo} and~\ref{sec:reinfo} give the formulas.

% =====================================================================
\section{Model requirements and the model package}\label{sec:model}
% =====================================================================

\subsection{What a model must have}\label{sec:modelreq}

\begin{itemize}
  \item \emph{An IRIS model file} that \code{Model.fromFile} loads. IRIS 20230622 is the
        tested release.
  \item \emph{At least one forward-looking transition variable}: a declared transition
        variable with a lead in a transition equation or a substitution. Without one,
        \code{hetlearn.extend_model} stops with \code{hetlearn:extend_model:noLeads}.
  \item \emph{A first-order solution.} The engine solves the model every time the beliefs or
        shares change and needs IRIS's state-space form~\eqref{eq:sspace}. Linear models are
        the tested case (\code{spec.linear = true}). For a nonlinear model the code calls
        \code{spec.sstate} (or \code{sstate(m)}) before each \code{solve}
        (\code{+engine/solve_model.m}), so the path exists, but no package uses it and no test
        covers it. Two parts are linear-only: the anchored RE variant is refused for
        nonlinear models (\code{hetlearn:load_model:anchored}), and the backward simulator
        (Section~\ref{sec:backward}) declines a nonlinear model. Since anchored RE is the
        default, a nonlinear model needs \code{expectations.re_variant = 'model_consistent'}.
  \item \emph{Plain state names.} The anchored RE variant, the fast step and the Kalman
        history read IRIS's transition vector element by element and accept only \code{v} and
        \code{v{-k}}. Log variables (printed as \code{log(v)}) disable the anchored block
        (\code{meta.anchor_ok = false}) and stop the fast step and the Kalman history with an
        error.
  \item \emph{At least as many shocks as observables} for the RE filter's fit check to apply,
        and for \code{history.mode = 'inversion'}.
\end{itemize}

The model's own equations are not restricted otherwise: hybrid terms, lags of any order,
substitutions, measurement equations, the model's own \code{!if} and \code{!switch} blocks,
\code{!for}, \code{!import} and \code{!include} are all accepted (Section~\ref{sec:leads}).

\subsection{The package contract: \texorpdfstring{\codem{model_spec.m}}{model\_spec.m}}\label{sec:spec}

A package is a folder \code{models/<name>/} (or any folder given by path) with a function
\code{model_spec()} that returns a struct. \code{hetlearn.load_spec} evaluates it with the
folder as the current folder, adds \code{spec.folder} and fills the optional fields
(\code{load_spec.m}).

\begin{longtable}{>{\raggedright\arraybackslash}p{2.6cm}>{\raggedright\arraybackslash}p{2.3cm}>{\raggedright\arraybackslash}p{10cm}}
\toprule
Field & Required & Meaning \\
\midrule
\endhead
\code{name} & yes & package name, used in generated file names \\
\code{model_file} & yes & IRIS model file, relative to the folder \\
\code{params} & yes & name of a function in the folder returning a struct of every parameter, \code{std_<shock>} and \code{corr_*} included \\
\code{linear} & yes & passed to \code{Model.fromFile} \\
\code{sstate} & no (\code{''}) & function name or handle \code{m = f(m)} called before every solve; \code{''} for linear models \\
\code{assign} & no (\code{struct()}) & values of the model file's own \code{!if}/\code{!switch} switches \\
\code{database} & for runs & CSV written by \code{databank.toCSV} or a \code{.mat} file (one struct, or the series as variables) \\
\code{filter_range} & for runs & \code{[first, last]} quarter of the RE filter and smoother \\
\code{observables} & for runs & cellstr of measurement variables \\
\code{policy} & no & Ramsey policy section (Section~\ref{sec:policy}) \\
\code{report} & no & \code{labels}, \code{units}, \code{scale} per observable; \code{scale} converts model units to report units \\
\bottomrule
\end{longtable}

Parameters are assigned by name: fields of the parameter struct that are not parameters, shock
standard deviations (\code{std_<shock>}) or correlations (\code{corr_*}) of the model are
ignored (\code{load_model.m}).

\subsection{Package options}

\subsubsection{\texorpdfstring{\codem{model}}{model}}\label{opt:model}
\begin{optspec}
  \item[Default] \code{''} (must be set; \code{hetlearn:run:model} otherwise).
  \item[Values] a folder name under \code{models/}, or the path of any folder that holds a
        \code{model_spec.m} (absolute or relative to the current folder).
  \item[Math] none.
  \item[Assumptions] the package folder is self-contained; generated files go to
        \code{<folder>/generated/}.
  \item[Interactions] \code{hetlearn.evaluate} requires every variant to use the same package
        (\code{hetlearn:evaluate:model}).
  \item[Code] \code{load_spec.m}, \code{run.m}.
  \item[Example]
\begin{lstlisting}
cfg.model = 'sw07';                    % models/sw07/model_spec.m
cfg.model = fullfile('tests', 'toy_nk');  % any folder with a model_spec.m
\end{lstlisting}
\end{optspec}

\subsubsection{\texorpdfstring{\codem{model_params}}{model\_params}}\label{opt:model-params}
\begin{optspec}
  \item[Default] \code{''}: the package's \code{spec.params}. (Project baseline:
        \code{'sw07_params_original'}.)
  \item[Values] \code{''}, or a valid MATLAB name of a parameter function that is a
        \code{.m} file in the package folder. Another name stops the load with
        \code{hetlearn:load_model:params}.
  \item[Math] the parameter vector $\theta$ of the original model.
  \item[Assumptions] parameters are given, never estimated by the toolbox.
  \item[Interactions] the chosen function replaces \code{spec.params} everywhere downstream:
        the generator's collapse test (when it runs), the anchored coefficients, the policy
        discount, the RE
        filter (whose cache key hashes the parameter values, so each set has its own caches),
        and \code{hetlearn.evaluate}, which keeps parameter sets apart when it shares filters.
  \item[Code] \code{load_model.m}, \code{re_filter.m}, \code{evaluate.m}.
  \item[Example]
\begin{lstlisting}
cfg.model_params = 'sw07_params_original';   % published SW07 mode
cfg.model_params = '';                        % back to spec.params
\end{lstlisting}
\end{optspec}

% =====================================================================
\section{The extended model}\label{sec:generator}
% =====================================================================

\code{hetlearn.extend_model(spec, plm)} reads the package's model file and writes an extended
file in \code{<folder>/generated/}. \code{hetlearn.load_model} calls it before every run; it
regenerates only when something changed (Section~\ref{sec:gencache}). The user never edits the
generated file.

\subsection{Lead detection}\label{sec:leads}

Code: \code{extend_model.m} (functions \code{mask_code}, \code{parse_structure},
\code{find_occurrences}, \code{active_leads}, \code{check_pseudofunctions}).

\begin{itemize}
  \item Leads are detected only on \emph{declared transition variables}, only inside
        \code{!transition_equations} (alias \code{!equations}) and \code{!substitutions}, and
        only in the dynamic part of an equation (the text before \code{!!}). A lead in a
        measurement equation is normalized to IRIS notation but not replaced.
  \item Comments (\texttt{\%}, \texttt{\%\{ \%\}}, the rest of a line after \code{...}) and
        quoted labels are never changed. IRIS 20230622 accepts only \texttt{\%} and
        \texttt{\%\{ \%\}}, so the delimiters \code{//}, \code{/*}, \code{*/} become
        \texttt{\%\%}, \texttt{\%\{}, \texttt{\%\}} with the comment text unchanged.
  \item Accepted lead notations, whitespace tolerant, for an integer $k\ge 1$:
        \code{x{+k}}, \code{x{k}}, \code{x{t+k}}, \code{x[+k]}, \code{x(+k)}. Non-IRIS lags
        (\code{x[-k]}, \code{x{t-k}}, \code{x(-k)}) are normalized to \code{x{-k}}. Because IRIS
        cannot load the non-IRIS forms, a reference file \code{<name>_norm.model} with only the
        subscripts normalized is written and used by the collapse test.
  \item The pseudo-functions \code{diff}, \code{dot}, \code{roc}, \code{pct}, \code{difflog},
        \code{movsum}, \code{movprod}, \code{movavg}, \code{movgeom} with a lead inside or a
        forward window stop the generator (\code{hetlearn:extend_model:pseudoLead}), naming the
        equation. Rewrite them with explicit leads.
  \item Files with \code{!for}, \code{!import} or \code{!include} are expanded first with the
        IRIS preparser (\code{parser.Preparser.parse} with \code{spec.assign}); comments are lost
        and the model's own \code{!if} blocks are resolved in that case.
  \item A file with its own \code{!if}/\code{!switch} blocks keeps every branch, but the set
        $\mathcal{F}$ of forward-looking variables and their maximum leads are detected on the
        branches that \code{spec.assign} selects (the preparser resolves a copy of the file).
        A variable whose leads sit only in inactive branches is not forward-looking and keeps
        \code{x{+k}} there.
  \item Generated names that collide with declared names or with \code{spec.assign} fields stop
        the generator (\code{hetlearn:extend_model:nameCollision}).
\end{itemize}

For each $x\in\mathcal{F}$ the generator records the maximum lead $K_x$ (\code{meta.max_lead}).
$\mathcal{F}$ is listed in declaration order (\code{meta.forward_vars}).

\subsection{Rewriting and the expectation block}\label{sec:block}

Every lead \code{x{+k}} with $x\in\mathcal{F}$ becomes \code{x_expk}. The generator then appends
an expectation block (new declaration and equation sections at the end of the file; IRIS allows
repeated sections, so the user's sections are not touched). For the default AR(1) PLM the block
of one variable reads (\code{make_block} in \code{extend_model.m}):

\begin{lstlisting}
!transition_variables
    "Aggregate expectation of x{+1}" x_exp1
    "Aggregate expectation of x{+2}" x_exp2              % only if K_x >= 2
    !if br_x
        "RE expectation of x{+1}" x_exp1_RE
        "BR (adaptive learning) expectation of x{+1}" x_exp1_BR
    !end
!parameters
    n1_x, n2_x, alpha_x, beta_x, are_x__const, are_x__<state>, ...
!transition_equations
    !if br_x
        x_exp1 = n1_x*x_exp1_RE + n2_x*x_exp1_BR;
        !if re_anchored
            x_exp1_RE = are_x__const + are_x__<s1>*s1 + ...;
        !else
            x_exp1_RE = x{+1};
        !end
        !if br_info_current
            x_exp1_BR = alpha_x + beta_x*x;
        !else
            x_exp1_BR = alpha_x + beta_x*(alpha_x + beta_x*x{-1});
        !end
    !else
        x_exp1 = x{+1};
    !end
    x_exp2 = x_exp1{+1};                                  % only if K_x >= 2
\end{lstlisting}

In equations, for $x\in\mathcal{L}$ (\code{br_x = true}),
\begin{align}
  \Eagg_t x_{t+1} &= n^m_{1,t}\,\fRE_t(x) + n^m_{2,t}\,\fBR_t(x),
      \qquad n^m_{2,t} = 1 - n^m_{1,t}, \label{eq:agg}\\
  \Eagg_t x_{t+k} &= E_t\bigl[\Eagg_{t+1} x_{t+k}\bigr], \qquad k = 2,\dots,K_x, \label{eq:aggk}
\end{align}
and for $x\in\mathcal{F}\setminus\mathcal{L}$ (\code{br_x = false}) the expectation is the
model-consistent one, $\Eagg_t x_{t+1} = E_t x_{t+1}$, as in the original model. The engine always
sets $n_{2}=1-n_{1}$ (\code{+engine/facility.m}, call \code{'params'}).

Equation~\eqref{eq:aggk} is the IRIS lead of an endogenous variable: the $k$-step aggregate
expectation is the model-consistent forecast of next quarter's $(k-1)$-step aggregate
expectation. Heterogeneity therefore enters at the one-step level, and longer expectations are
built by iterating it. A model with a lead of two or more keeps a lead even when every
forward-looking variable learns.

\paragraph{Derived regressors.} Each \code{beliefs.plm.derived.<name>} adds the transition
variable \code{<name>} with the equation \code{<name> = <expression>;}, outside every toggle
(Section~\ref{opt:beliefs-plm-derived}).

\paragraph{Log variables.} If $x$ is a log variable (\code{!log_variables}), the aggregate
expectations \code{x_exp1}, \dots, \code{x_expK} are declared as log variables too. With
\code{!all_but} a warning (\code{hetlearn:extend_model:logAllBut}) says that they inherit the
log status.

\paragraph{The lagged RE file.} With \code{expectations.re_info = 'lagged'} the generator writes a
second file whose RE expectation uses information up to $t-1$ (Section~\ref{sec:reinfo}).

\subsection{Model-file toggles}\label{sec:toggles}

The block sits behind IRIS \code{!if} switches. \code{hetlearn.load_model} sets them from the
options through the \code{'assign'} argument of \code{Model.fromFile}:

\begin{center}
\begin{tabular}{>{\raggedright\arraybackslash}p{4.4cm}>{\raggedright\arraybackslash}p{4.0cm}>{\raggedright\arraybackslash}p{6.4cm}}
\toprule
Switch & Set from & Effect when true \\
\midrule
\code{br_<x>} (one per $x\in\mathcal{F}$) & \code{expectations.variables} & $x$'s expectation is the RE/learner mix~\eqref{eq:agg}; false: $E_t x_{t+1}$ \\
\code{re_anchored} (global) & \code{expectations.re_variant} & anchored RE forecast (Section~\ref{sec:anchored}) instead of \code{x{+1}} \\
\code{br_info_current} (global) & \code{beliefs.info} & learners' forecast on data up to $t$ \\
\code{ramsey} (global) & \code{policy.mode} & the Ramsey loss replaces the policy rule (Section~\ref{sec:policy}) \\
\bottomrule
\end{tabular}
\end{center}

The model's own switches come from \code{spec.assign}. \code{expectations.re_info} is not a
switch: it selects another generated file.

\subsection{Parameters of the extended model}\label{sec:genparams}

Every generated parameter is declared unconditionally (an unused parameter is harmless):

\begin{itemize}[nosep]
  \item shares \code{n1_x}, \code{n2_x} for every $x\in\mathcal{F}$;
  \item the belief parameters of every PLM equation of the generator's PLM system
        (Section~\ref{sec:plmnames}). The generator does not know the scenario's learning set,
        so it gives \emph{every} forward-looking variable a primary PLM, built from
        \code{beliefs.plm}, plus the auxiliary PLMs that close the system;
  \item the anchored coefficients \code{are_x__const}, \code{are_x__<v>} and
        \code{are_x__<v>_l<k>}, one per state of the original model (Section~\ref{sec:anchored});
        left out when \code{meta.anchor_ok} is false;
  \item with \code{spec.policy}: the loss weights \code{w_<k>}, \code{ramsey_disc_on} and
        \code{ramsey_disc} (Section~\ref{sec:policy}).
\end{itemize}

At load time every \code{n1_x} is 1, every \code{n2_x} 0, every belief parameter 0, and the
anchored coefficients take their computed values (\code{load_model.m}). The loaded model is then
the RE model (Section~\ref{sec:collapse}); the engine assigns beliefs and shares each quarter.

\subsection{Generated files and the cache}\label{sec:gencache}

\begin{longtable}{>{\raggedright\arraybackslash}p{5.6cm}>{\raggedright\arraybackslash}p{9.4cm}}
\toprule
File (in \code{generated/}) & Content \\
\midrule
\endhead
\code{<name>_ext_<tag>.model} & the extended model (generator version 4) \\
\code{<name>_ext_<tag>_meta.mat} & \code{meta}: source file and MD5, normalized PLM spec and key, $\mathcal{F}$, maximum leads, generated names, toggles, PLM system, states of the original model, anchored parameter names, policy section, collapse-test results, generator version \\
\code{<name>_ext_<tag>_relag.model}, \code{_relag_meta.mat} & the lagged RE file (version 5) \\
\code{<name>_norm.model} & the source with subscripts in IRIS notation (only when needed) \\
\bottomrule
\end{longtable}

The tag is \code{ar1} when every PLM equation is an AR(1) with constant on the own lag and there
are no derived regressors; otherwise \code{p} followed by the first 8 hexadecimal digits of the
MD5 of the normalized PLM spec. A cached file is reused only if the generator version, the
source MD5, the PLM key, the \code{spec.assign} key and the policy key all match; otherwise it is
rewritten.

\subsection{The collapse test}\label{sec:collapse}

At $n_1 = 1$, $n_2 = 0$ and zero beliefs, equation~\eqref{eq:agg} gives $\Eagg_t x_{t+1} =
\fRE_t(x)$. With \code{re_anchored = false} this is \code{x{+1}}, so the extended model is the
original model with renamed leads. With \code{re_anchored = true} and the anchored coefficients
of the original model, it is the RE forecast at the current state, which on the RE path is again
$E_t x_{t+1}$. Every generated file must pass (\code{collapse_test} in \code{extend_model.m}):

\begin{enumerate}[nosep]
  \item every \code{br_x = true}, \code{re_anchored = false};
  \item every \code{br_x = false};
  \item every \code{br_x = true}, \code{re_anchored = true}, coefficients from
        \code{hetlearn.re_anchor_coefficients} (only when the file has the anchored block);
  \item every \code{br_x = true}, \code{br_info_current = true}.
\end{enumerate}

Each pass compares the impulse responses of every variable of the reference model to every
shock over 40 quarters; the largest absolute difference must be below $10^{-10}$
(\code{hetlearn:extend_model:collapseFailed} otherwise). The results are stored in
\code{meta.collapse_maxdiff} (passes 1--3) and \code{meta.collapse_info_maxdiff} (pass 4). The
lagged file has its own passes (Section~\ref{sec:reinfo}).

\subsection{Generator call options}\label{opt:extend-model}

\code{meta = hetlearn.extend_model(spec, plm, ...)} takes the call options \code{'force'}
(default false), which regenerates the file even when the cache matches, and \code{'validate'}
(default true), which runs the collapse test. \code{'re_info'} selects the lagged file and is set
by \code{hetlearn.load_model}. A direct call with a vector autoregression (VAR) over
\code{'learning'} gives every forward-looking variable the VAR, because the learning set is resolved only by
\code{hetlearn.load_model} (Section~\ref{opt:beliefs-plm-var}).

% =====================================================================
\section{Expectations: the aggregate and the modes}\label{sec:expectations}
% =====================================================================

Two options decide which expectations are heterogeneous and how the population is split: the
learning set $\mathcal{L}$, and the mode, which fixes the RE share $n^m_{1,t}$ in the aggregate
expectation~\eqref{eq:agg}.

\subsection{Options}

\subsubsection{\texorpdfstring{\codem{expectations.mode}}{expectations.mode}}\label{opt:expectations-mode}
\begin{optspec}
  \item[Default] \code{'rational'}.
  \item[Values] \code{'rational'}, \code{'learning'}, \code{'fixed_shares'}, \code{'switching'}.
  \item[Math] the RE share in the model in quarter $t$ (\code{+engine/shares_for_mode.m}):
\begin{equation}\label{eq:modes}
  n^m_{1,t} =
  \begin{cases}
    1 & \text{\codem{'rational'}}\\
    0 & \text{\codem{'learning'}}\\
    1 - \bar n_2 & \text{\codem{'fixed_shares'}, } \bar n_2 = \text{\codem{shares.fixed_n2}}\\
    n_{1,t-1} & \text{\codem{'switching'}, the share of the variable's group}
  \end{cases}
\end{equation}
        In every mode the end-of-quarter update still computes the fitness and the
        discrete-choice shares $n_{1,t}$ (Section~\ref{sec:switching}) and, under
        \code{beliefs.update = 'rls'}, the beliefs; outside \code{'switching'} they are
        diagnostics that do not enter the model.
  \item[Assumptions] \code{'rational'}: every agent is RE, so the run is exactly the RE model
        of the scenario (the collapse argument of Section~\ref{sec:collapse}); with the default
        unanticipated shocks it equals one IRIS \code{simulate} of the RE model over the
        horizon (test \code{test_engine} E1). \code{'switching'}: shares follow the
        discrete-choice rule; the share in force in $t$ was computed at the end of $t-1$.
  \item[Interactions] \code{'fixed_shares'} needs \code{shares.fixed_n2}
        (\code{hetlearn:options:missing}). The projection facility skips its ALM check when
        every learners' model share is 0 (\code{'rational'}, or fixed shares of 0). The density
        fast path needs \code{'rational'} (Section~\ref{sec:density}); \code{hetlearn.decompose}
        with \code{method = 'auto'} uses IRIS contributions for \code{'rational'}
        (Section~\ref{sec:decompose}). With an empty learning set the mode has no effect. A
        share kept by the facility (Section~\ref{sec:keptshares}) overrides~\eqref{eq:modes}
        for one quarter.
  \item[Code] \code{+engine/shares_for_mode.m}, \code{+engine/update_learning.m}.
  \item[Example]
\begin{lstlisting}
cfg.expectations.mode = 'switching';
cfg.expectations.variables = {'pinf'};
\end{lstlisting}
\end{optspec}

\subsubsection{\texorpdfstring{\codem{expectations.variables}}{expectations.variables}}\label{opt:expectations-variables}
\begin{optspec}
  \item[Default] \code{{}} (no learning variable: the run is RE whatever the mode).
  \item[Values] a cellstr of forward-looking transition variables, or \code{'all'} for every
        $x\in\mathcal{F}$. A char is turned into a one-element cellstr.
  \item[Math] the learning set $\mathcal{L}$; \code{br_x = true} for $x\in\mathcal{L}$.
  \item[Assumptions] each $x\in\mathcal{L}$ gets its own learners' forecast and its own PLM;
        variables outside $\mathcal{L}$ keep $E_t x_{t+1}$, the model-consistent expectation
        of the heterogeneous economy (the anchored variant applies only inside the
        \code{br_x} blocks).
  \item[Interactions] a name that is not forward-looking stops the load
        (\code{hetlearn:load_model:notForward}; \code{info.meta.forward_vars} lists the
        candidates). $\mathcal{L}$ is ordered as declared in the model, not as listed. It
        defines the share groups (\code{shares.structure}), the default of
        \code{gain.koz.master}, the members of a VAR over \code{'learning'}, and, with
        \code{expectations.re_info = 'lagged'}, the RE(t-1) model of Section~\ref{sec:reinfo}
        (an empty set is then an error). With anchored RE and every forward-looking variable of maximum lead 1 in
        $\mathcal{L}$, the model has no leads (Section~\ref{sec:backward}).
  \item[Code] \code{load_model.m} (\code{learning_vars}), \code{extend_model.m}.
  \item[Example]
\begin{lstlisting}
cfg.expectations.variables = {'pinf', 'w'};
cfg.expectations.variables = 'all';          % every forward-looking variable
\end{lstlisting}
\end{optspec}

% =====================================================================
\section{RE agents}\label{sec:re}
% =====================================================================

\subsection{Model-consistent RE agents}\label{sec:mc}

With \code{re_anchored = false} the generated equation is \code{x_exp1_RE = x{+1}}:
\begin{equation}\label{eq:mc}
  \fRE_t(x) = E_t\bigl[x_{t+1} \,\big|\, \phi_{t-1},\, n^m_{t}\bigr],
\end{equation}
the rational forecast in the heterogeneous economy whose ALM IRIS solves in quarter $t$, under
the anticipated-utility assumption of Section~\ref{sec:timing}. RE agents know the model, the
learners' PLM coefficients and the shares.

\subsection{Anchored RE agents}\label{sec:anchored}

Following \citet[sec.~2.3.1]{ozden2025}, anchored RE agents use the forecast of the pure-RE
model, evaluated at the current state of the heterogeneous economy. Write the RE solution of the
original model at its parameters as~\eqref{eq:sspace} and let $s_t = x^b_t$ be its predetermined
state vector (IRIS \code{xbVector}: current values of variables that appear with lags, auxiliary
lags such as \code{c{-1}}, the exogenous processes). Then $E^{\RE}_t x_{t+1} = \K_x + \T_{x\cdot}
s_t$, and
\begin{equation}\label{eq:anchored}
  \fRE_t(x) = k_x + c_x\transp s_t, \qquad k_x = \K_x,\quad c_x = \T_{x\cdot}\transp ,
\end{equation}
with $s_t$ taken from the heterogeneous economy. The generated equation is
\code{x_exp1_RE = are_x__const + sum_j are_x__<s_j>*s_j}, where a state \code{v} gets the
parameter \code{are_x__v} and an auxiliary lag \code{v{-k}} the parameter \code{are_x__v_lk}
(\code{+plm/anchor_names.m}).

The coefficients come from \code{hetlearn.re_anchor_coefficients}
(\code{re_anchor_coefficients.m}, through \code{anchor_coefficients.m} with lag 0 and horizon 1)
on the original model at \code{spec.params} (or \code{model_params}). They are checked by
unit-perturbation simulation: one IRIS \code{simulate} of quarter $t+1$ with zero shocks from a
databank with $n_b+1$ columns, the state at 0 and at each unit vector, must return $k_x$ and
$k_x + c_{x,j}$ (relative tolerance $10^{-9}$; \code{hetlearn:anchor:check}).
\code{hetlearn.load_model} recomputes them at every load, so they follow the parameters.

\paragraph{Assumptions.} Anchored agents believe the economy is the pure-RE economy and keep its
reduced form; they do not internalize the learners or the shares. At $n_1=1$ the anchored model
is the RE model (collapse pass 3). The forecast~\eqref{eq:anchored} uses $s_t$, which holds
values dated $t$, so it is determined jointly with the other variables dated $t$. Under Ramsey
policy the coefficients come from the rule-based RE model (Section~\ref{sec:policy}).

\subsection{Existence and models without leads}\label{sec:backward}

When every forward-looking variable is in $\mathcal{L}$, every maximum lead is 1 and the RE agents
are anchored, the ALM has no leads (\code{get(m, 'maxLead') = 0}). A linear model without leads
has a unique solution for any beliefs and shares whenever its system in the variables dated $t$
is nonsingular,
possibly explosive. IRIS rejects explosive solutions even in purely backward models, so when IRIS
fails on such a model, \code{+engine/solve_model.m} builds a backward system
(\code{+engine/backward_simulate.m}) from \code{system(m)}:
\begin{gather}
  A\begin{bmatrix} x^f_{t+1} \\ x^b_t\end{bmatrix} + B\begin{bmatrix} x^f_t \\ x^b_{t-1}\end{bmatrix} + C + D\varepsilon_t = 0,
  \qquad F y_t + G x^b_t + H + J\varepsilon_t = 0,\\
  M = \bigl[B_{\cdot,1:n_f},\ A_{\cdot,n_f+1:n}\bigr],\qquad
  \begin{bmatrix} x^f_t \\ x^b_t\end{bmatrix} = -M^{-1}\bigl(B_b x^b_{t-1} + C + D\varepsilon_t\bigr),\\
  y_t = -F^{-1}\bigl(G x^b_t + H + J\varepsilon_t\bigr),
\end{gather}
with $B_b = B_{\cdot,n_f+1:n}$ and no leads meaning $A_{\cdot,1:n_f}=0$. The roots are the
eigenvalues of $\Phi$, the rows of $-M^{-1}B_b$ that belong to $x^b$. What happens next depends on
the projection facility (Section~\ref{sec:facility}): under \code{'none'} the quarter is simulated
with this system and an \code{explosive} event records the largest root; under \code{'alm'} and
\code{'plm_alm'} a backward system with a root above 1 in modulus is an unstable ALM. The backward
system equals IRIS \code{simulate} whenever IRIS has a solution (\code{test_anchored} A3). A model
with a lead of two or more keeps \code{x_exp2 = x_exp1{+1}} and so keeps a lead.

\subsection{Option}

\subsubsection{\texorpdfstring{\codem{expectations.re_variant}}{expectations.re\_variant}}\label{opt:expectations-re-variant}
\begin{optspec}
  \item[Default] \code{'anchored'}.
  \item[Values] \code{'anchored'}, \code{'model_consistent'}.
  \item[Math] \code{'anchored'}: equation~\eqref{eq:anchored}; with
        \code{re_info = 'lagged'}, equation~\eqref{eq:anchoredlag}. \code{'model_consistent'}:
        equation~\eqref{eq:mc}; with \code{'lagged'}, the chain~\eqref{eq:chain}.
  \item[Assumptions] see Sections~\ref{sec:mc}--\ref{sec:anchored}. Under the Taylor rule the
        two variants give the same RE path (both are the RE model at $n_1 = 1$).
  \item[Interactions] the variant applies to the learning set only. \code{'anchored'} needs a
        linear model with plain state names (\code{hetlearn:load_model:anchored});
        \code{'model_consistent'} works for any model IRIS solves. With anchored RE and every
        lead-1 forward-looking variable learning, the model has no leads
        (Section~\ref{sec:backward}). For an RE Ramsey benchmark see
        Section~\ref{opt:policy-mode}. With many
        learning variables a model-consistent ALM can lose its stable solution: on SW07 with
        all twelve forward-looking variables learning, such a run completes only because the
        facility keeps the model shares (look for \code{projection_alm_shares} in
        \code{events.csv}). The RE filter of the history uses the scenario's variant (its own
        cache file).
  \item[Code] \code{extend_model.m} (\code{anchor_states}, \code{anchor_expr}),
        \code{re_anchor_coefficients.m}, \code{anchor_coefficients.m}, \code{load_model.m}.
  \item[Example]
\begin{lstlisting}
cfg.expectations.re_variant = 'model_consistent';
cfg.policy.mode = 'ramsey';                 % RE Ramsey benchmark
\end{lstlisting}
\end{optspec}

% =====================================================================
\section{Information timing}\label{sec:info}
% =====================================================================

\subsection{The learners' information set}\label{sec:beliefsinfo}

The learners' forecast of $x_{t+1}$ is built by iterating the PLMs forward
(\code{+plm/forecast_expr.m} writes it as text for the model file, \code{+plm/forecast_value.m}
evaluates it numerically). With the PLM equation of variable $v$,
$v_t = \alpha_v + \sum_j \beta_{v,j}\, u_{j,t-l_j}$ (Section~\ref{sec:plm}), define
\begin{equation}\label{eq:recursion}
  \mathcal{E}(v,h) = \alpha_v + \sum_j \beta_{v,j}\; \mathcal{T}(u_j,\, h - l_j),
  \qquad
  \mathcal{T}(u,k) =
  \begin{cases}
    u_{t+k} & k \le k_{\max}\\
    \mathcal{E}(u,k) & k > k_{\max}
  \end{cases}
\end{equation}
and $\fBR_t(x) = \mathcal{E}(x,1)$, all with the coefficients $\phi_{t-1}$. The information set
fixes $k_{\max}$: $-1$ under \code{'lagged'} (data up to $t-1$), $0$ under \code{'current'} (data
up to $t$). Since every regressor lag is at least 1, under \code{'current'} $\mathcal{E}(x,1)$
never recurses. Examples:

\begin{center}
\begin{tabular}{lll}
\toprule
PLM & \code{'lagged'} & \code{'current'} \\
\midrule
AR(1) & $\alpha + \beta(\alpha + \beta x_{t-1})$ & $\alpha + \beta x_t$ \\
AR(2) & $\alpha + \beta_1(\alpha + \beta_1 x_{t-1} + \beta_2 x_{t-2}) + \beta_2 x_{t-1}$ & $\alpha + \beta_1 x_t + \beta_2 x_{t-1}$ \\
VAR(1) & $a + B(a + B X_{t-1})$ & $a + B X_t$ \\
\bottomrule
\end{tabular}
\end{center}

Under \code{'current'} the forecast contains $x_t$ and is solved jointly with it
\citep[eq.~12]{ozden2025}. The beliefs used in $t$ are $\phi_{t-1}$ under both settings, and the
RLS update at the end of $t$ uses the same regressors (dated $t-1$ or earlier), so the beliefs do
not depend on the information set; the forecasts, and through them the fitness and the shares,
do. Under \code{'current'} the error that enters fitness at $t$ compares $x_t$ with
$\fBR_{t-1}(x) = \alpha_{t-2} + \beta_{t-2} x_{t-1}$ (AR(1)).

\subsubsection{\texorpdfstring{\codem{beliefs.info}}{beliefs.info}}\label{opt:beliefs-info}
\begin{optspec}
  \item[Default] \code{'lagged'}.
  \item[Values] \code{'lagged'}, \code{'current'}.
  \item[Math] $k_{\max} = -1$ or $0$ in~\eqref{eq:recursion}.
  \item[Assumptions] \code{'lagged'}: learners observe data up to $t-1$ when they forecast in
        $t$. \code{'current'}: they observe $x_t$ and the other regressors dated $t$, as in
        \citet{ozden2025}.
  \item[Interactions] the generated file holds both forms; \code{br_info_current} selects one
        (a file from generator version 3 or older with \code{'current'} is
        \code{hetlearn:load_model:info}). Under \code{gain.koz.error = 'forecast'} the error that
        drives the gain is $x_t - \fBR_{t-1}(x)$, a two-step error under \code{'lagged'}. In
        the history stage the forecast reads the history databank at $t+k$ with $k\le k_{\max}$;
        in the Kalman history it reads the real-time filtered values.
  \item[Code] \code{+plm/forecast_expr.m}, \code{+plm/forecast_value.m}, \code{load_model.m}.
  \item[Example]
\begin{lstlisting}
cfg.beliefs.info = 'current';               % alpha + beta*x_t
\end{lstlisting}
\end{optspec}

\subsection{The RE agents' information set}\label{sec:reinfo}

With \code{expectations.re_info = 'lagged'} the RE expectation of every $x\in\mathcal{L}$ is
$E_{t-1}x_{t+1}$, formed with information up to $t-1$; expectations of other variables stay
$E_t$. This is the RE counterpart of \code{beliefs.info = 'lagged'}: both forecasts made in $t$
then use data up to $t-1$, and both fitness errors are two-step errors. Every RE object of a
lagged scenario refers to the \emph{RE(t-1) model of its learning set}: the original model with
$E_{t-1}x_{t+1}$ in place of $E_t x_{t+1}$ for $x\in\mathcal{L}$, solved as an RE fixed point. That
model is the RE benchmark (mode \code{'rational'}), the anchored agents' model, the model of the RE
filter and smoother, and the source of the RE-moment initial beliefs.

\paragraph{Model-consistent form.} The lagged file (\code{make_block_lagged} in
\code{extend_model.m}) writes, inside each \code{br_x} block,
\begin{equation}\label{eq:chain}
  x^{(1)}_t = x_{t+1}, \qquad x^{(2)}_t = x^{(1)}_{t+1}, \qquad \fRE_t(x) = x^{(2)}_{t-1},
\end{equation}
(\code{x_REt1 = x{+1}}, \code{x_REt2 = x_REt1{+1}}, \code{x_exp1_RE = x_REt2{-1}}), so
$x^{(1)}_t = E_tx_{t+1}$, $x^{(2)}_t = E_tx_{t+2}$ and $\fRE_t(x) = E_{t-1}x_{t+1}$. Every lead stays
at $+1$ (IRIS's own \code{x{+2}} would add a lead element the fast step cannot read). With
model-consistent agents, $x^{(2)}_{t-1}$ comes from the model in force at $t-1$, so $\fRE_t(x)$ is
predetermined in $t$.

\paragraph{Anchored form.} Write the solution of the model-consistent lagged model at $n_1=1$ and
zero beliefs as~\eqref{eq:sspace}, with $x^b = (s, z, o)$: $s$ the original model's states, $z$ the
chain states $z_{x,t} = x^{(2)}_t$, $o$ any other state (PLM lags, derived regressors). Iterating
two quarters from the state at $t-1$,
\begin{equation}
  E_{t-1}x_{t+1} = \K_x + \T_{x\cdot}\bigl(\K_b + \T_b x^b_{t-1}\bigr) = k + C_s s_{t-1} + C_z z_{t-1},
\end{equation}
where the loadings on $o$ must be below $10^{-12}$ (\code{hetlearn:anchor:lagged}). In the
heterogeneous economy the chain state $z_{t-1}$, the RE agents' own forecast of $x_{t+1}$ formed
with information up to $t-1$, is $\fRE_t$ itself. Stacking the learning set,
$f_t = k + C_s s_{t-1} + C_z f_t$, so
\begin{equation}\label{eq:anchoredlag}
  \fRE_t = (I - C_z)^{-1}\bigl(k + C_s s_{t-1}\bigr),
\end{equation}
a linear function of states dated $t-1$ and earlier (\code{re_anchor_coefficients.m},
\code{lagged_coefficients}; a near-singular $I-C_z$, reciprocal condition below $10^{-12}$, is an
error). The generated equation dates every state one quarter earlier:
\code{x_exp1_RE = are_x__const + sum_j are_x__<s_j>*s_j{-1}}, a lag-$k$ state becoming
\code{v{-k-1}}. On the RE(t-1) path $z_{t-1}$ satisfies the same fixed point, so at $n_1 = 1$ the
anchored and model-consistent lagged models coincide; \code{hetlearn.load_model} checks it at
every load (impulse responses of every common variable, 40 quarters, below $10^{-10}$,
\code{info.re_info_check}). The coefficients depend on the learning set, so the model-consistent
lagged model is loaded and solved for every anchored lagged scenario.

\paragraph{Collapse test of the lagged file.} (\code{collapse_test_lagged}.) (1) every
\code{br_x = false} reproduces the original model; (2) with every lead-1 forward-looking variable
learning and \code{re_anchored = false}, the RE(t-1) model of that set must solve; (3) the same
with \code{re_anchored = true} and the lagged coefficients reproduces (2); (4) (2) with
\code{br_info_current = true} reproduces (2). If (2) has no unique stable solution for the whole
set, passes (2)--(4) run one variable at a time. \code{meta.collapse_maxdiff(1)} is NaN, since the
lagged file differs from the original by construction.

\paragraph{History.} The RE filter of a lagged scenario keeps the scenario's learning set and
variant (own cache), and also returns the one-step predictions $x_{t|t-1}$ (\code{H.predicted}).
Under \code{history.re_states = 'smoothed'} the RE forecast of the history is the smoothed
\code{x_exp1_RE} at $t$, which equals the two-step forecast from the smoothed state at $t-1$.
Under \code{'filtered'} it is \code{H.predicted.x_exp1_RE(t)}
$= E[x_{t+1}\mid\text{data to } t-1]$, the filter's two-step prediction
(\code{+engine/history_db.m}); the learners' data and the realized values stay $x_{t|t}$.

\subsubsection{\texorpdfstring{\codem{expectations.re_info}}{expectations.re\_info}}\label{opt:expectations-re-info}
\begin{optspec}
  \item[Default] \code{'current'}. (Project baseline: \code{'lagged'}.)
  \item[Values] \code{'current'}, \code{'lagged'}.
  \item[Math] \code{'current'}: $\fRE_t(x)$ from~\eqref{eq:mc} or~\eqref{eq:anchored}.
        \code{'lagged'}: from~\eqref{eq:chain} or~\eqref{eq:anchoredlag}.
  \item[Assumptions] under \code{'lagged'} the RE agents see the same data as lagged learners;
        the RE benchmark becomes the RE(t-1) model of $\mathcal{L}$, not the original model.
  \item[Interactions] \code{'lagged'} needs a non-empty learning set
        (\code{hetlearn:options:value}), learning variables with a maximum lead of 1
        (\code{hetlearn:load_model:reInfo}), and \code{policy.mode = 'rule'}
        (\code{hetlearn:options:unsupported}). It changes the RE filter, the RE forecasts of the
        history (\code{H.predicted} under \code{'filtered'}), the anchored coefficients, the
        RE-moment initial beliefs and the Kalman history (RE forecasts from the predictions
        $x_{t|t-1}$). \code{hetlearn.evaluate} gives a lagged variant its own real-time filters
        (the group key adds \code{re_info}, the learning set and \code{re_variant}). Compare a
        lagged learning run with a lagged RE run.
  \item[Code] \code{extend_model.m} (\code{make_block_lagged}, \code{collapse_test_lagged}),
        \code{re_anchor_coefficients.m}, \code{load_model.m}, \code{re_filter.m},
        \code{+engine/history_db.m}.
  \item[Example]
\begin{lstlisting}
cfg.expectations.re_info = 'lagged';        % E_{t-1} x_{t+1} for the learning set
cfg.beliefs.info = 'lagged';                % learners on the same data
\end{lstlisting}
\end{optspec}

% =====================================================================
\section{Perceived laws of motion}\label{sec:plm}
% =====================================================================

\subsection{The PLM equation}\label{sec:plmeq}

Each learning variable $x$ has a primary PLM, and every regressor without its own PLM gets an
auxiliary one, so that the system is closed (\code{+plm/build_system.m}). A PLM equation for a
target $v$ is
\begin{equation}\label{eq:plm}
  v_t = \alpha_v + \sum_{j=1}^{J_v} \beta_{v,j}\, u_{j,t-l_j} + e_{v,t},
  \qquad l_j \ge 1,
\end{equation}
with regressor vector and coefficients
\begin{equation}\label{eq:zphi}
  z_{v,t-1} = \bigl(1,\, u_{1,t-l_1},\, \dots,\, u_{J,t-l_J}\bigr)\transp, \qquad
  \phi_v = \bigl(\alpha_v,\, \beta_{v,1},\, \dots,\, \beta_{v,J}\bigr)\transp .
\end{equation}
Without a constant the leading 1 and $\alpha_v$ are dropped. Every regressor is dated $t-1$ or
earlier, whatever the information set; \code{beliefs.info} only changes which data enter the
forecast~\eqref{eq:recursion}. Regressors are lags of transition variables of the extended
model, including derived regressors (Section~\ref{opt:beliefs-plm-derived}).

\paragraph{Closing the system.} Starting from the learning variables, the builder walks a queue:
a VAR member gets the VAR equation; a variable with \code{beliefs.plm.by_var.<v>} gets that
specification; a learning variable otherwise gets \code{beliefs.plm.default}; any other regressor
gets an AR(1) with constant. Each new regressor joins the queue until every regressor has an
equation. Equations added this way are \emph{auxiliary} (\code{isAux = true}); they are learned
by the same RLS and recorded like the primary ones, but no agent uses their forecast except
inside the recursion~\eqref{eq:recursion}.

\paragraph{Parameter names.}\label{sec:plmnames} An AR(1) with constant on the own first lag has
\code{alpha_x} and \code{beta_x}. Any other PLM has \code{alpha_x} (if it has a constant) and one
\code{beta_x__<u>_l<l>} per regressor, e.g.\ \code{beta_pinf__ygap_l1}. Coefficient vectors are
ordered like the names: constant first, then the regressors in the order given (a VAR: lag first,
then members in the order listed). The same names are the model parameters, the fields of
\code{res.learning.<target>} and the rows of \code{learning.csv}.

\paragraph{Consistency of the generated file and the run.} The generator builds its PLM system
with every forward-looking variable as a primary target (Section~\ref{sec:genparams}); the run
builds it with the scenario's learning set. \code{hetlearn.load_model} compares, for every
$x\in\mathcal{L}$, the forecast text the run's system gives with the one in the active branch of
the generated file, and checks that every belief parameter is declared
(\code{hetlearn:load_model:plmMismatch} otherwise). The two differ when a learning variable's PLM
uses a forward-looking variable $u\notin\mathcal{L}$ as a regressor and \code{beliefs.plm.default}
is not \code{'ar1'}: the generator gives $u$ the default PLM as a primary target, the run gives
it an auxiliary AR(1). Setting \code{beliefs.plm.by_var.<u>} makes the two agree.

\subsection{Options}

\subsubsection{\texorpdfstring{\codem{beliefs.plm.default}}{beliefs.plm.default}}\label{opt:beliefs-plm-default}
\begin{optspec}
  \item[Default] \code{'ar1'}.
  \item[Values] \code{'ar1'}, \code{'ar2'}, \code{'ar1_noconst'}, \code{'var1'}.
  \item[Math] for a learning variable $x$ without a \code{by_var} entry and outside a VAR:
        \code{'ar1'} $x_t = \alpha + \beta x_{t-1}$; \code{'ar2'}
        $x_t = \alpha + \beta_1 x_{t-1} + \beta_2 x_{t-2}$; \code{'ar1_noconst'}
        $x_t = \beta x_{t-1}$; \code{'var1'} one VAR(1) with constant over $\mathcal{L}$,
        $X_t = a + B X_{t-1}$ with $X = (x)_{x\in\mathcal{L}}$.
  \item[Assumptions] learners fit a univariate (or VAR) reduced form; they do not know the
        model's structure.
  \item[Interactions] auxiliary PLMs stay AR(1) with constant whatever the default (set them
        with \code{by_var}); see the consistency check above. \code{'var1'} together with
        \code{beliefs.plm.var.variables} is an error (\code{hetlearn:options:value});
        \code{'var1'} ignores \code{beliefs.plm.var.lags} and \code{.constant} (it is always lag
        1 with a constant) and is resolved to $\mathcal{L}$ before generation, so the generated
        file depends on the learning set. Any default other than \code{'ar1'} gives a new file
        tag.
  \item[Code] \code{+plm/build_system.m} (\code{one_spec}, \code{var_spec}),
        \code{load_model.m} (\code{resolve_learning_set}).
  \item[Example]
\begin{lstlisting}
cfg.beliefs.plm.default = 'ar2';
cfg.beliefs.plm.default = 'var1';           % VAR(1) over the learning set
\end{lstlisting}
\end{optspec}

\subsubsection{\texorpdfstring{\codem{beliefs.plm.by_var}}{beliefs.plm.by\_var}}\label{opt:beliefs-plm-by-var}
\begin{optspec}
  \item[Default] \code{struct()} (no entries).
  \item[Values] \code{by_var.<v>} is a shorthand (\code{'ar1'}, \code{'ar2'},
        \code{'ar1_noconst'}) or a struct with fields \code{regressors} (cellstr of
        \code{'u{-l}'}, $l\ge1$; \code{u(-l)}, \code{u[-l]} and \code{u{t-l}} are read too) and
        \code{constant} (logical, default true). \code{<v>} is any transition variable or
        derived name except a VAR member.
  \item[Math] equation~\eqref{eq:plm} with the listed regressors. With \code{regressors = {}}
        and a constant the PLM is mean-only, $v_t = \alpha_v + e_{v,t}$, and the forecast is
        $\alpha_{v,t-1}$. With $z=1$ and $R_0=1$ (the RE-moment value) RLS keeps $R_t=1$ and
        becomes exponential smoothing, $\alpha_t = \alpha_{t-1} + g_t(v_t - \alpha_{t-1})$:
        adaptive expectations.
  \item[Assumptions] the learners of $v$ regress on the listed lags only.
  \item[Interactions] an entry for a learning variable sets its primary PLM; an entry for any
        other variable sets its auxiliary PLM. Errors: a VAR member
        (\code{hetlearn:plm:varByVar}), an unknown variable (\code{hetlearn:plm:unknownVar}), a
        regressor that is not a transition variable or has a lag below 1
        (\code{hetlearn:plm:regressor}), duplicate regressors, neither constant nor regressors
        (\code{hetlearn:plm:empty}), unknown fields (\code{hetlearn:plm:field}). The subtree is
        free-form (merged without name checks, Section~\ref{sec:precedence}); the builder checks
        the names against the model. A regressor that is a forward-looking variable outside $\mathcal{L}$ may need
        its own entry (consistency check above).
  \item[Code] \code{+plm/build_system.m}.
  \item[Example]
\begin{lstlisting}
cfg.beliefs.plm.by_var.pinf = struct('regressors', {{'pinf{-1}', 'ygap{-1}', 'r{-1}'}}, 'constant', true);
cfg.beliefs.plm.by_var.r = 'ar2';           % auxiliary PLM of the regressor r
\end{lstlisting}
\end{optspec}

\subsubsection{\texorpdfstring{\codem{beliefs.plm.derived}}{beliefs.plm.derived}}\label{opt:beliefs-plm-derived}
\begin{optspec}
  \item[Default] \code{struct()}.
  \item[Values] \code{derived.<name> = '<expression>'}: a non-empty string of current-dated
        declared names, numbers, \code{+ - * /} and parentheses, with at least one transition
        variable; no time subscripts and no functions. \code{<name>} must be a new name.
  \item[Math] a new transition variable $d_t = \text{expression}_t$, written as
        \code{<name> = <expression>;} outside every toggle, so it exists in the RE model and in
        every learning model.
  \item[Assumptions] learners can observe the derived quantity (e.g.\ an output gap) with the
        same timing as other variables.
  \item[Interactions] a derived variable is a regressor at any lag $\ge 1$ or a VAR member;
        it gets an auxiliary AR(1) PLM like any other regressor. It is dated $t$ in the
        forecast under \code{beliefs.info = 'current'}. It changes the model text, hence the
        file tag and the RE filter's cache key; its smoothed path is in the history. Errors:
        bad syntax or unknown names (\code{hetlearn:extend_model:derived}), a name collision
        (\code{hetlearn:extend_model:nameCollision}), a non-string value
        (\code{hetlearn:options:type}).
  \item[Code] \code{extend_model.m} (\code{derived_defs}), \code{+plm/build_system.m}.
  \item[Example]
\begin{lstlisting}
cfg.beliefs.plm.derived.ygap = 'y - yf';
cfg.beliefs.plm.by_var.pinf = struct('regressors', {{'pinf{-1}', 'ygap{-1}'}}, 'constant', true);
\end{lstlisting}
\end{optspec}

\subsubsection{\texorpdfstring{\codem{beliefs.plm.var}}{beliefs.plm.var}}\label{opt:beliefs-plm-var}
\begin{optspec}
  \item[Default] \code{struct('variables', {{}}, 'lags', 1, 'constant', true)}: no VAR.
  \item[Values] \code{variables}: a cellstr of transition or derived variables, or
        \code{'learning'} ($=\mathcal{L}$); \code{lags} $p$: a positive integer; \code{constant}:
        true or false.
  \item[Math] for every member $v\in V$,
        \begin{equation}\label{eq:var}
          v_t = \alpha_v + \sum_{l=1}^{p}\sum_{u\in V} \beta_{v,u,l}\, u_{t-l} + e_{v,t},
        \end{equation}
        parameters \code{alpha_v} and \code{beta_v__u_l<l>}. Learning members use it for their
        forecast; the forecast iterates the VAR two steps under \code{'lagged'}
        information~\eqref{eq:recursion}. Members outside $\mathcal{L}$ are auxiliary
        equations of the same VAR.
  \item[Assumptions] learners use a joint reduced form. Each equation has its own RLS and
        moment matrix; with common regressors, a common gain and equal initial moment
        matrices, this equals system RLS.
  \item[Interactions] \code{'learning'} is resolved to $\mathcal{L}$ by
        \code{hetlearn.load_model} before generation (a direct
        \code{hetlearn.extend_model} call would use every forward-looking variable). A
        \code{by_var} entry for a member is an error; a member listed twice or not in the model
        is an error; \code{default = 'var1'} with \code{variables} set is an error. The VAR is
        one block of the projection facility. RE-moment initial beliefs use every
        cross-moment. Under KOZ \code{'per_variable'} each equation has its own gain process.
  \item[Code] \code{+plm/build_system.m} (\code{var_spec}, \code{var_eq}),
        \code{load_model.m}.
  \item[Example]
\begin{lstlisting}
cfg.beliefs.plm.derived.ygap = 'y - yf';
cfg.beliefs.plm.var = struct('variables', {{'pinf', 'ygap', 'r'}}, 'lags', 1, 'constant', true);
\end{lstlisting}
\end{optspec}

% =====================================================================
\section{Belief updating}\label{sec:updating}
% =====================================================================

\subsection{Recursive least squares}\label{sec:rls}

After $x_t$ is realized, each PLM equation $v$ moves its beliefs by RLS \citep[eq.~13]{ozden2025},
\citep{evans2001}:
\begin{align}
  R_{v,t} &= R_{v,t-1} + g_{v,t}\bigl(z_{v,t-1} z_{v,t-1}\transp - R_{v,t-1}\bigr), \label{eq:rlsR}\\
  \phi_{v,t} &= \phi_{v,t-1} + g_{v,t}\, R_{v,t}^{-1} z_{v,t-1}\bigl(v_t - \phi_{v,t-1}\transp z_{v,t-1}\bigr). \label{eq:rlsphi}
\end{align}
$R$ is updated first and $\phi$ uses the new $R_t$ (a linear solve, no inverse;
\code{+rls/update.m}). Equation by equation, each with its own $R$. $\phi_t$ is in force in $t+1$.
The residual $e_{v,t} = v_t - \phi_{v,t-1}\transp z_{v,t-1}$ is stored (\code{S.fe_hist}) and is the
error that drives a KOZ gain under \code{gain.koz.error = 'rls'}.

\paragraph{Order inside the update} (\code{+engine/update_learning.m}). The update runs step 5
of Section~\ref{sec:timing} in the order listed there. The gain of each equation, $g_{v,t}$, comes
from the gain state before this quarter's update and is computed before the RLS candidates; the
gain state itself is updated last, whatever the projection facility did.

\paragraph{Missing data.} An equation whose target or any regressor is not finite at $t$ keeps
its beliefs and logs an \code{rls_missing} event. This happens only where the databank the
stage reads holds a NaN; the smoothed and filtered RE histories, the inversion and the Kalman
history fill quarters without data with model values.

\paragraph{The first quarter.} At $t$ = \code{beliefs.init.date} there are no forecast errors
(the forecasts of $t-1$ were not made), so fitness and shares keep their initial values, but the
RLS update uses the data of that quarter.

\subsubsection{\texorpdfstring{\codem{beliefs.update}}{beliefs.update}}\label{opt:beliefs-update}
\begin{optspec}
  \item[Default] \code{'rls'}.
  \item[Values] \code{'rls'}, \code{'fixed'}. (The removed values \code{'sac'} and \code{'kf'}
        raise \code{hetlearn:options:removed}.)
  \item[Math] \code{'rls'}: \eqref{eq:rlsR}--\eqref{eq:rlsphi}. \code{'fixed'}:
        $\phi_t = \phi_0$, $R_t = R_0$ for all $t$.
  \item[Assumptions] \code{'fixed'}: learners keep their initial rule; forecasts still move with
        the data and shares still switch.
  \item[Interactions] with \code{'fixed'} no gain is computed, the gain state stays unchanged, and the
        PLM part of the facility has nothing to check; the ALM check can still keep the
        previous model shares (the belief revert is skipped because nothing changed).
        \code{'fixed'} with \code{beliefs.init.method = 'user'} gives a run with given
        forecasting rules.
  \item[Code] \code{+engine/update_learning.m}.
  \item[Example]
\begin{lstlisting}
cfg.beliefs.update = 'fixed';
cfg.beliefs.init.values.pinf = struct('alpha', 0.5, 'beta', 0.9);
\end{lstlisting}
\end{optspec}

\subsection{Gain rules}\label{sec:gain}

The gain $g_{v,t}$ of each equation comes from the function \code{hetlearn.gain.<rule>} named by
\code{gain.rule}. Every rule answers four calls (\code{+gain/constant.m} documents the
interface): \code{'fields'} (extra series it records), \code{'init'} (initial state),
\code{'gain'} (the gains of this quarter's RLS updates, called before them) and \code{'update'}
(the state at the end of the quarter, called after the RLS and the projection facility). The
state lives in the learning state (\code{S.gain}), so every stage that carries the learning state
carries the gain: the history, the inversion and its snapshots, the forecast, density draws,
evaluation origins and the Ramsey loss continuation. With \code{'constant'} the gain is
$g_{v,t} = \gamma_0$ (or the \code{gain.by_var} value of the equation).

\subsubsection{\texorpdfstring{\codem{gain.rule}}{gain.rule}}\label{opt:gain-rule}
\begin{optspec}
  \item[Default] \code{'constant'}. (Project baseline: \code{'koz'}.)
  \item[Values] \code{'constant'}, \code{'koz'}; the name must be a function in
        \code{+gain/} (the removed rule \code{'decreasing'} raises
        \code{hetlearn:options:removed}).
  \item[Math] \code{'constant'}: $g_{v,t} = \gamma_0$ (or the \code{by_var} value).
        \code{'koz'}: $g_{v,t} = \gamma_{p(v),t-1}$, equations~\eqref{eq:koz}.
  \item[Assumptions] see the entries below.
  \item[Interactions] \code{gain.value} is the constant gain and the KOZ start value;
        \code{gain.koz.*} matter only for \code{'koz'}; with \code{beliefs.update = 'fixed'} no
        rule is called.
  \item[Code] \code{+gain/constant.m}, \code{+gain/koz.m}, \code{+engine/update_learning.m}.
  \item[Example]
\begin{lstlisting}
cfg.gain.rule = 'koz';
cfg.gain.value = 0.05;                      % start gain
\end{lstlisting}
\end{optspec}

\subsubsection{\texorpdfstring{\codem{gain.value}}{gain.value}}\label{opt:gain-value}
\begin{optspec}
  \item[Default] \code{0.0585}, the posterior mean of the constant gain in
        \citet[Table~1]{ozden2025}. (Project baseline: 0.05 as the KOZ start.)
  \item[Values] a number in $(0, 1]$.
  \item[Math] $\gamma_0$ in \eqref{eq:rlsR}--\eqref{eq:rlsphi} under \code{'constant'}; the
        initial KOZ gain when \code{gain.koz.gamma0} is empty.
  \item[Assumptions] a constant gain weights past observations geometrically with factor
        $1-\gamma_0$ (perpetual learning, tracking a drifting reduced form).
  \item[Interactions] under \code{'koz'} it must lie inside \code{gain.koz.bounds}.
  \item[Code] \code{+gain/constant.m}, \code{+gain/koz.m}.
  \item[Example]
\begin{lstlisting}
cfg.gain.rule = 'constant';
cfg.gain.value = 0.02;
\end{lstlisting}
\end{optspec}

\subsubsection{\texorpdfstring{\codem{gain.by_var}}{gain.by\_var}}\label{opt:gain-by-var}
\begin{optspec}
  \item[Default] \code{struct()}.
  \item[Values] \code{by_var.<target> = g} with $g\in(0,1]$; \code{<target>} is a PLM target
        (a learning variable or an auxiliary one).
  \item[Math] $g_{v,t} = g$ for that equation under \code{'constant'}; the start gain of its
        process under KOZ \code{'per_variable'}.
  \item[Assumptions] learners of different variables may discount the past at different rates.
  \item[Interactions] ignored with a warning (\code{hetlearn:gain:byVarIgnored}) under KOZ
        \code{'common'}. Free-form subtree: a name that is not a PLM target is ignored silently
        by the constant rule.
  \item[Code] \code{+gain/constant.m}, \code{+gain/koz.m}.
  \item[Example]
\begin{lstlisting}
cfg.gain.by_var.pinf = 0.02;
\end{lstlisting}
\end{optspec}

\subsection{The endogenous gain of Kostyshyna, \"Ozden and Zhang}\label{sec:koz}

\citet[eq.~9]{koz2024} let the gain follow the adaptive step-size algorithm of
\citet{kostyshyna2012} \citep[see also][]{kushner2003}. For each gain process $p$, after $x_t$ is
realized (\code{+gain/koz.m}):
\begin{align}\label{eq:koz}
  \gamma_{p,t} &= \Pi_{[\underline{\gamma},\,\overline{\gamma}]}\Bigl(\gamma_{p,t-1} + \mu\, R^g_{p,t-1}\, e_{p,t}\Bigr),\\
  R^g_{p,t} &= (1-\gamma_{p,t-1})\, R^g_{p,t-1} + e_{p,t}, \notag
\end{align}
and the RLS update at the end of $t$ uses $g_{v,t} = \gamma_{p(v),t-1}$, the gain of the process of
equation $v$ before this quarter's update. $\Pi$ projects onto \code{gain.koz.bounds}; $R^g$ is
not projected. $R^g$ tracks the derivative of the beliefs with respect to the gain: it adds up
past errors with weight $1-\gamma$, so $R^g_{t-1}e_t > 0$ and the gain rises while the errors keep
their sign, and the gain falls when they alternate.

The driving error is signed and scaled,
\begin{equation}\label{eq:kozerr}
  e_{p,t} = \frac{\sigma}{|D_p|}\sum_{d\in D_p} e_{d,t},\qquad
  e_{d,t} =
  \begin{cases}
    v_{d,t} - \phi_{d,t-1}\transp z_{d,t-1} & \text{\codem{error = 'rls'}}\\
    x_{d,t} - \fBR_{t-1}(x_d) & \text{\codem{error = 'forecast'}}
  \end{cases}
\end{equation}
with $\sigma$ = \code{gain.koz.error_scale} and $D_p$ the drivers of process $p$. A missing error
(the first quarter under \code{'forecast'}, missing data, or a missing error of any driver of the
list) leaves $\gamma_p$ and $R^g_p$ unchanged. A bound that binds sets $\gamma$ to the bound and logs
a \code{gain_bound} event; a non-finite update keeps the state and logs \code{gain_bound} too.

\paragraph{Timing.} The \code{'gain'} call returns $\gamma_{t-1}$ before the RLS candidates; the
\code{'update'} call applies~\eqref{eq:koz} after the RLS and the projection facility. A reverted
or rejected belief update does not revert the gain: the error was observed either way. With
$\mu = 0$, $\gamma_{p,t} = \gamma_{p,0}$ and the run equals the constant rule bit for bit
(\code{test_gain} G2, G7).

\paragraph{Outputs.} \code{res.learning.<target>.gain} is the gain used in $t$;
\code{gain_R} is $R^g_t$ of the target's process; \code{res.state.gain} is the final state.
Constant-gain runs record no \code{gain_R}.

\subsubsection{\texorpdfstring{\codem{gain.koz.mu}}{gain.koz.mu}}\label{opt:gain-koz-mu}
\begin{optspec}
  \item[Default] \code{0.001}, the step size of \citet{kostyshyna2012} used by \citet{koz2024}.
  \item[Values] a finite number $\ge 0$.
  \item[Math] $\mu$ in~\eqref{eq:koz}, the gain on the gain.
  \item[Assumptions] $\mu$ is scale dependent: $e$ enters both equations, so the effective step
        scales with $\sigma^2$.
  \item[Interactions] \code{error_scale}; $\mu=0$ reproduces \code{'constant'}.
  \item[Code] \code{+gain/koz.m}.
  \item[Example]
\begin{lstlisting}
cfg.gain.rule = 'koz';
cfg.gain.koz.mu = 0.0005;
\end{lstlisting}
\end{optspec}

\subsubsection{\texorpdfstring{\codem{gain.koz.gamma0}}{gain.koz.gamma0}}\label{opt:gain-koz-gamma0}
\begin{optspec}
  \item[Default] \code{[]}: \code{gain.value}, and under \code{'per_variable'}
        \code{gain.by_var.<target>} where set.
  \item[Values] \code{[]} or a number in $(0,1]$ inside the bounds.
  \item[Math] $\gamma_{p,0}$, the gain of the first RLS update.
  \item[Assumptions] KOZ start the process at their estimated constant gain.
  \item[Interactions] the initial gain (this value, or \code{gain.value} and the
        \code{by_var} values) must lie inside \code{gain.koz.bounds}
        (\code{hetlearn:options:value}).
  \item[Code] \code{+gain/koz.m}, \code{load_options.m} (\code{validate_gain}).
  \item[Example]
\begin{lstlisting}
cfg.gain.rule = 'koz';
cfg.gain.koz.gamma0 = 0.03;
\end{lstlisting}
\end{optspec}

\subsubsection{\texorpdfstring{\codem{gain.koz.Rg0}}{gain.koz.Rg0}}\label{opt:gain-koz-rg0}
\begin{optspec}
  \item[Default] \code{0}.
  \item[Values] a finite number.
  \item[Math] $R^g_{p,0}$ in~\eqref{eq:koz}. With $R^g_0=0$ the first update leaves the gain
        unchanged and sets $R^g_1 = e_1$.
  \item[Assumptions] KOZ are silent; 0 states that the initial estimate does not depend on
        the gain. KOZ's Table~1 ``$R_0 = 0.001$'' is the RLS moment matrix, not $R^g_0$; it is set
        through \code{beliefs.init.values.<x>.R} (Section~\ref{opt:beliefs-init-values}).
  \item[Interactions] none beyond \code{'koz'}.
  \item[Code] \code{+gain/koz.m}.
  \item[Example]
\begin{lstlisting}
cfg.gain.rule = 'koz';
cfg.gain.koz.Rg0 = 0;
\end{lstlisting}
\end{optspec}

\subsubsection{\texorpdfstring{\codem{gain.koz.bounds}}{gain.koz.bounds}}\label{opt:gain-koz-bounds}
\begin{optspec}
  \item[Default] \code{[0.001 0.25]}.
  \item[Values] \code{[lower, upper]} with $0 < \text{lower} \le \text{upper} \le 1$.
  \item[Math] $[\underline{\gamma},\overline{\gamma}]$ of the projection $\Pi$ in~\eqref{eq:koz}.
  \item[Assumptions] \citet{kostyshyna2012} projects the gain onto an interval; KOZ give none,
        so the interval is a toolbox choice.
  \item[Interactions] the initial gain must lie inside; each binding quarter logs a
        \code{gain_bound} event (counted in density outputs).
  \item[Code] \code{+gain/koz.m}.
  \item[Example]
\begin{lstlisting}
cfg.gain.rule = 'koz';
cfg.gain.koz.bounds = [0.005 0.15];
\end{lstlisting}
\end{optspec}

\subsubsection{\texorpdfstring{\codem{gain.koz.error}}{gain.koz.error}}\label{opt:gain-koz-error}
\begin{optspec}
  \item[Default] \code{'rls'}.
  \item[Values] \code{'rls'}, \code{'forecast'}.
  \item[Math] the case in~\eqref{eq:kozerr}: the one-step PLM residual (the error of the RLS
        update itself, KOZ eq.~8) or the learners' expectation error $x_t - \fBR_{t-1}(x)$.
  \item[Assumptions] KOZ do not say which error; the residual is the default.
  \item[Interactions] \code{'forecast'} is a two-step error under \code{beliefs.info =
        'lagged'}; it is missing in the first learning quarter. Auxiliary equations have no
        expectation and use their residual. With \code{'common'} the master names PLM targets
        (\code{'rls'}) or learning variables (\code{'forecast'}).
  \item[Code] \code{+gain/koz.m}, \code{+engine/update_learning.m}.
  \item[Example]
\begin{lstlisting}
cfg.gain.rule = 'koz';
cfg.gain.koz.error = 'forecast';
\end{lstlisting}
\end{optspec}

\subsubsection{\texorpdfstring{\codem{gain.koz.error_scale}}{gain.koz.error\_scale}}\label{opt:gain-koz-error-scale}
\begin{optspec}
  \item[Default] \code{1} (model units, percent per quarter for SW07).
  \item[Values] a finite number $>0$.
  \item[Math] $\sigma$ in~\eqref{eq:kozerr}.
  \item[Assumptions] KOZ's data are annualized; a factor of 4 converts quarterly to annualized
        percent and acts like a 16 times larger $\mu$.
  \item[Interactions] \code{gain.koz.mu}.
  \item[Code] \code{+gain/koz.m}.
  \item[Example]
\begin{lstlisting}
cfg.gain.rule = 'koz';
cfg.gain.koz.error_scale = 4;               % annualized errors
\end{lstlisting}
\end{optspec}

\subsubsection{\texorpdfstring{\codem{gain.koz.structure}}{gain.koz.structure}}\label{opt:gain-koz-structure}
\begin{optspec}
  \item[Default] \code{'common'}.
  \item[Values] \code{'common'}, \code{'per_variable'}.
  \item[Math] \code{'common'}: one process for every PLM equation (auxiliary ones included),
        driven by the errors of \code{gain.koz.master} \citep[fn.~7]{koz2024}.
        \code{'per_variable'}: one process per PLM equation, driven by its own error.
  \item[Assumptions] \code{'common'}: one learning speed for all beliefs.
  \item[Interactions] \code{'common'} ignores \code{gain.by_var} (warning);
        \code{'per_variable'} ignores \code{gain.koz.master} (warning
        \code{hetlearn:gain:masterIgnored}) and takes start values from \code{gain.by_var}.
  \item[Code] \code{+gain/koz.m} (\code{init}).
  \item[Example]
\begin{lstlisting}
cfg.gain.rule = 'koz';
cfg.gain.koz.structure = 'per_variable';
\end{lstlisting}
\end{optspec}

\subsubsection{\texorpdfstring{\codem{gain.koz.master}}{gain.koz.master}}\label{opt:gain-koz-master}
\begin{optspec}
  \item[Default] \code{''}: the first learning variable (the first PLM target if there is no
        learning variable).
  \item[Values] a name or a cellstr of valid names.
  \item[Math] $D_p$ in~\eqref{eq:kozerr}: a list gives the mean of the scaled errors.
  \item[Assumptions] the common learning speed reacts to the errors of the named variables.
  \item[Interactions] \code{'common'} only. Under \code{'rls'} the names must be PLM targets,
        under \code{'forecast'} learning variables (\code{hetlearn:gain:master}, checked when
        the gain state is built). A missing error of any driver freezes the process that
        quarter.
  \item[Code] \code{+gain/koz.m}.
  \item[Example]
\begin{lstlisting}
cfg.gain.rule = 'koz';
cfg.gain.koz.master = {'pinf', 'w'};
\end{lstlisting}
\end{optspec}

% =====================================================================
\section{Initial beliefs}\label{sec:init}
% =====================================================================

Initial beliefs $(\phi_0, R_0)$ are set once per run, by \code{+engine/state_init.m}, for every
PLM equation (auxiliary ones included). They are the beliefs in force at
$t$ = \code{beliefs.init.date}: they make the learners' forecast in that quarter, and the first
RLS update uses that quarter's data.

\subsection{Options}

\subsubsection{\texorpdfstring{\codem{beliefs.init.method}}{beliefs.init.method}}\label{opt:beliefs-init-method}
\begin{optspec}
  \item[Default] \code{'re_moments'} \citep[sec.~3.1]{ozden2025}.
  \item[Values] \code{'re_moments'}, \code{'training_ols'}, \code{'user'}.
  \item[Math] \code{'re_moments'} (\code{+init/re_moments.m}): for each equation,
        \begin{equation}\label{eq:remoments}
          R_{v,0} = E\bigl[z_{v} z_{v}\transp\bigr], \qquad \phi_{v,0} = E\bigl[z_v z_v\transp\bigr]^{-1} E\bigl[z_v v\bigr],
        \end{equation}
        the population projection of $v_t$ on $z_{v,t-1}$ in the RE model. The moments are
        $E[a_{t-l_a} b_{t-l_b}] = \mathrm{Cov}(a_{t-l_a}, b_{t-l_b}) + \mu_a\mu_b$ with the
        autocovariances of IRIS \code{acf(m, 'order', max(maxlag, 1))} and the means $\mu$ from
        the steady state. With a constant, $E[zz\transp] = \bigl[\begin{smallmatrix}1 & \mu\transp\\ \mu & E[uu\transp]\end{smallmatrix}\bigr]$.
        For a zero-mean AR(1) PLM this gives $\alpha_0 = 0$, $\beta_0 = \rho_1(x)$ and
        $R_0 = \mathrm{diag}(1, \mathrm{var}\,x)$. A near-singular $E[zz\transp]$ (reciprocal
        condition below $10^{-14}$) uses the pseudo-inverse with a warning.
        \code{'training_ols'} (\code{+init/training_ols.m}): OLS over the training range,
        $\phi_{v,0} = (Z\transp Z)^{-1}Z\transp y$, $R_{v,0} = Z\transp Z / T$, with rows $z_{v,t-1}\transp$ and
        targets $v_t$ for $t$ in the range; rows with a missing value are dropped.
        \code{'user'} (\code{+init/user.m}): the values given, the RE moments for anything not
        given.
  \item[Assumptions] \code{'re_moments'}: learners start at the least-squares projection that the
        moments of the RE economy imply. The model is the loaded model of
        the run at $n_1=1$: the RE model of the scenario, the RE(t-1) model under
        \code{re_info = 'lagged'}, and the rule model under Ramsey policy.
        \code{'training_ols'}: learners start from a regression on the history they watch.
  \item[Interactions] \code{'training_ols'} needs \code{beliefs.init.training_range}
        (\code{hetlearn:options:missing}) and reads the history databank of
        \code{history.re_states} (smoothed or filtered); fewer usable rows than coefficients is
        \code{hetlearn:init:shortSample}. \code{'user'} reads \code{beliefs.init.values}. With
        \code{history.learn = false} the initial beliefs apply at the origin.
  \item[Code] \code{+engine/state_init.m} and, in \code{+init/}, \code{re_moments.m},
        \code{training_ols.m}, \code{user.m}.
  \item[Example]
\begin{lstlisting}
cfg.beliefs.init.method = 'training_ols';
cfg.beliefs.init.training_range = [qq(1966,1), qq(1979,4)];
\end{lstlisting}
\end{optspec}

\subsubsection{\texorpdfstring{\codem{beliefs.init.date}}{beliefs.init.date}}\label{opt:beliefs-init-date}
\begin{optspec}
  \item[Default] \code{[]}: the first filter date plus 4 quarters (SW07: 1966Q1).
  \item[Values] an IRIS date with room for the PLM lags after the filter start and, with
        \code{history.learn = true}, not after the origin
        (\code{hetlearn:engine:initDate}).
  \item[Math] the quarter $t_0$ at which $\phi_{t_0-1} = \phi_0$ and the learning record starts.
  \item[Assumptions] learning starts here; shares start at \code{shares.init_n2}.
  \item[Interactions] the default start of \code{history.inversion.start}; the start of the
        Kalman likelihood (\code{loglik_start}) and of a \code{'smoothed'} Kalman start; ignored
        with \code{history.learn = false} (the record then starts at the origin).
  \item[Code] \code{+engine/core.m}.
  \item[Example]
\begin{lstlisting}
cfg.beliefs.init.date = qq(1985,1);
cfg.history.learn = true;
\end{lstlisting}
\end{optspec}

\subsubsection{\texorpdfstring{\codem{beliefs.init.values}}{beliefs.init.values}}\label{opt:beliefs-init-values}
\begin{optspec}
  \item[Default] \code{struct()}.
  \item[Values] \code{values.<target>} with any of \code{.phi} (the full vector, in parameter
        order), \code{.alpha}, \code{.beta} (the slope vector, regressor order), \code{.R} (a
        $k\times k$ matrix). \code{.phi} takes precedence over \code{.alpha}/\code{.beta}.
  \item[Math] $\phi_{v,0}$, $R_{v,0}$ for the targets given; \eqref{eq:remoments} for the rest.
  \item[Assumptions] the user's values are the learners' prior beliefs.
  \item[Interactions] \code{beliefs.init.method = 'user'} only. \code{.alpha} for a PLM without a
        constant and wrong sizes are errors (\code{hetlearn:init:noConstant},
        \code{hetlearn:init:size}); entries for variables without a PLM equation are ignored with
        a warning. Setting only \code{.R} keeps the RE-moment $\phi$ and replaces $R$, e.g.\ KOZ's
        $R_0 = 0.001\,I$.
  \item[Code] \code{+init/user.m}.
  \item[Example]
\begin{lstlisting}
cfg.beliefs.init.method = 'user';
cfg.beliefs.init.values.pinf = struct('R', 0.001*eye(2));
\end{lstlisting}
\end{optspec}

\subsubsection{\texorpdfstring{\codem{beliefs.init.training_range}}{beliefs.init.training\_range}}\label{opt:beliefs-init-training-range}
\begin{optspec}
  \item[Default] \code{[]}.
  \item[Values] \code{[start, end]} or a date vector of target dates.
  \item[Math] the set of $t$ in the OLS of \code{'training_ols'}; regressors are read at
        $t - l_j$.
  \item[Assumptions] the training sample is in the history the learners watch.
  \item[Interactions] required by \code{'training_ols'}, unused otherwise. It can overlap the
        quarters of the learning record, in which case the learners start from beliefs fitted on data they
        later learn from again.
  \item[Code] \code{+init/training_ols.m}.
  \item[Example]
\begin{lstlisting}
cfg.beliefs.init.method = 'training_ols';
cfg.beliefs.init.training_range = [qq(1966,1), qq(1984,4)];
\end{lstlisting}
\end{optspec}

% =====================================================================
\section{Switching between RE and learning}\label{sec:switching}
% =====================================================================

\subsection{Forecast errors, fitness and shares}\label{sec:dc}

At the end of $t$ the toolbox scores the two forecasts of $x_t$ made in $t-1$
(\code{+engine/update_learning.m}):
\begin{equation}\label{eq:fe}
  \FE^{i}_t(x) = \bigl(x_t - f^{i}_{t-1}(x)\bigr)^2, \qquad i\in\{\RE,\BR\}.
\end{equation}
The squared errors enter unscaled, in model units. Each share group $g$ sums them over its
variables $G_g$, $\FE^i_{g,t} = \sum_{v\in G_g}\FE^i_t(v)$. Fitness has memory $\omega$ and the RE
share follows a discrete choice with intensity $\chi$ \citep[eqs.~14--15]{ozden2025},
\citep{brock1997}:
\begin{align}
  \zeta^{i}_{g,t} &= -(1-\omega)\,\FE^{i}_{g,t} + \omega\,\zeta^{i}_{g,t-1}, \label{eq:fitness}\\
  n_{1,g,t} &= \frac{\exp(\chi\zeta^{\RE}_{g,t})}{\exp(\chi\zeta^{\RE}_{g,t}) + \exp(\chi\zeta^{\BR}_{g,t})}
  = \frac{1}{1 + \exp\bigl(-\chi(\zeta^{\RE}_{g,t} - \zeta^{\BR}_{g,t})\bigr)}, \qquad n_{2,g,t} = 1 - n_{1,g,t}. \label{eq:logit}
\end{align}
The logistic form is evaluated without overflow for any value of $\chi(\zeta^{\RE}-\zeta^{\BR})$
(\code{+switching/discrete_choice.m}). With $\omega = 0$ the memory term is dropped, so an infinite
initial fitness does not produce $0\cdot\infty$; $\chi = 0$ gives $n_1 = 1/2$; two infinite
fitness values of the same sign give $1/2$. A non-finite share keeps the previous fitness and
shares and logs a \code{shares_nan} event.

\paragraph{Timing.} The shares computed at the end of $t$ enter the model in $t+1$
\citep[eq.~16]{ozden2025}: under \code{'switching'}, $n^m_{1,t+1} = n_{1,t}$
(Section~\ref{opt:expectations-mode}). The forecasts scored at $t$ were made in $t-1$, so the
model in force in $t$ did not see the errors of $t$.

\paragraph{Initial fitness.} \code{+switching/init_fitness.m} sets
\begin{equation}\label{eq:initfit}
  \zeta^{\RE}_{0} = 0, \qquad \zeta^{\BR}_{0} = \frac{\ln \bar n_2 - \ln(1-\bar n_2)}{\chi},
  \qquad n_{1,0} = 1 - \bar n_2,
\end{equation}
with $\bar n_2$ = \code{shares.init_n2}, so the logit~\eqref{eq:logit} returns $\bar n_2$ at the
start. $\bar n_2 = 0$ gives $\zeta^{\BR}_0 = -\infty$ and $\bar n_2 = 1$ gives $+\infty$. With
$\chi = 0$ both fitness values are 0 and a warning says the shares will be $1/2$ after the first
quarter.

\paragraph{Diagnostics outside switching.} The fitness and the shares are updated in every mode;
outside \code{'switching'} they are recorded (\code{n1}, \code{n2}, \code{zetaRE},
\code{zetaBR}) but do not enter the model, whose share is \code{n1_model}.

\paragraph{Zero future shocks.} With zero shocks over the horizon both squared errors shrink, and
the fitness of both predictors decays at rate $\omega$; the shares drift toward $1/2$ by
construction (there is no cost term for RE).

\paragraph{Units of $\chi$.} $\chi$ multiplies squared errors. \citet{ozden2025} estimates $\chi = 0.51$ on data
in percent per quarter, the units of SW07, so it applies as estimated; for annualized data the same
switching needs $\chi/16$. A published intensity $\beta$ of the unnormalized form
$U_t = \omega U_{t-1} - \FE_t$ corresponds to $\chi = \beta/(1-\omega)$ here.

\subsection{Options}

\subsubsection{\texorpdfstring{\codem{switching.rule}}{switching.rule}}\label{opt:switching-rule}
\begin{optspec}
  \item[Default] \code{'discrete_choice'}.
  \item[Values] the name of a function \code{hetlearn.switching.<rule>} with the interface of
        \code{discrete_choice} (\code{[n1, zRE, zBR] = rule(FE_RE, FE_BR, zRE_prev,
        zBR_prev, chi, omega)}, vectorized over share groups).
  \item[Math] equations~\eqref{eq:fitness}--\eqref{eq:logit}.
  \item[Assumptions] agents choose the predictor with the higher past fitness with probability
        given by a logit.
  \item[Interactions] \code{switching.intensity}, \code{switching.memory}; the initial fitness
        always comes from \code{init_fitness}.
  \item[Code] \code{+switching/discrete_choice.m}, \code{load_options.m} (the function must
        exist).
  \item[Example]
\begin{lstlisting}
cfg.switching.rule = 'discrete_choice';
\end{lstlisting}
\end{optspec}

\subsubsection{\texorpdfstring{\codem{switching.intensity}}{switching.intensity}}\label{opt:switching-intensity}
\begin{optspec}
  \item[Default] \code{0.51}, the estimate of \citet{ozden2025} on quarterly data. (Project
        baseline: 2.)
  \item[Values] a number $\ge 0$.
  \item[Math] $\chi$ in~\eqref{eq:logit} and~\eqref{eq:initfit}.
  \item[Assumptions] $\chi = 0$: shares stay at $1/2$ after the first quarter; $\chi\to\infty$:
        everyone takes the better predictor.
  \item[Interactions] the units of the squared errors (above); \code{shares.init_n2} through
        \eqref{eq:initfit}.
  \item[Code] \code{+switching/discrete_choice.m}, \code{+switching/init_fitness.m}.
  \item[Example]
\begin{lstlisting}
cfg.switching.intensity = 1;
cfg.switching.memory = 0.6;
\end{lstlisting}
\end{optspec}

\subsubsection{\texorpdfstring{\codem{switching.memory}}{switching.memory}}\label{opt:switching-memory}
\begin{optspec}
  \item[Default] \code{0.6} \citep[sec.~4; the Table~1 posterior mean is 0.62]{ozden2025}.
        (Project baseline: 0.8.)
  \item[Values] a number in $[0, 1)$.
  \item[Math] $\omega$ in~\eqref{eq:fitness}.
  \item[Assumptions] fitness is an exponentially weighted average of past squared errors with
        weight $1-\omega$ on the newest.
  \item[Interactions] $\omega = 0$ drops the memory term.
  \item[Code] \code{+switching/discrete_choice.m}.
  \item[Example]
\begin{lstlisting}
cfg.switching.memory = 0.8;
\end{lstlisting}
\end{optspec}

\subsubsection{\texorpdfstring{\codem{shares.init_n2}}{shares.init\_n2}}\label{opt:shares-init-n2}
\begin{optspec}
  \item[Default] \code{0.5}.
  \item[Values] a number in $[0, 1]$.
  \item[Math] $\bar n_2$ in~\eqref{eq:initfit}: the learners' share at \code{beliefs.init.date}
        (and at the origin with \code{history.learn = false}).
  \item[Assumptions] the initial fitness difference is chosen to return this share.
  \item[Interactions] used by \code{'switching'} through the fitness; under \code{'switching'} it
        is also the model share in the first forecast quarter when \code{history.learn = false}.
        The other modes ignore it for the model share.
  \item[Code] \code{+engine/state_init.m}, \code{+switching/init_fitness.m}.
  \item[Example]
\begin{lstlisting}
cfg.shares.init_n2 = 0.2;
cfg.expectations.mode = 'switching';
\end{lstlisting}
\end{optspec}

\subsubsection{\texorpdfstring{\codem{shares.fixed_n2}}{shares.fixed\_n2}}\label{opt:shares-fixed-n2}
\begin{optspec}
  \item[Default] \code{[]}.
  \item[Values] a number in $[0,1]$, or a struct by learning variable with values in $[0,1]$.
  \item[Math] $\bar n_2$ in~\eqref{eq:modes}: $n^m_{1,t} = 1-\bar n_2$ (or $1-\bar n_{2,x}$) in
        every quarter.
  \item[Assumptions] the population is fixed; switching is a diagnostic.
  \item[Interactions] required by \code{expectations.mode = 'fixed_shares'}
        (\code{hetlearn:options:missing}); a struct must have an entry for every learning
        variable (\code{hetlearn:engine:fixedShares}). Not together with the N-type
        \code{shares.fixed}. Fixed shares of 0 make the ALM check skip.
  \item[Code] \code{+engine/shares_for_mode.m}.
  \item[Example]
\begin{lstlisting}
cfg.expectations.mode = 'fixed_shares';
cfg.shares.fixed_n2 = struct('pinf', 0.3, 'w', 0.5);
\end{lstlisting}
\end{optspec}

\subsubsection{\texorpdfstring{\codem{shares.structure}}{shares.structure}}\label{opt:shares-structure}
\begin{optspec}
  \item[Default] \code{'per_variable'}.
  \item[Values] \code{'per_variable'}, \code{'universal'}.
  \item[Math] \code{'per_variable'}: one group per learning variable, $G_x = \{x\}$, each with
        its own fitness and share. \code{'universal'}: one group for all learning variables,
        $G = $ \code{shares.master}; every learning variable uses the group's share.
  \item[Assumptions] \code{'universal'}: an agent is RE or a learner for every variable at once,
        and judges predictors by the summed squared errors of the master variables, as
        \citet{ozden2025} does for inflation and the output gap.
  \item[Interactions] \code{'universal'} needs \code{shares.master}
        (\code{hetlearn:options:missing}). The recorded \code{FE_RE}, \code{FE_BR} stay each
        variable's own errors; the group sums drive the shares.
  \item[Code] \code{+engine/state_init.m} (share groups), \code{+engine/update_learning.m}.
  \item[Example]
\begin{lstlisting}
cfg.shares.structure = 'universal';
cfg.shares.master = {'pinf', 'y'};
\end{lstlisting}
\end{optspec}

\subsubsection{\texorpdfstring{\codem{shares.master}}{shares.master}}\label{opt:shares-master}
\begin{optspec}
  \item[Default] \code{{}}.
  \item[Values] a name or a cellstr of learning variables.
  \item[Math] $G$ of the universal group: $\FE^i_{t} = \sum_{v\in G}\FE^i_t(v)$.
  \item[Assumptions] the share responds to the errors of these variables only.
  \item[Interactions] \code{'universal'} only; every name must be a learning variable
        (\code{hetlearn:engine:master}).
  \item[Code] \code{+engine/state_init.m}.
  \item[Example]
\begin{lstlisting}
cfg.shares.structure = 'universal';
cfg.shares.master = 'pinf';
\end{lstlisting}
\end{optspec}

% =====================================================================
\section{The projection facility}\label{sec:facility}
% =====================================================================

Learning models can lose their stable solution. The projection facility is the toolbox's only
stabilizing device (the earlier devices, partial updates, shrinking, a root threshold and error
scaling, were removed; their option names raise \code{hetlearn:options:removed}). It restricts
beliefs in the spirit of \citet{marcet1989} and \citet{evans2001}.

\subsection{PLM blocks and their roots}\label{sec:blocks}

Two PLM equations are linked when one has the other's target among its regressors (any lag). A
\emph{block} is a connected component of that graph (\code{+plm/blocks.m}): a VAR is one block;
an auxiliary equation belongs to the block of the equations that use it, so two learning variables
that share a regressor are one block; separate AR(1) PLMs are separate blocks. A block is closed,
so its companion matrix is defined (\code{+plm/companion.m}): with $y_t$ the $n$ targets of the
block and $L$ its largest lag,
\begin{equation}\label{eq:companion}
  X_t = F X_{t-1},\qquad X_t = \bigl(1,\, y_t\transp,\, y_{t-1}\transp,\, \dots,\, y_{t-L+1}\transp\bigr)\transp,
  \qquad \rho = \max\bigl|\mathrm{eig}(F_{2:,2:})\bigr|,
\end{equation}
the spectral radius without the unit root of the constant ($\rho = |\beta|$ for an AR(1); NaN if a
coefficient is not finite; 0 for a mean-only PLM). The companion of the whole system is block
diagonal after a permutation, so its largest root is the largest block root.

\subsection{The three settings}\label{sec:facilitysettings}

\subsubsection{\texorpdfstring{\codem{projection.facility}}{projection.facility}}\label{opt:projection-facility}
\begin{optspec}
  \item[Default] \code{'plm_alm'}. (Project baseline: \code{'alm'}.)
  \item[Values] \code{'plm_alm'}, \code{'alm'}, \code{'none'}.
  \item[Math] at the end of $t$, after the RLS candidates $(\tilde\phi_t, \tilde R_t)$ are formed
        from $(\phi_{t-1}, R_{t-1})$ (\code{+engine/facility.m}):
        \begin{enumerate}[nosep,label=(\alph*)]
          \item \emph{PLM part} (\code{'plm_alm'} only). For each block with an updated
                equation, compute $\rho$ of the block at the candidate beliefs (the other
                equations of the block at their previous values). If $\rho$ is finite, $\rho\le 1$
                and every candidate $R$ is finite, the block takes its candidates; otherwise all
                its equations keep $\phi_{t-1}$ and $R_{t-1}$ (event \code{projection_plm}, with the
                block and its root). Every candidate comes from the previous beliefs and each block
                is judged on its own equations, so the result does not depend on the equation
                order.
          \item \emph{ALM check} (\code{'alm'} and \code{'plm_alm'}). Solve the model of $t+1$:
                beliefs $\phi_t$ and the mode's shares for $t+1$. It is stable if IRIS finds a
                unique stable solution, or, for a model without leads, if the backward system's
                largest root is at most 1. The check is skipped when every learners' model share
                is 0.
          \item \emph{Revert.} If the ALM is unstable and this quarter changed any belief, every
                belief ($\phi$ and $R$ of every PLM equation, auxiliary ones included) goes back to
                its value before this quarter's update, $(\phi_{t-1}, R_{t-1})$, and the model is
                solved again (event \code{projection_alm}). There is no search for the equation
                that causes the instability.
          \item \emph{Keep the shares.} If the model is still unstable and the new model shares
                differ from those in force in $t$, the shares in force in $t$ are kept for $t+1$
                with the reverted beliefs, i.e.\ the model of $t$, which solved (event
                \code{projection_alm_shares}; state \code{S.n1m_keep}, one quarter only). The
                fitness and the discrete-choice shares $n_{1,t}$ keep their update; the next
                quarter tries the new shares again.
          \item \emph{Stop.} If no step gives a stable model, the run stops with
                \code{hetlearn:engine:almUnstable} (date, stage, shares, beliefs, IRIS flag). This
                happens only when no earlier solved model exists: the initial beliefs and shares
                are unstable at the first check of a run, or a stage checks a model other than the
                one checked before (the Ramsey forecast model at the origin) and the previous
                beliefs and shares do not solve it either.
        \end{enumerate}
        \code{'none'} restricts nothing: no PLM part and no ALM check. A quarter that must be
        simulated without a solution stops with \code{hetlearn:engine:noSolution} (date, stage,
        learners' shares, the largest root of every PLM block, beliefs, IRIS flag), except a
        linear model without leads, which runs on its backward system with \code{explosive}
        events.
  \item[Assumptions] \code{'alm'} allows explosive PLMs as long as the economy is stable;
        \code{'plm_alm'} also rules out explosive perceived dynamics block by block. Beliefs and
        shares last in force are always admissible because they solved.
  \item[Interactions] the facility never touches the gain state: a KOZ gain updates from the
        observed error whatever the facility did. A revert detected at the start of a quarter
        (a model that the end-of-quarter check did not solve, e.g.\ the first forecast quarter
        with \code{history.learn = false}) rewrites the previous quarter's record, so
        \code{res.learning} holds the beliefs actually in force; kept shares become that
        quarter's \code{n1_model}. With \code{beliefs.update = 'fixed'} the revert step finds
        nothing to revert and goes to the share step. With fixed shares the share step finds
        nothing to keep. Under \code{'none'} a comparison target that stops does not fail the run
        that compares with it (warning \code{hetlearn:save_results:compare}).
  \item[Code] \code{+engine/facility.m}, \code{+plm/blocks.m}, \code{+plm/companion.m},
        \code{+engine/update_learning.m}, \code{+engine/shares_for_mode.m},
        \code{+engine/solve_model.m}.
  \item[Example]
\begin{lstlisting}
cfg.projection.facility = 'alm';            % restrict the ALM only
cfg.expectations.variables = 'all';
\end{lstlisting}
\end{optspec}

\subsection{Kept shares}\label{sec:keptshares}

The kept shares (step d) apply for one quarter: \code{+engine/update_learning.m} clears
\code{S.n1m_keep} before the next check and records the shares in force in $t$ as
\code{S.n1m_last}. In the record, \code{n1}, \code{n2} are the discrete-choice shares and
\code{n1_model} the share in force, kept ones included. Over the RE-smoothed history
\code{n1_model} stays NaN (no model is simulated there), but kept shares are carried in the state,
so the next check and the first forecast quarter use them.

\subsection{Where the check runs}\label{sec:wherecheck}

Beliefs formed at the end of $t$ are in force in $t+1$, so the ALM check at the end of $t$ solves
the model of $t+1$, and the stage reuses that solution in $t+1$ without another solve. The same
rules apply in every stage that updates beliefs:

\begin{longtable}{>{\raggedright\arraybackslash}p{4.2cm}>{\raggedright\arraybackslash}p{10.8cm}}
\toprule
Stage & Model checked at the end of $t$ \\
\midrule
\endhead
RE-smoothed history & the history model before the origin (the rule model of a Ramsey scenario); the forecast model at the origin, whose solution the forecast's first quarter reuses (\code{fopt.warm}). The model is solved only for the check \\
Inversion history & the model being inverted; the forecast model at the origin \\
Kalman history & the ALM being filtered; the forecast model at the origin \\
Forecast & the model of the next quarter \\
Density draws, evaluation runs, Ramsey loss continuation & as the forecast \\
\bottomrule
\end{longtable}

\code{res.meta.history_alm_check_s}, \code{.history_alm_n_check} and \code{.history_alm_n_solve}
measure the history-stage checks; \code{res.meta.forecast_n_solve} counts the forecast's solves
and \code{res.meta.projection_facility} records the setting.

\subsection{Events}\label{sec:events}

\begin{longtable}{>{\raggedright\arraybackslash}p{4.2cm}>{\raggedright\arraybackslash}p{10.8cm}}
\toprule
Event & Meaning \\
\midrule
\endhead
\code{projection_plm} & a PLM block's candidate had a root above 1 (or was not finite); the block kept its beliefs \\
\code{projection_alm} & the ALM with the new beliefs had no stable solution; every belief was reverted \\
\code{projection_alm_shares} & still unstable with the old beliefs; the previous quarter's model shares stayed in force \\
\code{explosive} & a model without leads was simulated by the backward system (no stable IRIS solution) \\
\code{rls_missing} & missing data in a PLM equation; it kept its beliefs \\
\code{shares_nan} & non-finite shares; the previous fitness and shares were kept \\
\code{gain_bound} & the KOZ gain hit a bound of \code{gain.koz.bounds} (or a non-finite gain update was skipped) \\
\bottomrule
\end{longtable}

Events are returned in \code{res.events} and \code{events.csv} with columns \code{date},
\code{type}, \code{variable}, \code{detail} and \code{stage} (\code{history} up to the origin,
\code{inversion} from \code{history.inversion.start} to the origin in an inversion run,
\code{forecast} after the origin; \code{+engine/assemble.m}).

% =====================================================================
\section{The history stage}\label{sec:history}
% =====================================================================

The history stage runs from \code{beliefs.init.date} $t_0$ to the origin $T_0$. It produces the
learning state at the end of $T_0$ (beliefs, moments, fitness, shares, gain state, the forecasts of
$x_{T_0+1}$ made at $T_0$) and the databank the forecast starts from. Four designs are available
through \code{history.mode} and \code{history.re_states}:

\begin{center}
\begin{tabular}{p{3.3cm}p{5.3cm}p{5.8cm}}
\toprule
Setting & What the agents watch & What the forecast starts from \\
\midrule
\code{'re_smoother'}, \code{'smoothed'} & the RE-smoothed history (two-sided) & the RE-smoothed state at $T_0$ \\
\code{'re_smoother'}, \code{'filtered'} & the RE-filtered history (one-sided) & the RE-smoothed state at $T_0$ \\
\code{'inversion'} & the history backed out with the learning model & the backed-out state at $T_0$ \\
\code{'kalman'} & the real-time filtered state of the ALM in force & the filtered state at $T_0$ \\
\bottomrule
\end{tabular}
\end{center}

Under \code{'re_smoother'} the path of the economy does not depend on the learning state: agents
learn while watching the RE economy. Under \code{'inversion'} and \code{'kalman'} the learning
model generates the history.

\subsection{The RE filter and smoother}\label{sec:refilter}

\code{hetlearn.re_filter(spec, opts)} (\code{re_filter.m}) runs the IRIS Kalman filter and smoother
of the RE model over \code{spec.filter_range} and caches the result.

\paragraph{The filter model.} Under \code{expectations.re_info = 'current'} it is the extended model
with every \code{br_x = true} (\code{expectations.variables = 'all'}), $n_1 = 1$, zero beliefs, the
scenario's \code{re_variant} and \code{beliefs.info}, and the policy rule: the RE model, with every
\code{x_exp1_RE} in the smoothed history. Under \code{'lagged'} it keeps the scenario's learning set
(the RE(t-1) model of $\mathcal{L}$). Ramsey scenarios are always filtered with the rule model.

\paragraph{Outputs.} One call \code{filter(m, d, range, 'output', 'filter,smooth', 'meanOnly',
true)} (\code{'pred,filter,smooth'} for lagged models) gives the databanks the stages read:
\begin{itemize}[nosep]
  \item \code{H.smooth}: $x_{t|T}$, two-sided, every transition and measurement variable and shock;
  \item \code{H.filtered}: $x_{t|t}$, one-sided (data up to $t$);
  \item \code{H.predicted}: $x_{t|t-1}$, lagged models only.
\end{itemize}
The other fields of \code{H} hold the database (clipped in real time), the range, the fit error
and the cache information.
Observed measurement variables equal the data in every output; where every observable is missing,
the filtered value is the one-step prediction and the filtered shocks are 0.

\paragraph{Fit check.} The observables are rebuilt from the smoothed states and shocks through the
measurement equations (\code{lhsmrhs} with the data on the left). With at least as many shocks as
observables the largest absolute error must be below $10^{-8}$ (\code{hetlearn:re_filter:fit}).
Comparing the smoothed observables with the data would be vacuous, since IRIS copies the data into
them.

\paragraph{Cache.} \code{generated/refilter_<hash>.mat}, with the hash over the extended model text,
the parameter values, the assign struct (toggles), the database file bytes and the filter range;
in real time \code{refilter_rt_<date>_<hash>.mat}, with the clipped data values in place of the file
bytes, so a data change after the end date does not touch it. A cache that lacks the filtered
databank, or a lagged cache without predictions, is filtered again and rewritten.

\subsection{Learning on the RE history}\label{sec:resmoother}

With \code{history.mode = 're_smoother'} and \code{history.learn = true}, the sequence of
Section~\ref{sec:timing} runs for $t = t_0,\dots,T_0$ on the history databank $D$ of
\code{history.re_states} (\code{+engine/history.m}, \code{+engine/history_db.m}), with three
differences. Steps 2 and 3 are skipped: the model is never simulated, and it is solved only for
the ALM check of step 5 under \code{'alm'} and \code{'plm_alm'} (Section~\ref{sec:wherecheck}).
In step 4 the forecasts come from $D$: $\fRE_t(x) = D.\codem{x_exp1_RE}(t)$, and
$\fBR_t(x) = \mathcal{E}(x,1)$ of~\eqref{eq:recursion} with $\phi_{t-1}$ and the data $D.u(t+k)$,
$k\le k_{\max}$. In step 5 the realized value is $x_t = D.x(t)$. \code{res.paths.x_exp1_BR} holds the learners' forecasts over the
record dates (NaN before $t_0$), not the zero-belief values of the RE filter.

\subsubsection{\texorpdfstring{\codem{history.mode}}{history.mode}}\label{opt:history-mode}
\begin{optspec}
  \item[Default] \code{'re_smoother'}.
  \item[Values] \code{'re_smoother'}, \code{'inversion'}, \code{'kalman'}.
        \code{'hetero_filter'} is a reserved name (\code{hetlearn:options:notImplemented},
        pointing to \code{'kalman'}).
  \item[Math] Sections~\ref{sec:resmoother}, \ref{sec:inversion}, \ref{sec:kalman}.
  \item[Assumptions] \code{'re_smoother'}: the data are RE data, the agents' learning does not
        feed back into the history. \code{'inversion'}: the learning model explains the data
        exactly through the endogenized shocks. \code{'kalman'}: the learning model explains the
        data up to filtering uncertainty, quarter by quarter.
  \item[Interactions] \code{'inversion'} and \code{'kalman'} need \code{history.learn = true} and
        \code{history.re_states = 'smoothed'}; \code{history.realtime} applies to
        \code{'re_smoother'} only (Section~\ref{opt:history-realtime}); \code{hetlearn.decompose}
        refuses \code{'kalman'} results; \code{hetlearn.evaluate} runs one inversion per
        inversion variant.
  \item[Code] \code{+engine/core.m}, \code{+engine/history.m}, \code{+engine/kalman_history.m}.
  \item[Example]
\begin{lstlisting}
cfg.history.mode = 'inversion';
cfg.history.re_states = 'smoothed';
\end{lstlisting}
\end{optspec}

\subsubsection{\texorpdfstring{\codem{history.re_states}}{history.re\_states}}\label{opt:history-re-states}
\begin{optspec}
  \item[Default] \code{'smoothed'}. (Project baseline: \code{'filtered'}.)
  \item[Values] \code{'smoothed'}, \code{'filtered'}.
  \item[Math] $D$ = \code{H.smooth} or \code{H.filtered}. Under \code{'filtered'}
        $\fRE_t(x) = x^{\RE}_{t+1|t}$, the filter's one-step prediction of the RE model, and
        $\FE^{\RE}_t$ is its squared one-step prediction error; the learners' data and the
        realized $x_t$ are $x_{t|t}$. With \code{re_info = 'lagged'} the RE forecasts come from
        \code{H.predicted} instead (Section~\ref{sec:reinfo}).
  \item[Assumptions] \code{'smoothed'}: the RE forecast at $t$ uses states estimated with data
        after $t$, which gives RE agents hindsight in the fitness comparison. \code{'filtered'} removes
        it. For an observed variable the two coincide (filtered = smoothed = data); in a quarter
        without data the filtered value is the RE prediction, so $\FE^{\RE}_t = 0$ there.
  \item[Interactions] \code{'filtered'} with \code{'inversion'} or \code{'kalman'} is an error
        (\code{hetlearn:options:value}): both are one-sided already. It also sets the data of
        \code{history.learn = false} and of \code{beliefs.init.method = 'training_ols'}. The
        forecast still starts from the smoothed state at the origin; with
        \code{history.realtime = true} that state is real time too, which makes
        \code{'filtered'} plus \code{realtime} the fully real-time setup. The filtered values up
        to the origin do not depend on later data, so \code{realtime} does not change the history
        loop under \code{'filtered'}.
  \item[Code] \code{+engine/history_db.m}, \code{re_filter.m}.
  \item[Example]
\begin{lstlisting}
cfg.history.re_states = 'filtered';
cfg.history.realtime = true;                % fully real time
\end{lstlisting}
\end{optspec}

\subsubsection{\texorpdfstring{\codem{history.realtime}}{history.realtime}}\label{opt:history-realtime}
\begin{optspec}
  \item[Default] \code{false}. (\code{hetlearn.evaluate} sets it from \code{opts.realtime},
        default true.)
  \item[Values] true, false.
  \item[Math] with true, \code{hetlearn.run} calls \code{hetlearn.re_filter(..., 'end', T0)}:
        every database series is clipped at $T_0$ and the filter and smoother run over
        $[\text{filter start}, T_0]$, so $x_{t|T}$ becomes $x_{t|T_0}$.
  \item[Assumptions] no data after the origin reach the history or the starting state.
  \item[Interactions] applies only with \code{history.mode = 're_smoother'} (\code{run.m});
        under \code{'inversion'} (one-sided by construction, the full-sample smoother supplies
        only the state before the inversion start) and \code{'kalman'} the full-sample filter is
        used and the option has no effect. Real-time caches are kept per end date.
  \item[Code] \code{run.m}, \code{re_filter.m} (\code{clip_data}).
  \item[Example]
\begin{lstlisting}
cfg.history.realtime = true;
cfg.forecast.origin = qq(2015,1);
\end{lstlisting}
\end{optspec}

\subsubsection{\texorpdfstring{\codem{history.learn}}{history.learn}}\label{opt:history-learn}
\begin{optspec}
  \item[Default] \code{true}.
  \item[Values] true, false.
  \item[Math] false: no history loop; $\phi_{T_0} = \phi_0$, $R_{T_0} = R_0$, $n_{1,T_0} =
        1 - \bar n_2$, and the forecasts made at $T_0$ come from the same formulas at $t = T_0$ on
        the history databank. No update at the origin.
  \item[Assumptions] the initial beliefs and shares apply directly at the origin.
  \item[Interactions] the record starts at $T_0$; \code{beliefs.init.date} is not used;
        \code{'inversion'} and \code{'kalman'} are errors. The first forecast quarter is the
        first ALM check, so an unstable initial model stops with
        \code{hetlearn:engine:almUnstable}.
  \item[Code] \code{+engine/core.m}, \code{+engine/history.m}.
  \item[Example]
\begin{lstlisting}
cfg.history.learn = false;
cfg.beliefs.init.method = 'user';
\end{lstlisting}
\end{optspec}

\subsection{The inversion history}\label{sec:inversion}

With \code{history.mode = 'inversion'} the history is simulated one quarter at a time with the
learning model, as a forecast with conditions (\code{+engine/history.m}, function
\code{inversion}). Let $t_s$ = \code{history.inversion.start}. Before $t_s$ the state is the
RE-smoothed one; if $t_s > t_0$ the learners first learn on the RE-smoothed history over
$[t_0, t_s-1]$ (Section~\ref{sec:resmoother}). For $t = t_s,\dots,T_0$ the sequence of
Section~\ref{sec:timing} runs (stage \code{'inversion'}) with one difference: the simulation of
step 3 exogenizes every observable that has data at $t$ at its data value and endogenizes the
shocks of \code{history.inversion.shocks}. With $y^0_t$ the observables of the simulation with
zero shocks and $M$ their response to the endogenized shocks,
\begin{equation}\label{eq:inversion}
  M\delta = y^{\text{obs}}_t - y^0_t,
\end{equation}
solved exactly (as many shocks as observables) or with minimum norm in standard-deviation units
(more shocks; Section~\ref{sec:conditions}); observables without data at $t$ are not exogenized.
Every variable at $t$ is stored. The update at $t_s$ is the first-quarter update of
Section~\ref{sec:rls} only when $t_s = t_0$.

Quarter $t$ uses data up to $t$, so the inversion is one-sided by construction. The forecast starts
from the backed-out history: \code{+engine/core.m} replaces the RE-smoothed variables and shocks
over $[t_s, T_0]$ by the backed-out ones. Data come from \code{H.data}. Under mode
\code{'rational'} the inversion from the smoothed state reproduces the RE-smoothed shocks (test
target $10^{-8}$; $2\times10^{-13}$ on SW07 over 1966Q1--2026Q2 when every quarter has data). In a
quarter without data nothing is exogenized and the endogenized shocks
stay 0, so from such a quarter on the inversion departs from the smoother, which infers shocks
there from later data. A Ramsey scenario backs out its history with the rule model (untested).

\paragraph{Diagnostics} (\code{res.inversion}). \code{shocks} holds the endogenized shocks and
\code{params} the parameter values in the model each quarter, which \code{hetlearn.decompose}
uses. \code{shock_diff_vs_smoothed} and \code{path_diff_vs_smoothed} give the largest absolute
distance from the RE smoother (\code{x_exp1_BR} excluded), and \code{max_obs_miss} the largest
miss of an observable target. The other fields hold dates, counts and run times.

\subsubsection{\texorpdfstring{\codem{history.inversion.start}}{history.inversion.start}}\label{opt:history-inversion-start}
\begin{optspec}
  \item[Default] \code{[]}: \code{beliefs.init.date}.
  \item[Values] a date in $[t_0, T_0]$ (\code{hetlearn:engine:inversion} otherwise).
  \item[Math] $t_s$ above.
  \item[Assumptions] before $t_s$ the economy is the RE-smoothed one.
  \item[Interactions] \code{'inversion'} only; with $t_s > t_0$ the learners learn on the
        smoothed history first; the events after $t_s$ carry the stage \code{inversion}.
  \item[Code] \code{+engine/core.m}, \code{+engine/history.m}.
  \item[Example]
\begin{lstlisting}
cfg.history.mode = 'inversion';
cfg.history.inversion.start = qq(1985,1);
\end{lstlisting}
\end{optspec}

\subsubsection{\texorpdfstring{\codem{history.inversion.shocks}}{history.inversion.shocks}}\label{opt:history-inversion-shocks}
\begin{optspec}
  \item[Default] \code{{}}: every transition shock of the model.
  \item[Values] a cellstr of shocks; at least as many as observables.
  \item[Math] the shocks whose values $\delta$ solve~\eqref{eq:inversion}.
  \item[Assumptions] every data surprise is attributed to these shocks; the others stay 0.
  \item[Interactions] fewer shocks than observables or an unknown name is
        \code{hetlearn:engine:inversion}; a rank-deficient response is
        \code{hetlearn:engine:conditions}; with more shocks than observables
        \code{engine.conditions_weighting} sets the norm.
  \item[Code] \code{+engine/history.m}, \code{+engine/step.m}.
  \item[Example]
\begin{lstlisting}
cfg.history.mode = 'inversion';
cfg.history.inversion.shocks = {'ea', 'eb', 'eg', 'eqs', 'em', 'epinf', 'ew'};
\end{lstlisting}
\end{optspec}

\subsection{The Kalman history}\label{sec:kalman}

With \code{history.mode = 'kalman'} the history is a time-varying Kalman filter of the ALM in force
each quarter \citep[see][for the filter]{durbin2012}; beliefs and shares move with the same update
as everywhere else (\code{+engine/kalman_history.m}). For quarter $t$ from the start date to $T_0$:

\paragraph{(1) The ALM in force.} Before $t_0$ the RE limit ($n_1 = 1$); from $t_0$ the model of
steps 1 and 2 of Section~\ref{sec:timing} (stage \code{'kalman'}). Its state space on the full transition vector
$\xi = (x^f, x^b)$ is
\begin{equation}\label{eq:kalss}
  \xi_t = A\xi_{t-1} + \K + \Rs\varepsilon_t,\quad A = [0,\ \T],\qquad
  y_t = Z_f\xi_t + \Hs\varepsilon_t + \D,\quad Z_f = [0,\ \Z],\qquad \mathrm{Cov}(\varepsilon_t) = \Omega,
\end{equation}
with $\Omega$ from the model's \code{std_*} and \code{corr_*} parameters (a backward system is used
the same way).

\paragraph{(2) Prediction and update.} With $G = Z_f\Rs + \Hs$,
\begin{align}
  a_{t|t-1} &= A a_{t-1|t-1} + \K, & P_{t|t-1} &= A P_{t-1|t-1} A\transp + \Rs\Omega\Rs\transp, \notag\\
  \hat y_{t|t-1} &= Z_f a_{t|t-1} + \D, & C_t &= P_{t|t-1}Z_f\transp + \Rs\Omega\Hs\transp, \notag\\
  F_t &= Z_f C_t + \Hs\Omega G\transp, & v_t &= y_t - \hat y_{t|t-1}, \label{eq:kalman}
\end{align}
and, on the rows $o$ with data,
\begin{gather}
  \ell_t = -\tfrac12\Bigl(n_t\ln 2\pi + \ln\det F_{oo,t} + v_{o,t}\transp F_{oo,t}^{-1}v_{o,t}\Bigr),\\
  a_{t|t} = a_{t|t-1} + C_{o,t}F_{oo,t}^{-1}v_{o,t},\quad
  P_{t|t} = P_{t|t-1} - C_{o,t}F_{oo,t}^{-1}C_{o,t}\transp,\quad
  \varepsilon_{t|t} = \Omega G_o\transp F_{oo,t}^{-1}v_{o,t}.
\end{gather}
A quarter without data is a prediction only, with $\ell_t = 0$. A covariance $F_{oo}$ that is not
positive definite stops the run. The log likelihood is $\sum_{t\ge t_0}\ell_t$.

\paragraph{(3) The agents' information} (from $t_0$). RE forecasts made with data to $t-1$
(\code{re_info = 'lagged'}) are the predictions $\codem{x_exp1_RE}_{t|t-1}$; with data up to $t$ the
filtered $\codem{x_exp1_RE}_{t|t}$. The learners' PLM forecasts~\eqref{eq:recursion}, their
regressors and the realized values are the filtered values $x_{s|s}$ stored when each quarter was
filtered, in real time, not revised later.

\paragraph{(4) Update.} Step 5 of Section~\ref{sec:timing}, with the first-quarter update at
$t_0$.

\paragraph{The forecast} starts from the filtered state at $T_0$: the lags of the state vector at
the origin take their values in $a_{T_0|T_0}$.

\paragraph{Outputs} (\code{res.kalman}).

\begin{center}
\begin{tabular}{>{\raggedright\arraybackslash}p{5.0cm}>{\raggedright\arraybackslash}p{9.8cm}}
\toprule
Field & Content \\
\midrule
\code{dates}, \code{observables} & the filtered quarters and the measurement variables \\
\code{y}, \code{pred}, \code{err}, \code{F} & data $y_t$, predictions $\hat y_{t|t-1}$, errors $v_t$, covariances $F_t$ \\
\code{llt}, \code{n_obs} & $\ell_t$ and the number of observables with data, by quarter \\
\code{loglik}, \code{loglik_start} & $\sum_{t\ge t_0}\ell_t$ and its first quarter $t_0$ \\
\code{filtered}, \code{predicted} & databanks of $x_{t|t}$ and $x_{t|t-1}$ \\
\code{state_names}, \code{a_origin}, \code{P_origin} & the state vector $\xi$, $a_{T_0|T_0}$ and $P_{T_0|T_0}$ \\
\code{fbr_model_maxdiff} & the largest gap between the learners' forecast and the model's \code{x_exp1_BR} (below) \\
\code{n_solve}, \code{runtime_s} & solves and run time \\
\bottomrule
\end{tabular}
\end{center}

\code{fbr_model_maxdiff} compares $\fBR_t$ on the real-time data with the model's
\code{x_exp1_BR} (predicted under \code{'lagged'}, filtered under \code{'current'}). The gap is
rounding for a PLM whose forecast uses data
of one date only, such as an AR(1); for an AR(2) it is not, because the model's lagged states are
the filter's revised values and the agents' data the real-time ones. \code{hetlearn.run} writes
\code{kalman_errors.csv} (date, \code{n_obs}, \code{llt}, and \code{err_}, \code{pred_}, \code{var_}
of every measurement variable); \code{hetlearn.engine.kalman_history('rmse', res, subsamples)}
gives one-step RMSEs and log likelihoods by subsample.

\subsubsection{\texorpdfstring{\codem{history.kalman.init}}{history.kalman.init}}\label{opt:history-kalman-init}
\begin{optspec}
  \item[Default] \code{'stochastic'}.
  \item[Values] \code{'stochastic'}, \code{'smoothed'}.
  \item[Math] \code{'stochastic'}: the filter starts at the filter start date from the asymptotic
        distribution of the first quarter's ALM (the RE limit), as IRIS \code{filter} does:
        $a_0 = (I - \T_b)^{-1}\K_b$, $P_0$ the solution of $P = \T_b P \T_b\transp + Q_{bb}$
        (stationary roots only: a root of modulus $\ge 1-10^{-10}$ is an error).
        \code{'smoothed'}: the filter starts at $t_0$ from the RE-smoothed state of $t_0-1$ with
        $P_0 = 0$; the agents' data before $t_0$ are RE-smoothed.
  \item[Assumptions] \code{'smoothed'} with as many shocks as observables and no measurement error
        makes every update exact, so the filter equals the inversion until a quarter without data.
  \item[Interactions] \code{'kalman'} only; needs \code{history.learn = true} and
        \code{history.re_states = 'smoothed'}; under \code{projection.facility = 'none'} a quarter
        without a model solution stops with \code{hetlearn:engine:noSolution}.
  \item[Code] \code{+engine/kalman_history.m} (\code{initial_state}).
  \item[Example]
\begin{lstlisting}
cfg.history.mode = 'kalman';
cfg.history.kalman.init = 'smoothed';
\end{lstlisting}
\end{optspec}

% =====================================================================
\section{The forecast stage}\label{sec:forecast}
% =====================================================================

\subsection{The quarter-by-quarter loop}\label{sec:forecastloop}

\code{+engine/forecast.m} runs the sequence of Section~\ref{sec:timing} for
$t = T_0+1,\dots,T_0+H$, starting from the learning state at the end of $T_0$ and the databank of
\code{+engine/build_shocks.m}: the history up to $T_0$ (RE-smoothed, or backed out, or filtered,
depending on \code{history.mode}) and the scenario's shocks over the horizon. After each quarter the
simulated values are written into the databank, so quarter $t+1$ starts from the scenario's own
$t$. The result is an \emph{unconditional} forecast when the shocks are zero and there are no
conditions.

Three properties hold by construction and are tested:
\begin{itemize}[nosep]
  \item with mode \code{'rational'} the forecast equals one IRIS \code{simulate} of the RE model
        over the horizon with unanticipated shocks (\code{test_engine} E1);
  \item with mode \code{'rational'}, \code{shocks.source = 'smoothed'} and the smoothed starting
        state, the forecast reproduces the data (\code{test_sw07});
  \item the learners' forecast read from the model equals the PLM formula with $\phi_{t-1}$ in
        every quarter (the timing check).
\end{itemize}

\subsection{Origin and horizon}

\subsubsection{\texorpdfstring{\codem{forecast.origin}}{forecast.origin}}\label{opt:forecast-origin}
\begin{optspec}
  \item[Default] \code{[]} (must be set: \code{hetlearn:engine:origin}).
  \item[Values] an IRIS date inside the filter range, with room for the PLM lags after the filter
        start.
  \item[Math] $T_0$, the last quarter of data used; the forecast covers $T_0+1,\dots,T_0+H$.
  \item[Assumptions] without \code{history.realtime} the smoothed state at $T_0$ uses the whole
        sample, including data after $T_0$.
  \item[Interactions] \code{beliefs.init.date} must not be later (with learning);
        \code{history.inversion.start} must not be later; \code{hetlearn.evaluate} overrides it.
  \item[Code] \code{+engine/core.m}.
  \item[Example]
\begin{lstlisting}
cfg.forecast.origin = qq(2022,2);
cfg.forecast.horizon = 16;
\end{lstlisting}
\end{optspec}

\subsubsection{\texorpdfstring{\codem{forecast.horizon}}{forecast.horizon}}\label{opt:forecast-horizon}
\begin{optspec}
  \item[Default] \code{16}.
  \item[Values] a positive integer.
  \item[Math] $H$.
  \item[Assumptions] none.
  \item[Interactions] conditions and \code{shocks.add} dates must lie inside the horizon (dates of
        \code{shocks.range} and \code{density.range} outside it are ignored);
        \code{policy.loss_horizon} beyond $H$ triggers a continuation; \code{hetlearn.evaluate}
        sets it from \code{opts.horizon}.
  \item[Code] \code{+engine/core.m}.
  \item[Example]
\begin{lstlisting}
cfg.forecast.horizon = 8;
cfg.policy.loss_horizon = 80;
\end{lstlisting}
\end{optspec}

\subsection{Shocks over the horizon}\label{sec:shocks}

Every shock of the model is set to 0 over the horizon, then \code{shocks.source} and
\code{shocks.add} fill it (\code{+engine/build_shocks.m}). All shocks are unanticipated.

\subsubsection{\texorpdfstring{\codem{shocks.source}}{shocks.source}}\label{opt:shocks-source}
\begin{optspec}
  \item[Default] \code{'none'}.
  \item[Values] \code{'none'}, \code{'smoothed'}.
  \item[Math] \code{'none'}: $\varepsilon_t = 0$. \code{'smoothed'}: $\varepsilon_t =
        \varepsilon_{t|T}$, the RE-smoothed shocks (\code{H.smooth}) over \code{shocks.range},
        missing values set to 0.
  \item[Assumptions] \code{'smoothed'} replays the shocks the RE model attributes to the data.
        They are the RE-smoothed shocks under every history mode, not the backed-out ones.
  \item[Interactions] \code{shocks.range}, \code{shocks.exclude}; under Ramsey policy
        smoothed rule shocks run with no effect (warning \code{hetlearn:policy:ruleInputs}).
  \item[Code] \code{+engine/build_shocks.m}.
  \item[Example]
\begin{lstlisting}
cfg.shocks.source = 'smoothed';
cfg.shocks.exclude = {'em'};
\end{lstlisting}
\end{optspec}

\subsubsection{\texorpdfstring{\codem{shocks.range}}{shocks.range}}\label{opt:shocks-range}
\begin{optspec}
  \item[Default] \code{[]}: the horizon.
  \item[Values] a date vector, or \code{[start, end]}.
  \item[Math] the quarters that receive smoothed shocks; dates outside the horizon are ignored.
  \item[Assumptions] none.
  \item[Interactions] \code{shocks.source = 'smoothed'} only.
  \item[Code] \code{+engine/build_shocks.m}, \code{+util/to_dates.m}.
  \item[Example]
\begin{lstlisting}
cfg.shocks.source = 'smoothed';
cfg.shocks.range = [qq(2022,3), qq(2023,4)];
\end{lstlisting}
\end{optspec}

\subsubsection{\texorpdfstring{\codem{shocks.exclude}}{shocks.exclude}}\label{opt:shocks-exclude}
\begin{optspec}
  \item[Default] \code{{}}.
  \item[Values] a cellstr of shocks of the model (an unknown name is
        \code{hetlearn:engine:shocks}).
  \item[Math] the listed smoothed shocks are set to 0.
  \item[Assumptions] a counterfactual without those shocks.
  \item[Interactions] \code{shocks.source = 'smoothed'}; \code{shocks.add} still adds to an
        excluded shock.
  \item[Code] \code{+engine/build_shocks.m}.
  \item[Example]
\begin{lstlisting}
cfg.shocks.source = 'smoothed';
cfg.shocks.exclude = {'epinf', 'ew'};
\end{lstlisting}
\end{optspec}

\subsubsection{\texorpdfstring{\codem{shocks.add}}{shocks.add}}\label{opt:shocks-add}
\begin{optspec}
  \item[Default] \code{{}}.
  \item[Values] a cell array of triples \code{{name, dates, values}}: a shock, a date vector or
        \code{[start, end]}, and a scalar or one value per date.
  \item[Math] $\varepsilon_{j,t} \leftarrow \varepsilon_{j,t} + v_t$ at the given dates.
  \item[Assumptions] each added shock is a surprise in its quarter.
  \item[Interactions] adds to the source; under Ramsey a rule shock runs with no effect.
        Free-form subtree (checked by \code{build_shocks}: \code{hetlearn:engine:shocks}).
  \item[Code] \code{+engine/build_shocks.m}.
  \item[Example]
\begin{lstlisting}
cfg.shocks.add = {{'em', qq(2022,3), 0.25}, ...
                  {'epinf', [qq(2022,3), qq(2022,4)], [0.1 0.05]}};
\end{lstlisting}
\end{optspec}

\subsection{Conditions: exogenize and endogenize}\label{sec:conditions}

A condition puts a variable on a path in given quarters by endogenizing a shock. In quarter $t$
let $\mathcal{K}_t$ be the conditioned variables with targets $\tau_t$, $\mathcal{J}_t$ the
endogenized shocks, $s^0_t$ the values of the conditioned variables in the one-period simulation
with the scheduled shocks, and $M_t$ their response to a unit of each endogenized shock. The engine
solves (\code{+engine/step.m})
\begin{equation}\label{eq:cond}
  M_t\,\delta_t = \tau_t - s^0_t,\qquad
  \delta_t =
  \begin{cases}
    M_t^{-1}(\tau_t - s^0_t) & |\mathcal{J}_t| = |\mathcal{K}_t|\\
    S\,(M_tS)^{+}(\tau_t - s^0_t) & |\mathcal{J}_t| > |\mathcal{K}_t|,\ \text{\codem{'std'}}\\
    M_t^{+}(\tau_t - s^0_t) & |\mathcal{J}_t| > |\mathcal{K}_t|,\ \text{\codem{'none'}}
  \end{cases}
\end{equation}
with $S = \mathrm{diag}(\sigma_j)$, the standard deviations \code{std_<shock>}; the \code{'std'}
case is the minimum of $\sum_j(\delta_j/\sigma_j)^2$ subject to the targets. Then
$\varepsilon_{j,t}\leftarrow\varepsilon_{j,t}+\delta_{j,t}$ and the quarter is simulated again; every
target must be met within $10^{-9}$. Fewer shocks than targets, or $\mathrm{rank}\,M_t <
|\mathcal{K}_t|$, is \code{hetlearn:engine:conditions}. Under \code{engine.step = 'iris'} $M_t$
comes from unit-shock simulations; under \code{'fast'} it is exact: the rows of $\Rs$ for
transition variables and $\Z\Rs_b + \Hs$ for measurement variables. The procedure equals an IRIS
\code{Plan} with exogenize/endogenize pairs run one quarter at a time (\code{test_engine} E3).

The conditions are unanticipated: agents learn in quarter $t$ that the variable was hit, and the
endogenized shocks enter the data the learners see. A condition therefore moves beliefs and shares
in later quarters.

\subsubsection{\texorpdfstring{\codem{conditions}}{conditions}}\label{opt:conditions}
\begin{optspec}
  \item[Default] an empty struct array.
  \item[Values] a struct array, one element per pair, with fields
        \begin{itemize}[nosep]
          \item \code{variable}: a transition or measurement variable;
          \item \code{shock}: a shock of the model;
          \item \code{dates}: dates inside the horizon, or \code{[start, end]};
          \item \code{values}: a scalar or one value per date;
          \item \code{type}: \code{'level'} (default) or \code{'delta'};
          \item \code{reference}: \code{'self'} (default) or a scenario id.
        \end{itemize}
  \item[Math] the targets $\tau_t$ of~\eqref{eq:cond}: the value itself (\code{'level'}), or
        the reference path plus the value (\code{'delta'}).
  \item[Assumptions] \code{'self'} is the same scenario with every condition removed and
        \code{compare_with} emptied, run first and not saved (an inversion history is reused);
        a scenario id is read from \code{results_dir/<id>/results.mat}, and run first when it has
        no results.
  \item[Interactions] the option is replaced as a whole, not merged. A variable conditioned twice
        in a quarter, a date outside the horizon, an unknown variable or shock, or a values/dates
        mismatch are errors (\code{hetlearn:engine:conditions}). Under Ramsey policy, quarters
        whose endogenized shocks are all rule shocks keep only rule-only targets
        (Section~\ref{sec:rulenames}). In density draws the conditions stay in force, so
        endogenized shocks offset draws on the same shock. The density fast path needs no
        conditions.
  \item[Code] \code{+engine/conditions_prepare.m}, \code{+engine/step.m},
        \code{+engine/core.m} (references).
  \item[Example]
\begin{lstlisting}
cfg.conditions = struct('variable', 'ms', 'shock', 'em', 'dates', [qq(2022,3), qq(2023,2)], ...
    'values', -0.25, 'type', 'delta', 'reference', 'self');
\end{lstlisting}
\end{optspec}

\subsubsection{\texorpdfstring{\codem{engine.conditions_weighting}}{engine.conditions\_weighting}}\label{opt:engine-conditions-weighting}
\begin{optspec}
  \item[Default] \code{'std'}.
  \item[Values] \code{'std'}, \code{'none'}.
  \item[Math] the two minimum-norm cases of~\eqref{eq:cond}.
  \item[Assumptions] \code{'std'}: the smallest shocks in standard-deviation units are the most
        plausible.
  \item[Interactions] matters only with more endogenized shocks than targets (conditions, the
        inversion); \code{'std'} needs a positive \code{std_<shock>} for every endogenized shock.
  \item[Code] \code{+engine/step.m} (\code{solve_shocks}).
  \item[Example]
\begin{lstlisting}
cfg.engine.conditions_weighting = 'none';
cfg.history.mode = 'inversion';
\end{lstlisting}
\end{optspec}

\subsection{The one-period step}\label{sec:enginestep}

\subsubsection{\texorpdfstring{\codem{engine.step}}{engine.step}}\label{opt:engine-step}
\begin{optspec}
  \item[Default] \code{'fast'}.
  \item[Values] \code{'fast'}, \code{'iris'}.
  \item[Math] \code{'fast'}: quarter $t$ from the state-space form~\eqref{eq:sspace} of the solved
        model, cached after every successful solve:
        $\xi_t = \T x^b_{t-1} + \Rs\varepsilon_t + \K$, $y_t = \Z x^b_t + \Hs\varepsilon_t + \D$,
        with $x^b_{t-1}$ read from the databank (element \code{v{-k}} is \code{v} at $t-1-k$) and
        missing shocks set to 0. \code{'iris'}: \code{simulate(m, d, t:t, 'prependInput', true)}.
  \item[Assumptions] none; the two agree to rounding ($10^{-14}$ per quarter, about $10^{-12}$
        over 16 quarters).
  \item[Interactions] a stored backward system has priority over both. \code{'fast'} needs plain
        state names. Bit-identity tests pin \code{'iris'}. The fast step is 2--3 times faster
        without conditions and about 20 times with the observables exogenized (the inversion).
  \item[Code] \code{+engine/step.m}, \code{+engine/solve_model.m}, \code{+engine/core.m}.
  \item[Example]
\begin{lstlisting}
cfg.engine.step = 'iris';
cfg.history.mode = 'inversion';
\end{lstlisting}
\end{optspec}

\subsection{Identification and comparisons}

\subsubsection{\texorpdfstring{\codem{id}, \codem{title}, \codem{description}}{id, title, description}}\label{opt:id}
\begin{optspec}
  \item[Default] \code{''}.
  \item[Values] text. \code{id} must be set (\code{hetlearn:run:id}).
  \item[Math] none.
  \item[Assumptions] \code{id} names the results folder \code{results/<id>/}; scenario ids are
        function names, found by \code{compare_with} and \code{'delta'} references, so they must
        be unique on the path.
  \item[Interactions] \code{hetlearn.evaluate} requires unique variant ids.
  \item[Code] \code{run.m}, \code{+engine/save_results.m}.
  \item[Example]
\begin{lstlisting}
cfg.id = 'my_learn';
cfg.title = 'Switching on inflation';
\end{lstlisting}
\end{optspec}

\subsubsection{\texorpdfstring{\codem{compare_with}}{compare\_with}}\label{opt:compare-with}
\begin{optspec}
  \item[Default] \code{{}}.
  \item[Values] a cellstr of scenario ids.
  \item[Math] \code{compare_<id>.csv}: this run minus that run, every observable, over the
        horizon, model units.
  \item[Assumptions] the other run has the same horizon dates.
  \item[Interactions] a target without results is run first (\code{feval(id)}); a target that
        cannot be run leaves its file unwritten with a warning
        (\code{hetlearn:save_results:compare}). \code{hetlearn.evaluate} empties it.
  \item[Code] \code{+engine/save_results.m}, \code{+engine/load_result.m}.
  \item[Example]
\begin{lstlisting}
cfg.compare_with = {'a1_re'};
cfg.id = 'a1_learn_variant';
\end{lstlisting}
\end{optspec}

\subsubsection{\texorpdfstring{\codem{report.variables}}{report.variables}}\label{opt:report-variables}
\begin{optspec}
  \item[Default] \code{{}}.
  \item[Values] a cellstr; \code{report} is a free-form subtree.
  \item[Math] none. Reserved: converted to a cellstr and read by no toolbox code.
  \item[Assumptions] none.
  \item[Interactions] none.
  \item[Code] \code{load_options.m}.
  \item[Example]
\begin{lstlisting}
cfg.report.variables = {'pinfobs', 'robs'};
cfg.report.note = 'free-form';
\end{lstlisting}
\end{optspec}

% =====================================================================
\section{Monetary policy: rules and Ramsey}\label{sec:policy}
% =====================================================================

\subsection{The policy section of a package}\label{sec:specpolicy}

Ramsey policy needs \code{spec.policy} with four fields (\code{extend_model.m},
\code{policy_setup}):

\begin{itemize}[nosep]
  \item \code{instrument}: a transition variable (SW07: \code{r});
  \item \code{rule_label}: the exact, unique label of the policy rule in
        \code{!transition_equations};
  \item \code{targets}: a struct of target expressions (names, numbers, \code{+ - * / ^},
        parentheses and lags \code{x{-k}}; at least one transition variable). SW07:
        \code{pinf = 4*pinf}, \code{ygap = y - yf}, \code{r = 4*r}, \code{dr = 4*(r - r{-1})};
  \item \code{discount}: a number or an expression, in which substitutions written as
        \code{$name$} are expanded; SW07 uses \code{$cbetabar$*$cgamma$}.
\end{itemize}

The generator wraps the labeled rule:
\begin{lstlisting}
!if ramsey
    min((1 - ramsey_disc_on)*(<discount>) + ramsey_disc_on*ramsey_disc) w_pinf*(4*pinf)^2 + w_ygap*(y - yf)^2 + ...;
!else
    <the policy rule>;
!end
\end{lstlisting}
and declares \code{w_<k>}, \code{ramsey_disc_on} and \code{ramsey_disc} unconditionally. The
reference file is never wrapped; the collapse test runs with \code{ramsey = false}. A label that is
not found sets \code{meta.policy.ok = false} and is an error only for \code{policy.mode = 'ramsey'}
(\code{hetlearn:policy:rule}). Without \code{spec.policy} there is no Ramsey mode and no loss
accounting.

\subsection{The planner's problem}\label{sec:ramsey}

With \code{policy.mode = 'ramsey'} the model is loaded again with \code{ramsey = true} and
\code{Model.fromFile(..., 'optimal', {'type', policy.type})}, so IRIS derives the first-order
conditions of
\begin{equation}\label{eq:ramsey}
  \min\; E_t\sum_{h\ge0}\delta^h\sum_k w_k\,\bigl(\text{target}_{k,t+h}\bigr)^2
\end{equation}
subject to every other model equation, including the expectation block with the beliefs and shares
in force, which the planner, like every agent, treats as permanent (Section~\ref{sec:timing}). The Lagrange
multipliers of the plan (IRIS \code{Mu_Eq<k>}) are new transition variables of the Ramsey model.
Under \code{'commitment'} the planner honors past promises through their lags; under
\code{'discretion'} it re-optimizes each quarter and takes the private sector's leads as given. The toolbox ships a
corrected copy of one IRIS file, \code{toolbox/iris_patches/+model/Insertable.m}, which
\code{setup_paths} puts first on the path, because the optimal-policy parser of IRIS 20230622 fails
on MATLAB R2025a.

\paragraph{Timing.} The history is always the rule model: the RE filter is run with the rule and
the RE-moment initial beliefs come from the rule model (for anchored RE agents see
Section~\ref{opt:policy-mode}). The plan starts at the origin: the multipliers are 0 at and before $T_0$
(\code{+policy/prepare.m}, \code{+policy/forecast_setup.m}), and their values are stored every
quarter so their lags carry forward. With learning the planner re-optimizes each quarter with the
current beliefs and shares and honors past multipliers.

\subsection{Rule-only names}\label{sec:rulenames}

The generator finds the variables used only by the rule and their own laws of motion, and the shocks
used only there (SW07: \code{ms}, \code{em}; \code{meta.policy.rule_vars}, \code{.rule_shocks}).
Under Ramsey policy they move nothing else. Shocks and conditions on them still run, with one
warning per run (\code{hetlearn:policy:ruleInputs}); in quarters whose endogenized shocks are all
rule shocks, only conditions on rule-only variables are kept (a target on the instrument cannot be
met through a rule shock once the planner sets the instrument).

\subsection{Loss accounting}\label{sec:loss}

Every run of a package with \code{spec.policy}, rule and Ramsey alike, reports
(\code{+policy/loss.m})
\begin{equation}\label{eq:loss}
  \mathcal{L} = \sum_{h=1}^{N}\delta^h\sum_k w_k\,\bigl(\text{target}_{k,T_0+h}\bigr)^2,
  \qquad
  \mathcal{L}_k = w_k\sum_{h=1}^{N}\delta^h\,\bigl(\text{target}_{k,T_0+h}\bigr)^2,
\end{equation}
with the weights of \code{policy.loss} and $N$ = \code{policy.loss_horizon} (default $H$). If
$N > H$ the scenario is continued from \code{res.state} for $N-H$ quarters with zero shocks, no
conditions and learning on (a failed continuation gives NaN and a warning). Rule and Ramsey runs are
therefore compared on the same weights. \code{res.loss} holds the total, the contribution of each
target, the target paths, the per-quarter loss and the discount factors; \code{loss.csv} has one row
per target and a total row.

\subsection{Options}

\subsubsection{\texorpdfstring{\codem{policy.mode}}{policy.mode}}\label{opt:policy-mode}
\begin{optspec}
  \item[Default] \code{'rule'}.
  \item[Values] \code{'rule'}, \code{'ramsey'}.
  \item[Math] \code{'rule'}: the package's rule; \code{'ramsey'}: problem~\eqref{eq:ramsey} from
        the origin.
  \item[Assumptions] see Section~\ref{sec:ramsey}.
  \item[Interactions] \code{'ramsey'} needs \code{spec.policy} with a found label
        (\code{hetlearn:policy:noSpec}, \code{hetlearn:policy:rule}); it is not supported with
        \code{expectations.re_info = 'lagged'} (\code{hetlearn:options:unsupported}: the plan
        with predetermined RE expectations has not been derived) nor with the N-type options;
        \code{hetlearn.decompose} refuses Ramsey results. Anchored RE agents keep the Taylor-rule
        forecast, so mode \code{'rational'} with anchored agents is not an RE Ramsey plan; the RE
        benchmark is \code{re_variant = 'model_consistent'}.
  \item[Code] \code{load_model.m}, \code{+policy/setup.m}, \code{+policy/prepare.m},
        \code{+policy/forecast_setup.m}.
  \item[Example]
\begin{lstlisting}
cfg.policy.mode = 'ramsey';
cfg.policy.type = 'commitment';
\end{lstlisting}
\end{optspec}

\subsubsection{\texorpdfstring{\codem{policy.type}}{policy.type}}\label{opt:policy-type}
\begin{optspec}
  \item[Default] \code{'commitment'}.
  \item[Values] \code{'commitment'}, \code{'discretion'}.
  \item[Math] IRIS optimal policy of the given type.
  \item[Assumptions] commitment: the planner honors past promises (multipliers); discretion:
        re-optimizes taking expectations as given.
  \item[Interactions] \code{'ramsey'} only. Discretion is indeterminate in IRIS
        (\code{MULTIPLE_STABLE}) at some parameter points: SW07 under RE, and any learning
        variable at the RE point where the model is loaded; interior shares solve.
  \item[Code] \code{load_model.m}.
  \item[Example]
\begin{lstlisting}
cfg.policy.mode = 'ramsey';
cfg.policy.type = 'discretion';
\end{lstlisting}
\end{optspec}

\subsubsection{\texorpdfstring{\codem{policy.loss}}{policy.loss}}\label{opt:policy-loss}
\begin{optspec}
  \item[Default] \code{struct('pinf', 1, 'ygap', 0.5, 'r', 0.1, 'dr', 0)}.
  \item[Values] a struct of non-negative finite weights by target name of \code{spec.policy}.
  \item[Math] $w_k$ in~\eqref{eq:ramsey} and~\eqref{eq:loss}.
  \item[Assumptions] a quadratic loss in the targets.
  \item[Interactions] replaced as a whole, not merged: a target left out gets weight 0; an
        unknown name is \code{hetlearn:policy:unknownTarget}. The weights are model parameters
        (\code{w_<k>}), so changing them needs no new file.
  \item[Code] \code{+policy/setup.m}, \code{+policy/loss.m}.
  \item[Example]
\begin{lstlisting}
cfg.policy.mode = 'ramsey';
cfg.policy.loss = struct('pinf', 1, 'ygap', 0.25, 'r', 0.5);   % dr gets 0
\end{lstlisting}
\end{optspec}

\subsubsection{\texorpdfstring{\codem{policy.discount}}{policy.discount}}\label{opt:policy-discount}
\begin{optspec}
  \item[Default] \code{[]}: \code{spec.policy.discount} evaluated at the parameters.
  \item[Values] \code{[]} or a number in $(0,1]$.
  \item[Math] $\delta$ in~\eqref{eq:ramsey} and~\eqref{eq:loss}; in the model file the discount is
        $(1-\codem{ramsey_disc_on})\cdot\delta_{\text{spec}} + \codem{ramsey_disc_on}\cdot\codem{ramsey_disc}$,
        with \code{ramsey_disc_on} = 1 when the option is set.
  \item[Assumptions] none.
  \item[Interactions] used by the loss of rule runs too.
  \item[Code] \code{+policy/setup.m}, \code{+policy/eval_expr.m}.
  \item[Example]
\begin{lstlisting}
cfg.policy.mode = 'ramsey';
cfg.policy.discount = 0.99;
\end{lstlisting}
\end{optspec}

\subsubsection{\texorpdfstring{\codem{policy.loss_horizon}}{policy.loss\_horizon}}\label{opt:policy-loss-horizon}
\begin{optspec}
  \item[Default] \code{[]}: \code{forecast.horizon}.
  \item[Values] \code{[]} or a positive integer.
  \item[Math] $N$ in~\eqref{eq:loss}.
  \item[Assumptions] beyond $H$ the scenario continues with zero shocks and no conditions, with
        learning on.
  \item[Interactions] the continuation uses the forecast model and the end state; it is reported
        in \code{res.loss.continuation}.
  \item[Code] \code{+policy/loss.m}.
  \item[Example]
\begin{lstlisting}
cfg.forecast.horizon = 16;
cfg.policy.loss_horizon = 80;
\end{lstlisting}
\end{optspec}

% =====================================================================
\section{Density forecasts}\label{sec:density}
% =====================================================================

With \code{density.enabled = true}, \code{hetlearn.run} runs the history once and then repeats the
forecast stage for each of $N$ shock draws from the same origin state (\code{+density/run.m}).

\paragraph{Draws} (\code{+density/draws.m}). For the shocks $s$ of \code{density.shocks} and the
quarters of \code{density.range} inside the horizon, draw $j$ adds $E^{(j)}_{t,s}$ to the
scenario's own shock:
\begin{equation}\label{eq:draws}
  \varepsilon^{(j)}_{s,t} = \varepsilon^{\text{scen}}_{s,t} + E^{(j)}_{t,s},\qquad
  E^{(j)}_{t,s} =
  \begin{cases}
    \sigma_s\, z^{(j)}_{t,s},\ z \sim N(0,1)\ \text{i.i.d.} & \text{\codem{'normal'}}\\
    \bigl(p_{\iota^{(j)}_t} - \kappa\,\bar p\bigr)_s,\ \iota^{(j)}_t\sim U\{1,\dots,P\} & \text{\codem{'bootstrap'}}
  \end{cases}
\end{equation}
where $\sigma_s$ is the model's \code{std_<shock>} (\code{'calibrated'}) or the sample standard
deviation (normalized by the number of values minus one) of the RE-smoothed shock over
\code{density.sigma_range}
(\code{'smoothed'}), and $p_1,\dots,p_P$ are the vectors of RE-smoothed values of all drawn shocks at
the dates of \code{sigma_range} where all are finite, with mean $\bar p$; $\kappa = 1$ with \code{density.demean} and 0 otherwise. The bootstrap draws whole
dates, jointly across shocks, with replacement. Random numbers come from
\code{RandStream('mt19937ar', 'Seed', seed)}, draw by draw, so the first $k$ draws do not depend on
$N$ and the global stream is not touched. Draws are unanticipated, like every shock.

\paragraph{Paths.} The \emph{fast path} (\code{+density/fast_path.m}) applies when
\code{density.fast_path} is true, the mode is \code{'rational'} and there are no conditions: at
$n_1 = 1$ the solution does not depend on the beliefs or shares except through \code{x_exp1_BR},
which feeds nothing, so one \code{simulate(m, d, horizon, 'anticipate', false)} gives every other
variable; the learning diagnostics are then replayed on the simulated path. It equals the general
path to rounding (\code{test_density} D4). The \emph{general path} calls
\code{+engine/forecast.m} per draw, reusing the point forecast's first-quarter solve, so beliefs and
shares respond to the drawn shocks. Conditions keep their targets in every draw.

\paragraph{Exclusions.} A draw with a non-finite value, or any absolute value above $10^6$, in a
summarized series is \code{'explosive'}; a draw whose forecast stage stops (for example
\code{hetlearn:engine:almUnstable} or \code{noSolution}) is \code{'error'}. Both are left out of the
summaries and counted.

\paragraph{Summaries.} For every observable, the learners' model share \code{n2_model_<x>}
$= 1 - n^m_{1,t}$ and the belief parameters of each learning variable's PLM equation: the quantiles,
the mean, the standard deviation and the point forecast (the run without draws), in model units and,
for observables, in report units. The quantile of probability $p$ (\code{+density/quantile.m}, the
definition of MATLAB's Statistics Toolbox) places the sorted values $x_{(1)}\le\dots\le x_{(N)}$ at
probabilities $(i - 0.5)/N$ and interpolates linearly:
\begin{equation}\label{eq:quantile}
  q(p) = (1-w)\,x_{(\lfloor\pi\rfloor)} + w\,x_{(\lfloor\pi\rfloor+1)},\qquad
  \pi = \min\bigl(\max(Np + 0.5,\,1),\,N\bigr),\quad w = \pi - \lfloor\pi\rfloor .
\end{equation}

\paragraph{Outputs.} \code{res.density} (in memory, not in \code{results.mat}),
\code{density.mat} (variable \code{density}), \code{density_quantiles.csv} (variable, date,
quantile, \code{value_model}, \code{value_report}), \code{density_events.csv} (one row per draw:
draw, status, the counts of \code{projection_plm}, \code{projection_alm},
\code{projection_alm_shares}, \code{explosive_events}, \code{rls_missing}, \code{shares_nan},
\code{gain_bound}, then \code{max_abs} and the error message). A run without density deletes the
density files of an earlier run of the same id.

\subsection{Options}

\subsubsection{\texorpdfstring{\codem{density.enabled}}{density.enabled}}\label{opt:density-enabled}
\begin{optspec}
  \item[Default] \code{false}.
  \item[Values] true, false.
  \item[Math] switches the draws~\eqref{eq:draws} on.
  \item[Assumptions] the history is computed once; only the forecast stage is repeated.
  \item[Interactions] needs \code{density.shocks} (\code{hetlearn:options:missing});
        \code{hetlearn.evaluate} switches it off.
  \item[Code] \code{run.m}, \code{+density/run.m}.
  \item[Example]
\begin{lstlisting}
cfg.density.enabled = true;
cfg.density.shocks = {'epinf', 'ew'};
\end{lstlisting}
\end{optspec}

\subsubsection{\texorpdfstring{\codem{density.shocks}}{density.shocks}}\label{opt:density-shocks}
\begin{optspec}
  \item[Default] \code{{}}.
  \item[Values] a cellstr of shocks of the model (\code{hetlearn:density:shocks} otherwise).
  \item[Math] the shocks $s$ in~\eqref{eq:draws}; the others keep the scenario's values.
  \item[Assumptions] uncertainty comes from these shocks only.
  \item[Interactions] required when enabled; a shock that is also endogenized by a condition is
        offset by the condition.
  \item[Code] \code{+density/draws.m}.
  \item[Example]
\begin{lstlisting}
cfg.density.enabled = true;
cfg.density.shocks = {'eb', 'eg', 'eqs'};
\end{lstlisting}
\end{optspec}

\subsubsection{\texorpdfstring{\codem{density.sigma}}{density.sigma}}\label{opt:density-sigma}
\begin{optspec}
  \item[Default] \code{'calibrated'}.
  \item[Values] \code{'calibrated'}, \code{'smoothed'}.
  \item[Math] $\sigma_s$ in~\eqref{eq:draws}: \code{std_<shock>} of the model, or the sample
        standard deviation of the RE-smoothed shock over \code{sigma_range}.
  \item[Assumptions] normal, independent shocks.
  \item[Interactions] \code{density.method = 'normal'} only; \code{'smoothed'} needs at least two
        smoothed values (\code{hetlearn:density:sigma}).
  \item[Code] \code{+density/draws.m}.
  \item[Example]
\begin{lstlisting}
cfg.density.sigma = 'smoothed';
cfg.density.sigma_range = [qq(1966,1), qq(2019,4)];
\end{lstlisting}
\end{optspec}

\subsubsection{\texorpdfstring{\codem{density.sigma_range}}{density.sigma\_range}}\label{opt:density-sigma-range}
\begin{optspec}
  \item[Default] \code{[]}: the filter range.
  \item[Values] a date vector or \code{[start, end]}.
  \item[Math] the sample of smoothed shocks for \code{'smoothed'} and the bootstrap pool.
  \item[Assumptions] the shocks of this sample are representative of the future.
  \item[Interactions] \code{'smoothed'} and \code{'bootstrap'}; the sample means of the drawn shocks
        over it are reported in \code{draw_info.sample_mean}.
  \item[Code] \code{+density/draws.m}.
  \item[Example]
\begin{lstlisting}
cfg.density.method = 'bootstrap';
cfg.density.sigma_range = [qq(1985,1), qq(2019,4)];
\end{lstlisting}
\end{optspec}

\subsubsection{\texorpdfstring{\codem{density.method}}{density.method}}\label{opt:density-method}
\begin{optspec}
  \item[Default] \code{'normal'}.
  \item[Values] \code{'normal'}, \code{'bootstrap'}.
  \item[Math] the two cases of~\eqref{eq:draws}.
  \item[Assumptions] \code{'bootstrap'} keeps the joint distribution of the shocks of one date
        (their contemporaneous correlation and fat tails), not their serial dependence.
  \item[Interactions] \code{'bootstrap'} ignores \code{density.sigma}; it needs at least two pool
        dates (\code{hetlearn:density:bootstrap}).
  \item[Code] \code{+density/draws.m}.
  \item[Example]
\begin{lstlisting}
cfg.density.method = 'bootstrap';
cfg.density.demean = true;
\end{lstlisting}
\end{optspec}

\subsubsection{\texorpdfstring{\codem{density.demean}}{density.demean}}\label{opt:density-demean}
\begin{optspec}
  \item[Default] \code{true}.
  \item[Values] true, false.
  \item[Math] subtracts $\bar p$ from the pool in~\eqref{eq:draws}.
  \item[Assumptions] the smoothed shocks have non-zero sample means; demeaning centers the density
        on the scenario's own shocks.
  \item[Interactions] \code{'bootstrap'} only.
  \item[Code] \code{+density/draws.m}.
  \item[Example]
\begin{lstlisting}
cfg.density.method = 'bootstrap';
cfg.density.demean = false;
\end{lstlisting}
\end{optspec}

\subsubsection{\texorpdfstring{\codem{density.ndraws}}{density.ndraws}}\label{opt:density-ndraws}
\begin{optspec}
  \item[Default] \code{200}.
  \item[Values] a positive integer.
  \item[Math] $N$.
  \item[Assumptions] none.
  \item[Interactions] the general path costs one forecast stage per draw (SW07: about 0.7 s per
        draw over 16 quarters; the fast path about 0.07 s).
  \item[Code] \code{+density/draws.m}, \code{+density/run.m}.
  \item[Example]
\begin{lstlisting}
cfg.density.ndraws = 500;
cfg.density.seed = 1;
\end{lstlisting}
\end{optspec}

\subsubsection{\texorpdfstring{\codem{density.seed}}{density.seed}}\label{opt:density-seed}
\begin{optspec}
  \item[Default] \code{0}.
  \item[Values] a non-negative integer.
  \item[Math] the seed of the run's own \code{RandStream}.
  \item[Assumptions] none.
  \item[Interactions] the same seed gives the same draws in RE and learning runs, so their
        densities differ only through the model.
  \item[Code] \code{+density/draws.m}.
  \item[Example]
\begin{lstlisting}
cfg.density.seed = 42;
\end{lstlisting}
\end{optspec}

\subsubsection{\texorpdfstring{\codem{density.quantiles}}{density.quantiles}}\label{opt:density-quantiles}
\begin{optspec}
  \item[Default] \code{[0.05 0.16 0.5 0.84 0.95]}.
  \item[Values] a vector of numbers in $(0,1)$.
  \item[Math] the probabilities $p$ of~\eqref{eq:quantile}.
  \item[Assumptions] none.
  \item[Interactions] none.
  \item[Code] \code{+density/quantile.m}.
  \item[Example]
\begin{lstlisting}
cfg.density.quantiles = [0.1 0.25 0.5 0.75 0.9];
\end{lstlisting}
\end{optspec}

\subsubsection{\texorpdfstring{\codem{density.range}}{density.range}}\label{opt:density-range}
\begin{optspec}
  \item[Default] \code{[]}: the horizon.
  \item[Values] a date vector or \code{[start, end]}; dates outside the horizon are ignored
        (none inside is \code{hetlearn:density:range}).
  \item[Math] the quarters $t$ in~\eqref{eq:draws}.
  \item[Assumptions] uncertainty only in these quarters.
  \item[Interactions] none.
  \item[Code] \code{+density/draws.m}.
  \item[Example]
\begin{lstlisting}
cfg.density.range = [qq(2022,3), qq(2023,2)];
\end{lstlisting}
\end{optspec}

\subsubsection{\texorpdfstring{\codem{density.keep_draws}}{density.keep\_draws}}\label{opt:density-keep-draws}
\begin{optspec}
  \item[Default] \code{false}.
  \item[Values] true, false.
  \item[Math] none: stores every draw (\code{density.draws.series.<name>}, $H\times N$, and the
        drawn shocks) in \code{density.mat}.
  \item[Assumptions] none.
  \item[Interactions] file size grows with $N$.
  \item[Code] \code{+density/run.m}.
  \item[Example]
\begin{lstlisting}
cfg.density.keep_draws = true;
cfg.density.ndraws = 100;
\end{lstlisting}
\end{optspec}

\subsubsection{\texorpdfstring{\codem{density.fast_path}}{density.fast\_path}}\label{opt:density-fast-path}
\begin{optspec}
  \item[Default] \code{true}.
  \item[Values] true, false.
  \item[Math] one linear simulation per draw instead of the quarter-by-quarter loop, when the mode
        is \code{'rational'} and there are no conditions.
  \item[Assumptions] at $n_1=1$ beliefs and shares do not affect the path.
  \item[Interactions] false forces the general path (used to test their equality).
  \item[Code] \code{+density/fast_path.m}.
  \item[Example]
\begin{lstlisting}
cfg.expectations.mode = 'rational';
cfg.density.fast_path = false;
\end{lstlisting}
\end{optspec}

% =====================================================================
\section{Historical shock decompositions}\label{sec:decompose}
% =====================================================================

\code{D = hetlearn.decompose(res, vars, range, opts)} (\code{decompose.m}) decomposes the deviation
of each variable in \code{vars} (default \code{{'pinfobs', 'dy'}}, SW07 names) from its RE steady
state into one component per shock (or group), initial conditions and an intercept channel.
\code{res} is a result struct or a saved scenario id; \code{range} the reporting dates (default:
every date from the first one with initial conditions in \code{res.paths} to the end of the
horizon).

\paragraph{The RE steady state} comes from the RE model's state space:
$x^b_{\text{ss}} = (I-\T_b)^{-1}\K_b$, $\xi_{\text{ss}} = \T x^b_{\text{ss}} + \K$,
$y_{\text{ss}} = \Z x^b_{\text{ss}} + \D$.

\paragraph{Method \codem{'iris'}.} IRIS contributions of the RE model,
\code{simulate(m, res.paths, range, 'anticipate', false, 'contributions', true)}, from the first date
with initial conditions; \code{init} is the \code{Init+Const+Trends} column minus the steady state;
\code{beliefs_constants} is the rest (the \code{Nonlinear} column and rounding; 0 for a linear RE
model).

\paragraph{Method \codem{'tv'}.} A time-varying linear decomposition with the model in force in each
quarter: the RE model ($n_1=1$, zero beliefs) before the learning model applies (the RE-smoothed
history), the parameter values recorded by the inversion (\code{res.inversion.params}) over an
inversion history, and over the horizon the beliefs $\phi_{t-1}$ and model shares $n^m_{1,t}$ from
\code{res.learning} (the beliefs in force after the projection facility, reverted ones included).
The model is re-solved when the parameters change. With $\T_t,\Rs_t,\Z_t,\Hs_t$ its state-space
matrices (or its backward system),
\begin{equation}\label{eq:tv}
  \xi^c_t = \T_t x^{b,c}_{t-1} + \Rs_t\varepsilon^c_t, \qquad y^c_t = \Z_t x^{b,c}_t + \Hs_t\varepsilon^c_t,
\end{equation}
without constants, for $c$ = each shock ($\varepsilon^c$ = that shock only, zero start) and
$c$ = \code{init} (no shocks, start = the deviation of the initial state from the steady state). Then
\begin{equation}
  \text{\codem{beliefs_constants}}_t = (\xi_t - \xi_{\text{ss}}) - \sum_c \xi^c_t ,
\end{equation}
the time-varying intercepts (perceived means) and their propagation. The decomposition is additive
by construction. On the RE model \code{'tv'} equals \code{'iris'} (toy $10^{-13}$). The matrices with
their constants rebuild \code{res.paths} quarter by quarter (\code{D.check.recon_maxerr};
\code{x_exp1_BR} is left out in RE-model quarters, where it holds the history-stage forecasts).

\paragraph{Reading.} Slope beliefs and shares act inside every shock component, through each
quarter's $\T_t$ and $\Rs_t$, so they are not a component of their own; \code{beliefs_constants} is
the intercept channel only.

\paragraph{Output.} \code{D.values.<v>} holds one column per component (dates $\times$
components, model units) and \code{D.labels} their labels; \code{D.check} holds the additivity
error and, under \code{'tv'}, the reconstruction error \code{recon_maxerr}.
\code{decomposition.csv} has the same values in long form; \code{docs/USER_GUIDE.md}, Section 6.1,
lists its columns.

\paragraph{Restrictions.} Ramsey results are refused (\code{hetlearn:decompose:ramsey}), and so are
results of \code{history.mode = 'kalman'} (\code{hetlearn:options:unsupported}): the filtered history
is not a path of the model with the filtered shocks, because the state covariance never collapses.

\subsection{Call options}

\subsubsection{Call options of \texorpdfstring{\codem{hetlearn.decompose}}{hetlearn.decompose}}\label{opt:decompose}
\begin{optspec}
  \item[Default] \code{groups = struct()}, \code{method = 'auto'}, \code{model = []} (loaded from
        \code{res.cfg}), \code{save = true}, \code{results_dir = results/}, \code{file = ''}
        (\code{<results_dir>/<id>/decomposition.csv}), \code{quiet = false}.
  \item[Values] \code{groups}: a struct of shock lists; \code{method}: \code{'auto'},
        \code{'iris'}, \code{'tv'}; \code{model}: \code{struct('spec', spec, 'm', m, 'info', info)}.
  \item[Math] \code{'auto'} uses \code{'iris'} for mode \code{'rational'} and \code{'tv'}
        otherwise; a group's component is the sum of its members'.
  \item[Assumptions] see above.
  \item[Interactions] a shock in two groups or an unknown shock is \code{hetlearn:decompose:groups};
        an unknown option is \code{hetlearn:decompose:option}. A learning run with the RE-smoothed
        history is decomposed with the RE model over the history (it is the RE path) and with the
        learning model over the horizon.
  \item[Code] \code{decompose.m}.
  \item[Example]
\begin{lstlisting}
G = struct('demand', {{'eb', 'eg', 'eqs'}}, 'monetary', {{'em'}});
D = hetlearn.decompose('a6_inv_learn', {'pinfobs'}, [qq(2019,1), qq(2026,2)], struct('groups', G));
\end{lstlisting}
\end{optspec}

% =====================================================================
\section{Recursive forecast evaluation}\label{sec:evaluate}
% =====================================================================

\code{E = hetlearn.evaluate(name, variants, origins, opts)} (\code{evaluate.m}) runs every variant
(a scenario function name or a cfg struct) at every origin and scores the forecasts of the
observables.

\paragraph{Runs.} Each variant's cfg gets the origin, \code{forecast.horizon = opts.horizon}, zero
future shocks, no conditions, density off, \code{compare_with = {}} and
\code{history.realtime = opts.realtime}; only the expectations side of the scenario matters. Each
variant's model is loaded once. With \code{opts.realtime} the RE filter of each origin is run on data
up to that origin, once per origin and filter group, and shared by the variants of the group (the
group key holds the PLM spec and \code{model_params}; lagged variants add \code{re_info}, the
learning set and \code{re_variant}). An \code{'inversion'} variant gets one inversion to the last
origin on the full-sample filter, keeping its state at every origin, and each forecast starts from
the state at its origin (the inversion is one-sided). The parameters stay at their full-sample
values. A failing run is recorded in \code{settings.failures} and its errors are missing.

\paragraph{Errors and metrics.} For variant $v$, origin $i$, horizon $h$ and observable $k$, in report
units (scale $c_k$),
\begin{equation}
  e_{v,i,h,k} = c_k\bigl(y_{k,T_i+h} - \hat y_{v,k,T_i+h|T_i}\bigr),
\end{equation}
with the actuals from the package database (the current vintage); targets without data are skipped.
By sample (\code{full} and each subsample of target dates), variant, variable and $h$: $n$,
$\mathrm{RMSE} = \sqrt{\overline{e^2}}$, $\mathrm{ME} = \bar e$, $\mathrm{MAE} = \overline{|e|}$ over the
variant's errors; the relative RMSE $\mathrm{RMSE}_v/\mathrm{RMSE}_b$ against the benchmark $b$ on
the targets both have.

\paragraph{Diebold--Mariano test} \citep{diebold1995}. On the common targets, in origin order,
$d_i = e_{v,i}^2 - e_{b,i}^2$ (a negative mean favors the variant), $u = d - \bar d$,
$\hat\gamma_k = \frac1n\sum_{i>k}u_iu_{i-k}$, $L = \min(h-1, n-1)$, and
\begin{equation}\label{eq:dm}
  \widehat{\mathrm{LRV}} =
  \begin{cases}
    \hat\gamma_0 + 2\sum_{k=1}^{L}\hat\gamma_k & \text{\codem{'uniform'}}\\
    \hat\gamma_0 + 2\sum_{k=1}^{L}\bigl(1 - \tfrac{k}{L+1}\bigr)\hat\gamma_k & \text{\codem{'bartlett'}}
  \end{cases}
  \qquad
  \mathrm{DM} = \frac{\bar d}{\sqrt{\widehat{\mathrm{LRV}}/n}} .
\end{equation}
The uniform kernel is the rectangular window for which \citet{harvey1997} (HLN) derived their correction;
where its variance is not positive the cell falls back to Bartlett \citep{newey1987} (kernel column
\code{bartlett_fallback}, counted in \code{settings.dm_fallbacks}). With the HLN correction the
statistic is multiplied by $\sqrt{(n + 1 - 2h + h(h-1)/n)/n}$ and the two-sided p-value comes from
Student's $t$ with $n-1$ degrees of freedom, $p = I_{\nu/(\nu+\mathrm{DM}^2)}(\nu/2, 1/2)$; without it
$p = \mathrm{erfc}(|\mathrm{DM}|/\sqrt2)$. DM is NaN when $n<2$ or the variance is not positive.

\paragraph{Outputs} (\code{results/eval_<name>/}): the errors in \code{errors_long.csv}, the
tables \code{metrics.csv}, \code{relative_rmse.csv} and \code{dm.csv}, and \code{settings.mat};
\code{docs/USER_GUIDE.md}, Section 6.4, lists their columns. \code{E} also holds the forecasts and
actuals as arrays. \code{hetlearn.evaluate_tables(errors, opts)} rebuilds the three tables from a saved
\code{errors_long.csv} (or an \code{eval_<name>} folder) without rerunning any forecast, for example
with another kernel.

\subsection{Call options}

\subsubsection{Call options of \texorpdfstring{\codem{hetlearn.evaluate}}{hetlearn.evaluate}}\label{opt:evaluate}
\begin{optspec}
  \item[Default] \code{horizon = 8}, \code{benchmark = ''} (the first variant), \code{dm = true},
        \code{dm_hln = true}, \code{dm_kernel = ''} (\code{'uniform'} with HLN, else
        \code{'bartlett'}), \code{subsamples = struct()}, \code{realtime = true},
        \code{results_dir = results/}, \code{save = true}, \code{quiet = false}.
  \item[Values] \code{subsamples}: a struct of \code{[start, end]} target-date ranges;
        \code{dm_kernel}: \code{''}, \code{'uniform'}, \code{'bartlett'}; \code{origins}: a date
        vector or \code{[start, end]}.
  \item[Math] equations above.
  \item[Assumptions] real-time information for the history and the starting state (with
        \code{realtime}); full-sample parameters; actuals of the current vintage.
  \item[Interactions] all variants must use the same package (\code{hetlearn:evaluate:model}) and
        have unique ids; the benchmark must be a variant. Lagged and inversion variants get their
        own filters as described above. \code{evaluate_tables} takes the table options, with
        defaults from \code{settings.mat} next to the file (\code{docs/USER_GUIDE.md},
        Section 7).
  \item[Code] \code{evaluate.m}, \code{evaluate_tables.m}.
  \item[Example]
\begin{lstlisting}
E = hetlearn.evaluate('my_eval', {'e01_re', 'e02_learn'}, [qq(2010,1), qq(2024,2)], ...
    struct('horizon', 8, 'subsamples', struct('post2008', [qq(2008,1), qq(2019,4)])));
\end{lstlisting}
\end{optspec}

% =====================================================================
\section{Scenario files and caches}\label{sec:options}
% =====================================================================

\subsection{Scenario files}\label{sec:scenarios}

A scenario is a function that returns a partial option tree. Inheritance is a function call: a file
starts from its parent and sets what differs.
\begin{lstlisting}
function cfg = my_learn()
cfg = a1_re();                        % the RE parent: model, origin, horizon, learning set
cfg.id = 'my_learn';                  % results/my_learn/
cfg.expectations.mode = 'switching';
cfg.compare_with = {'a1_re'};
end
\end{lstlisting}
Ids are function names; \code{compare_with}, \code{'delta'} references and the runner find
scenarios by name, so names must be unique on the path. Each result stores its merged options
(\code{res.cfg}); \code{hetlearn.load_model} re-merges a saved tree, so it runs on later versions.
Saved results never update themselves: after an edit, rerun the scenario and its descendants.

\subsection{Caches}

Generated files and RE filter caches live in \code{<package>/generated/}
(Sections~\ref{sec:gencache} and~\ref{sec:refilter}). Every cache key includes what the cached object
depends on (model text, parameters, toggles, data, range), so a stale cache is never used; deleting
the folder only forces a rebuild.

% =====================================================================
\section{Outputs and the result struct}\label{sec:outputs}
% =====================================================================

\code{res = hetlearn.run(cfg)} returns the result struct and, unless \code{'save', false}, writes
\code{results/<id>/} (\code{+engine/assemble.m}, \code{+engine/save_results.m}).
\code{docs/USER_GUIDE.md}, Section 6, describes the same files for the project's scenarios.

\subsection{The result struct}

\begin{longtable}{>{\raggedright\arraybackslash}p{3.0cm}>{\raggedright\arraybackslash}p{12cm}}
\toprule
Field & Content \\
\midrule
\endhead
\code{cfg} & the merged option tree \\
\code{paths} & databank over [filter start, $T_0+H$] of every transition and measurement variable and shock: the history to the origin (RE-smoothed, backed out over the inversion, or filtered under the Kalman history), the scenario after it. \code{x_exp1_BR} holds the learners' forecasts over the record dates and NaN before \\
\code{learning} & struct by PLM target (Section~\ref{sec:learningfields}) \\
\code{events} & table: \code{date}, \code{type}, \code{variable}, \code{detail}, \code{stage} (Section~\ref{sec:events}) \\
\code{state} & the learning state at the end of the horizon: \code{phi}, \code{R}, \code{zetaRE}, \code{zetaBR}, \code{n1}, \code{fRE}, \code{fBR}, \code{fe_hist}, \code{phi_hist}, \code{gain}, \code{prev}, \code{n1m_last}, \code{n1m_keep} \\
\code{dates} & \code{origin}, \code{horizon}, \code{init_date}, \code{record} ($t_0,\dots,T_0+H$, or $T_0,\dots$ without history learning), \code{full} \\
\code{meta} & model, folder, learning variables, name lists, extended file and its MD5, collapse-test results, RE filter hash, fit error, end date and real-time flag, IRIS version, \code{history_mode}, \code{engine_step}, \code{projection_facility}, \code{history_runtime_s}, \code{history_alm_check_s}, \code{history_alm_n_check}, \code{history_alm_n_solve}, \code{forecast_n_solve}, \code{started}, \code{finished}, \code{runtime_s} \\
\code{loss} & the policy loss (Section~\ref{sec:loss}); empty without \code{spec.policy} \\
\code{history_mode} & \code{history.mode} \\
\code{inversion} & inversion diagnostics (Section~\ref{sec:inversion}) \\
\code{kalman} & Kalman history output (Section~\ref{sec:kalman}) \\
\code{density} & density output, in memory only (Section~\ref{sec:density}) \\
\bottomrule
\end{longtable}

\subsection{Learning fields}\label{sec:learningfields}

\code{res.learning.<target>} holds IRIS Series over the record dates for every PLM equation,
auxiliary ones included:
\begin{itemize}[nosep]
  \item one series per belief parameter (\code{alpha_pinf}, \code{beta_pinf},
        \code{beta_pinf__ygap_l1}, \dots): $\phi_t$, the beliefs at the end of $t$, in force in
        $t+1$ (after the facility, reverted ones included);
  \item \code{R_<param>}: the diagonal of $R_t$;
  \item \code{gain}: the gain used in $t$ (NaN where there was no update); \code{gain_R} under KOZ;
  \item \code{isAux}, \code{params} (the parameter names).
\end{itemize}
For learning variables also: \code{n1}, \code{n2} (discrete-choice shares at the end of $t$),
\code{n1_model} (the RE share in the model in $t$; NaN over the RE-smoothed history),
\code{zetaRE}, \code{zetaBR}, \code{FE_RE}, \code{FE_BR} (the variable's own squared errors;
under universal shares the group sums drive the shares) and \code{fRE}, \code{fBR} (forecasts of
$x_{t+1}$ made in $t$). Under switching, \code{n1_model}$(t)$ = \code{n1}$(t-1)$ unless the facility
kept the shares (\code{test_engine} E11 checks the timing from the paths alone).

\subsection{Share names}\label{sec:sharenames}

\begin{center}
\begin{tabular}{>{\raggedright\arraybackslash}p{5.2cm}>{\raggedright\arraybackslash}p{3.4cm}>{\raggedright\arraybackslash}p{6.2cm}}
\toprule
Name & Where & Meaning \\
\midrule
\code{shares.init_n2}, \code{shares.fixed_n2} & options & learners' share \\
\code{n1_<x>}, \code{n2_<x>} & model parameters & RE and learners' shares of $x$'s expectation in the model \\
\code{n1}, \code{n2} & \code{res.learning.<x>} & discrete-choice shares at the end of $t$ \\
\code{n1_model} & \code{res.learning.<x>} & RE share in the model in $t$ \\
\code{n2_model_<x>} & density outputs & $1-$\code{n1_model} \\
\bottomrule
\end{tabular}
\end{center}

\subsection{Files in \texorpdfstring{\codem{results/<id>/}}{results/<id>/}}

\begin{longtable}{>{\raggedright\arraybackslash}p{5.2cm}>{\raggedright\arraybackslash}p{9.8cm}}
\toprule
File & Content \\
\midrule
\endhead
\code{results.mat} & \code{res} (without \code{density}) \\
\code{observables.csv} & date and every observable, [filter start, $T_0+H$], model units \\
\code{paths.csv} & date and every transition variable over the horizon \\
\code{learning.csv} & long: date, variable (PLM target), field, value \\
\code{events.csv} & date, type, variable, detail, stage \\
\code{compare_<id>.csv} & this run minus run \code{<id>}, observables, horizon \\
\code{loss.csv} & target, expression, weight, contribution, share, discount, horizon, continuation, mode, type; a total row \\
\code{density.mat}, \code{density_quantiles.csv}, \code{density_events.csv} & Section~\ref{sec:density} \\
\code{decomposition.csv} & after \code{hetlearn.decompose} \\
\code{kalman_errors.csv} & Kalman history (a stale file is deleted otherwise) \\
\bottomrule
\end{longtable}

Dates are written as text (\code{2022Q3}); numbers with 17 significant digits. Observables and
learning series are in model units; \code{value_report} columns multiply by
\code{spec.report.scale}.

% =====================================================================
\section{Validation and the test suite}\label{sec:validation}
% =====================================================================

\subsection{Checks inside every run}

\begin{longtable}{>{\raggedright\arraybackslash}p{4.2cm}>{\raggedright\arraybackslash}p{7.4cm}>{\raggedright\arraybackslash}p{3cm}}
\toprule
Check & What it compares & Tolerance \\
\midrule
\endhead
Collapse test & extended model at $n_1=1$, zero beliefs, against the original (Section~\ref{sec:collapse}) & $10^{-10}$ \\
PLM consistency & the run's BR forecast text against the generated file & exact text \\
Anchored coefficients & state-space coefficients against unit-perturbation simulations & $10^{-9}$ relative \\
Lagged anchored model & anchored against model-consistent RE(t-1) model at $n_1 = 1$ & $10^{-10}$ \\
RE filter fit & observables rebuilt through the measurement equations & $10^{-8}$ \\
Timing check & model's \code{x_exp1_BR} against the PLM formula with $\phi_{t-1}$ & $10^{-9}$ \\
Conditions & every target met after the shock solve & $10^{-9}$ \\
Kalman & positive definite $F_t$; stationary first ALM under \code{'stochastic'} & roots below $1-10^{-10}$ \\
\bottomrule
\end{longtable}

\subsection{The test suite}

\code{tests/run_tests.m} runs the test suites of the toolbox and prints one PASS/FAIL line per
check; \code{run_tests({'test_options', 'test_rls'})} runs a subset. The suites check the equations
of this appendix against hand-coded loops, IRIS references and stored results in
\code{tests/reference/}, on the toy package \code{tests/toy_nk/} (a three-equation hybrid New
Keynesian model, \citealp[eq.~1]{ozden2025}) and on SW07 test scenarios
(\code{tests/scenarios/tsw07_*.m}).

\subsection{Known limitations}

\begin{itemize}[nosep]
  \item Shocks and conditions are unanticipated (Section~\ref{sec:timing}).
  \item Nonlinear models and log variables: Section~\ref{sec:modelreq}.
  \item Ramsey policy: Sections~\ref{opt:policy-mode} and~\ref{opt:policy-type}.
  \item \code{history.realtime} acts on the RE-smoother history only
        (Section~\ref{opt:history-realtime}).
  \item The evaluation uses the current data vintage and full-sample parameters
        (Section~\ref{sec:evaluate}).
  \item With zero future shocks the shares drift toward $1/2$ (Section~\ref{sec:dc}).
\end{itemize}

% =====================================================================
\appendix
\section{Option index}\label{app:index}
% =====================================================================

Every option of \code{hetlearn.defaults()} within the scope of this appendix, with its default and
the entry that documents it. Defaults in brackets are resolved later (\code{[]} or \code{''}).

\begin{longtable}{>{\raggedright\arraybackslash}p{5.2cm}>{\raggedright\arraybackslash}p{6.0cm}>{\raggedright\arraybackslash}p{1.6cm}>{\raggedright\arraybackslash}p{1.4cm}}
\toprule
Option & Default & Section & Page \\
\midrule
\endhead
\code{id}, \code{title}, \code{description} & \code{''} & \ref{opt:id} & \pageref{opt:id} \\
\code{model} & \code{''} & \ref{opt:model} & \pageref{opt:model} \\
\code{model_params} & \code{''} (= \code{spec.params}) & \ref{opt:model-params} & \pageref{opt:model-params} \\
\code{forecast.origin} & \code{[]} (required) & \ref{opt:forecast-origin} & \pageref{opt:forecast-origin} \\
\code{forecast.horizon} & \code{16} & \ref{opt:forecast-horizon} & \pageref{opt:forecast-horizon} \\
\code{expectations.mode} & \code{'rational'} & \ref{opt:expectations-mode} & \pageref{opt:expectations-mode} \\
\code{expectations.variables} & \code{{}} & \ref{opt:expectations-variables} & \pageref{opt:expectations-variables} \\
\code{expectations.re_variant} & \code{'anchored'} & \ref{opt:expectations-re-variant} & \pageref{opt:expectations-re-variant} \\
\code{expectations.re_info} & \code{'current'} & \ref{opt:expectations-re-info} & \pageref{opt:expectations-re-info} \\
\code{beliefs.update} & \code{'rls'} & \ref{opt:beliefs-update} & \pageref{opt:beliefs-update} \\
\code{beliefs.info} & \code{'lagged'} & \ref{opt:beliefs-info} & \pageref{opt:beliefs-info} \\
\code{beliefs.plm.default} & \code{'ar1'} & \ref{opt:beliefs-plm-default} & \pageref{opt:beliefs-plm-default} \\
\code{beliefs.plm.by_var} & \code{struct()} & \ref{opt:beliefs-plm-by-var} & \pageref{opt:beliefs-plm-by-var} \\
\code{beliefs.plm.derived} & \code{struct()} & \ref{opt:beliefs-plm-derived} & \pageref{opt:beliefs-plm-derived} \\
\code{beliefs.plm.var} & no VAR (lags 1, constant) & \ref{opt:beliefs-plm-var} & \pageref{opt:beliefs-plm-var} \\
\code{beliefs.init.method} & \code{'re_moments'} & \ref{opt:beliefs-init-method} & \pageref{opt:beliefs-init-method} \\
\code{beliefs.init.date} & \code{[]} (filter start + 4) & \ref{opt:beliefs-init-date} & \pageref{opt:beliefs-init-date} \\
\code{beliefs.init.values} & \code{struct()} & \ref{opt:beliefs-init-values} & \pageref{opt:beliefs-init-values} \\
\code{beliefs.init.training_range} & \code{[]} & \ref{opt:beliefs-init-training-range} & \pageref{opt:beliefs-init-training-range} \\
\code{gain.rule} & \code{'constant'} & \ref{opt:gain-rule} & \pageref{opt:gain-rule} \\
\code{gain.value} & \code{0.0585} & \ref{opt:gain-value} & \pageref{opt:gain-value} \\
\code{gain.by_var} & \code{struct()} & \ref{opt:gain-by-var} & \pageref{opt:gain-by-var} \\
\code{gain.koz.mu} & \code{0.001} & \ref{opt:gain-koz-mu} & \pageref{opt:gain-koz-mu} \\
\code{gain.koz.gamma0} & \code{[]} (= \code{gain.value}) & \ref{opt:gain-koz-gamma0} & \pageref{opt:gain-koz-gamma0} \\
\code{gain.koz.Rg0} & \code{0} & \ref{opt:gain-koz-rg0} & \pageref{opt:gain-koz-rg0} \\
\code{gain.koz.bounds} & \code{[0.001 0.25]} & \ref{opt:gain-koz-bounds} & \pageref{opt:gain-koz-bounds} \\
\code{gain.koz.error} & \code{'rls'} & \ref{opt:gain-koz-error} & \pageref{opt:gain-koz-error} \\
\code{gain.koz.error_scale} & \code{1} & \ref{opt:gain-koz-error-scale} & \pageref{opt:gain-koz-error-scale} \\
\code{gain.koz.structure} & \code{'common'} & \ref{opt:gain-koz-structure} & \pageref{opt:gain-koz-structure} \\
\code{gain.koz.master} & \code{''} (first learning variable) & \ref{opt:gain-koz-master} & \pageref{opt:gain-koz-master} \\
\code{shares.structure} & \code{'per_variable'} & \ref{opt:shares-structure} & \pageref{opt:shares-structure} \\
\code{shares.master} & \code{{}} & \ref{opt:shares-master} & \pageref{opt:shares-master} \\
\code{shares.init_n2} & \code{0.5} & \ref{opt:shares-init-n2} & \pageref{opt:shares-init-n2} \\
\code{shares.fixed_n2} & \code{[]} & \ref{opt:shares-fixed-n2} & \pageref{opt:shares-fixed-n2} \\
\code{switching.rule} & \code{'discrete_choice'} & \ref{opt:switching-rule} & \pageref{opt:switching-rule} \\
\code{switching.intensity} & \code{0.51} & \ref{opt:switching-intensity} & \pageref{opt:switching-intensity} \\
\code{switching.memory} & \code{0.6} & \ref{opt:switching-memory} & \pageref{opt:switching-memory} \\
\code{projection.facility} & \code{'plm_alm'} & \ref{opt:projection-facility} & \pageref{opt:projection-facility} \\
\code{history.learn} & \code{true} & \ref{opt:history-learn} & \pageref{opt:history-learn} \\
\code{history.mode} & \code{'re_smoother'} & \ref{opt:history-mode} & \pageref{opt:history-mode} \\
\code{history.re_states} & \code{'smoothed'} & \ref{opt:history-re-states} & \pageref{opt:history-re-states} \\
\code{history.realtime} & \code{false} & \ref{opt:history-realtime} & \pageref{opt:history-realtime} \\
\code{history.inversion.start} & \code{[]} (= init date) & \ref{opt:history-inversion-start} & \pageref{opt:history-inversion-start} \\
\code{history.inversion.shocks} & \code{{}} (transition shocks) & \ref{opt:history-inversion-shocks} & \pageref{opt:history-inversion-shocks} \\
\code{history.kalman.init} & \code{'stochastic'} & \ref{opt:history-kalman-init} & \pageref{opt:history-kalman-init} \\
\code{shocks.source} & \code{'none'} & \ref{opt:shocks-source} & \pageref{opt:shocks-source} \\
\code{shocks.range} & \code{[]} (horizon) & \ref{opt:shocks-range} & \pageref{opt:shocks-range} \\
\code{shocks.exclude} & \code{{}} & \ref{opt:shocks-exclude} & \pageref{opt:shocks-exclude} \\
\code{shocks.add} & \code{{}} & \ref{opt:shocks-add} & \pageref{opt:shocks-add} \\
\code{conditions} & empty struct array & \ref{opt:conditions} & \pageref{opt:conditions} \\
\code{density.enabled} & \code{false} & \ref{opt:density-enabled} & \pageref{opt:density-enabled} \\
\code{density.shocks} & \code{{}} & \ref{opt:density-shocks} & \pageref{opt:density-shocks} \\
\code{density.sigma} & \code{'calibrated'} & \ref{opt:density-sigma} & \pageref{opt:density-sigma} \\
\code{density.sigma_range} & \code{[]} (filter range) & \ref{opt:density-sigma-range} & \pageref{opt:density-sigma-range} \\
\code{density.method} & \code{'normal'} & \ref{opt:density-method} & \pageref{opt:density-method} \\
\code{density.demean} & \code{true} & \ref{opt:density-demean} & \pageref{opt:density-demean} \\
\code{density.ndraws} & \code{200} & \ref{opt:density-ndraws} & \pageref{opt:density-ndraws} \\
\code{density.seed} & \code{0} & \ref{opt:density-seed} & \pageref{opt:density-seed} \\
\code{density.quantiles} & \code{[0.05 0.16 0.5 0.84 0.95]} & \ref{opt:density-quantiles} & \pageref{opt:density-quantiles} \\
\code{density.range} & \code{[]} (horizon) & \ref{opt:density-range} & \pageref{opt:density-range} \\
\code{density.keep_draws} & \code{false} & \ref{opt:density-keep-draws} & \pageref{opt:density-keep-draws} \\
\code{density.fast_path} & \code{true} & \ref{opt:density-fast-path} & \pageref{opt:density-fast-path} \\
\code{engine.step} & \code{'fast'} & \ref{opt:engine-step} & \pageref{opt:engine-step} \\
\code{engine.conditions_weighting} & \code{'std'} & \ref{opt:engine-conditions-weighting} & \pageref{opt:engine-conditions-weighting} \\
\code{policy.mode} & \code{'rule'} & \ref{opt:policy-mode} & \pageref{opt:policy-mode} \\
\code{policy.type} & \code{'commitment'} & \ref{opt:policy-type} & \pageref{opt:policy-type} \\
\code{policy.loss} & \code{pinf} 1, \code{ygap} 0.5, \code{r} 0.1, \code{dr} 0 & \ref{opt:policy-loss} & \pageref{opt:policy-loss} \\
\code{policy.discount} & \code{[]} (= spec) & \ref{opt:policy-discount} & \pageref{opt:policy-discount} \\
\code{policy.loss_horizon} & \code{[]} (= horizon) & \ref{opt:policy-loss-horizon} & \pageref{opt:policy-loss-horizon} \\
\code{compare_with} & \code{{}} & \ref{opt:compare-with} & \pageref{opt:compare-with} \\
\code{report.variables} & \code{{}} (reserved) & \ref{opt:report-variables} & \pageref{opt:report-variables} \\
\bottomrule
\end{longtable}

Not covered here (bounded-rationality extensions, off by default): \code{expectations.types},
\code{expectations.transforms}, \code{shares.init}, \code{shares.fixed}. Call options of
\code{hetlearn.decompose} and \code{hetlearn.evaluate} are in Sections~\ref{opt:decompose}
and~\ref{opt:evaluate}; those of \code{hetlearn.extend_model} in Section~\ref{opt:extend-model}.

% =====================================================================
\bibliographystyle{plainnat}
\bibliography{references}

\end{document}
